% !TeX program = lualatex
% Tệp gốc ghép Chương 1--3. Hiện chỉ kích hoạt Chương 1.
\documentclass[12pt,a4paper,oneside]{report}
\usepackage[a4paper,left=25mm,right=22mm,top=24mm,bottom=24mm]{geometry}
\usepackage{fontspec}
\defaultfontfeatures{Ligatures=TeX}
\setmainfont{Noto Serif}
\setsansfont{Noto Sans}
\setmonofont{DejaVu Sans Mono}
\usepackage[vietnamese,english]{babel}
\babelprovide[import,main]{vietnamese}
\usepackage{microtype}
\usepackage{graphicx,float,booktabs,longtable,array,multirow,tabularx,pdfpages}
\usepackage{xcolor,amsmath,amssymb,enumitem,caption,fancyhdr,titlesec,tikz,hyperref}
\usetikzlibrary{arrows.meta,positioning,shapes,calc,fit}
\definecolor{iotblue}{HTML}{1B3A70}
\definecolor{iotcyan}{HTML}{087EAC}
\definecolor{iotorange}{HTML}{EA6B21}
\definecolor{iotlight}{HTML}{EAF2F8}
\hypersetup{colorlinks=true,linkcolor=iotblue,urlcolor=iotcyan,pdftitle={Tong hop kien thuc IoT - Chuong 1},pdfauthor={Bien soan tieng Viet tu bai giang cua TS. Nguyen Dac Cu}}
\titleformat{\chapter}[display]{\sffamily\bfseries\color{iotblue}}{\Large CHƯƠNG \thechapter}{1em}{\Huge}
\titleformat{\section}{\sffamily\Large\bfseries\color{iotblue}}{\thesection}{0.7em}{}
\titleformat{\subsection}{\sffamily\large\bfseries\color{iotcyan}}{\thesubsection}{0.65em}{}
\setlist[itemize]{topsep=3pt,itemsep=2pt,leftmargin=1.55em}
\setlist[enumerate]{topsep=3pt,itemsep=2pt,leftmargin=1.55em}
\captionsetup{font=small,labelfont=bf,justification=centering}
\setlength{\headheight}{14pt}
\pagestyle{fancy}\fancyhf{}\fancyhead[L]{\small\sffamily Cơ sở và thiết kế hệ thống IoT}\fancyhead[R]{\small\sffamily Chương \thechapter}\fancyfoot[C]{\thepage}\renewcommand{\headrulewidth}{0.4pt}
\newcommand{\chonesourcepage}[1]{\hfill{\footnotesize\color{gray}(Nguồn: slide #1, Chương 1)}}
% Hình gốc từ bài giảng: tệp JPEG được xuất bởi scripts/extract_slides.py.
% Nếu tài nguyên chưa được sinh, tài liệu vẫn biên dịch; hình tương ứng sẽ có ghi chú.
\newcommand{\choneslidefigure}[2]{\begin{figure}[H]\centering\IfFileExists{assets/chapter1/slide#1.jpg}{\includegraphics[width=0.91\linewidth]{assets/chapter1/slide#1.jpg}}{\fcolorbox{iotblue}{iotlight}{\parbox{.86\linewidth}{\centering\itshape Hình gốc ở slide #1.\\Đặt PDF bài giảng trong thư mục nguồn và chạy script trích xuất hình.}}}\caption{#2 (slide #1 của bài giảng gốc).}\end{figure}}
\newcommand{\choneemph}[1]{\textcolor{iotblue}{\textbf{#1}}}
\begin{document}
\emergencystretch=3em
\selectlanguage{vietnamese}
\renewcommand{\contentsname}{Mục lục}
\renewcommand{\figurename}{Hình}
\renewcommand{\tablename}{Bảng}
\begin{titlepage}
\centering\vspace*{22mm}
{\Large\sffamily TRƯỜNG ĐẠI HỌC PHENIKAA\\KHOA ĐIỆN -- ĐIỆN TỬ\par}
\vspace{28mm}{\Huge\sffamily\bfseries\color{iotblue}CƠ SỞ VÀ THIẾT KẾ HỆ THỐNG IoT\par}
\vspace{8mm}{\LARGE\color{iotorange}\bfseries CHƯƠNG 1 -- TỔNG QUAN VỀ IoT\par}
\vspace{5mm}{\large Bản tổng hợp và hệ thống hóa kiến thức bằng tiếng Việt\par}
\vfill
\begin{tabular}{rl}\textbf{Giảng viên tài liệu gốc:}& TS. Nguyễn Đắc Cử\\
\textbf{Mã học phần:}& EEE703006\\
\textbf{Nguồn chính:}& Chapter 1: An overview of IoT (50 slide)\\
\textbf{Trình biên dịch:}& LuaLaTeX\\
\textbf{Phạm vi:}& Chỉ Chương 1; sẵn sàng ghép Chương 2 và 3\end{tabular}
\vspace{16mm}\end{titlepage}
\pagenumbering{roman}\tableofcontents\clearpage\pagenumbering{arabic}
% chapter1.tex: chỉ chứa nội dung Chương 1; không có \documentclass/\begin{document}.
% Biên dịch bằng lualatex chapter1_main.tex hai lượt.
% Các macro có tiền tố chone tránh xung đột khi ghép với chương khác.
\providecommand{\chonesourcepage}[1]{\hfill{\footnotesize\color{gray}(Nguồn: slide Chương 1, #1)}}
\providecommand{\choneemph}[1]{\textcolor{iotblue}{\textbf{#1}}}
\providecommand{\choneslidefigure}[2]{\begin{figure}[H]\centering\IfFileExists{assets/chapter1/slide#1.jpg}{\includegraphics[width=0.91\linewidth]{assets/chapter1/slide#1.jpg}}{\parbox{.85\linewidth}{\centering Hình nguồn slide #1; chạy script extract\_chapter1\_slides.py.}}\caption{#2 (slide #1 của bài giảng gốc).}\end{figure}}
\chapter{Tổng quan về Internet vạn vật (IoT)}\label{chap:overview}
\noindent\textbf{Nguồn cơ sở:} 50 slide ``Chapter 1: An overview of IoT'', học phần \emph{Cơ sở và thiết kế hệ thống IoT} của TS. Nguyễn Đắc Cử, Khoa Điện -- Điện tử, Trường Đại học Phenikaa. Nội dung dưới đây được diễn giải lại bằng tiếng Việt và đối chiếu với từng slide. Phần hình từ slide được đặt cạnh nội dung tương ứng; các biểu đồ dự báo mang tính lịch sử được ghi đúng bối cảnh nguồn.\par
\medskip
\noindent\textbf{Sáu nhóm kiến thức của chương} (slide 2): Giới thiệu học phần; lịch sử Internet và IoT; các khái niệm IoT; kiến trúc IoT; ứng dụng IoT; thách thức và cơ hội của IoT.

\section{Giới thiệu học phần}\label{sec:course}
\subsection{Thông tin học phần và phạm vi nghiên cứu}
Tên tiếng Việt là \emph{Cơ sở và thiết kế hệ thống IoT}; tên tiếng Anh \emph{IoT System Design}; mã học phần \textbf{EEE703006}. Trọng tâm của môn học là các khái niệm IoT, nền tảng phần cứng, phần mềm ứng dụng, các cơ chế xử lý dữ liệu và thông tin, cũng như các công nghệ/giao thức truyền thông được nêu trong slide: \textbf{Zigbee, Bluetooth, IEEE 802.15.4, IEEE 802.15.6 và IEEE 802.15.11}. Dòng ``Syllabus: Click here'' trong slide 3 là liên kết giáo trình/đề cương trên bản trình chiếu; bản trích xuất PDF không cho biết nội dung đích nên không suy diễn. \chonesourcepage{3}

\subsection{Quy định học tập và thời khóa biểu được ghi trong slide}
Các yêu cầu học tập: tham dự tối thiểu \textbf{80\%} số buổi học trên lớp; tham dự \textbf{100\%} buổi thực hành; hoàn thành các nhiệm vụ của học phần. Thời khóa biểu dưới đây được giữ \emph{nguyên thông tin lịch sử của lớp năm 2022}, không phải lịch hiện hành.\chonesourcepage{4}
\begin{table}[H]\centering\small
\caption{Thời khóa biểu trong tài liệu gốc (slide 4).}
\begin{tabularx}{\linewidth}{@{}l l l l l X@{}}\toprule
Hình thức & Mã lớp & Thời gian & Thứ & Tiết & Phòng\\\midrule
Lý thuyết & IoT-1-1-22(N01) & 15/08--27/11/2022 & Hai & 10, 11 & A3-302\\
Thực hành & IoT-1-1-22(N01.TH1) & 19/09--27/11/2022 & Năm & 4, 5, 6 & A4-503\\
Thực hành & IoT-1-1-22(N01.TH2) & 19/09--27/11/2022 & Bảy & 1, 2, 3 & A4-503\\\bottomrule
\end{tabularx}\end{table}

\section{Lịch sử Internet và sự hình thành IoT}\label{sec:history}
\subsection{Khái niệm Internet vạn vật và bước chuyển của máy tính}
\emph{Internet of Things (IoT)} -- Internet vạn vật hay mạng lưới vạn vật kết nối Internet -- là hệ thống các thiết bị tính toán có liên hệ với nhau, trong đó các vật thể, động vật hay con người có thể được gắn mã nhận dạng riêng, có khả năng thu nhận và trao đổi thông tin qua mạng mà không cần con người trực tiếp tham gia vào mọi lần tương tác giữa người--người hoặc người--máy. Định nghĩa này xuất hiện hai lần trong slide 5 và 12, nhằm nhấn mạnh đặc tính \emph{định danh, kết nối, trao đổi dữ liệu và tự động hóa giao tiếp}. \chonesourcepage{5, 12}
\choneslidefigure{05}{Hình cây kết nối các biểu tượng thiết bị và dịch vụ, mô tả ý tưởng Internet vạn vật}

\subsection{Bốn giai đoạn phát triển Internet}
Các slide 6--7 mô tả sự phát triển của Internet theo bốn giai đoạn với tác động xã hội, kinh doanh và mức độ liên kết thông minh tăng dần:
\begin{enumerate}
\item \choneemph{Connectivity -- Digitize Access (Kết nối, số hóa truy cập):} kết nối con người tới thư điện tử (email), trình duyệt/dịch vụ Web và công cụ tìm kiếm để tiếp cận thông tin.
\item \choneemph{Networked Economy -- Digitize Business (Kinh tế nối mạng, số hóa kinh doanh):} xuất hiện thương mại điện tử, chuỗi cung ứng số và cộng tác qua mạng nhằm cải thiện hiệu quả quy trình kinh doanh.
\item \choneemph{Immersive Experiences -- Digitize Interactions (Trải nghiệm phong phú, số hóa tương tác):} video, mạng xã hội, tính di động và kết nối thường xuyên trở nên phổ biến; ngày càng nhiều ứng dụng chuyển lên đám mây.
\item \choneemph{Internet of Things -- Digitize the World (Internet vạn vật, số hóa thế giới thực):} các đối tượng, thiết bị và máy móc trước đây chưa kết nối được gắn thêm khả năng liên lạc, mở ra dịch vụ và trải nghiệm mới; liên kết con người, quy trình, dữ liệu và vạn vật.
\end{enumerate}
Sơ đồ tại slide 6 thể hiện trục hoành là mức độ \emph{kết nối thông minh (intelligent connections)}, trục tung là \emph{tác động tới kinh doanh và xã hội}, được biểu diễn bằng bốn cột tăng dần. Slide 7 có bảng định nghĩa từng giai đoạn.\chonesourcepage{6--7}
\choneslidefigure{06}{Bốn giai đoạn phát triển của Internet và tác động kinh tế--xã hội}

\subsection{Từ tiền Internet đến Internet of Things}
Sơ đồ mở rộng tại slide 8 trình bày thêm tiến trình \emph{Pre-Internet} (người tới người), \emph{Internet of Content / WWW} (nội dung), \emph{Internet of Services / Web 2.0} (dịch vụ), \emph{Internet of People / Social Media} (con người) và \emph{Internet of Things / Machine-to-Machine} (máy tới máy). Các ví dụ tiêu biểu lần lượt là:
\begin{itemize}
\item Trước Internet: điện thoại cố định/di động, SMS; phát triển mạng thông minh.
\item Internet nội dung: email, thông tin, giải trí; nền tảng và dịch vụ CNTT thông minh.
\item Internet dịch vụ: năng suất điện tử và thương mại điện tử; điện thoại và ứng dụng thông minh.
\item Internet con người: Skype, Facebook, YouTube, Twitter; thiết bị thông minh, đối tượng và thẻ định danh.
\item Internet vạn vật: định danh, theo dõi, giám sát, đo lường; dữ liệu chia sẻ được tổ chức theo ngữ nghĩa, kết hợp dữ liệu và ngữ cảnh môi trường.
\end{itemize}
\choneslidefigure{08}{Sơ đồ các giai đoạn từ Pre-Internet tới Internet of Things}

\subsection{Các cột mốc lịch sử và đóng góp của các nhà nghiên cứu}
Theo slide 9, thời kỳ IoT được nhiều tài liệu gọi là giai đoạn bắt đầu khoảng \textbf{2008--2009}; song thuật ngữ \emph{Internet of Things} được ghi nhận từ \textbf{1999}, do \textbf{Kevin Ashton} sử dụng khi làm việc tại \textbf{Procter \& Gamble (P\&G)}. Ý tưởng ban đầu gắn kết quản lý chuỗi cung ứng với Internet: dùng \emph{RFID -- Radio Frequency Identification} (nhận dạng bằng sóng vô tuyến), nhãn và mã vạch để theo dõi vật phẩm, qua đó cải thiện hiệu quả quản lý. Hai mốc này không đồng nghĩa: một mốc nói về tên gọi, mốc còn lại nói về thời kỳ phổ biến của IoT.\chonesourcepage{9}

Tại slide 10, Kevin Ashton nhấn mạnh rằng trong thế kỷ XX máy tính giống như ``bộ não không có giác quan'' và phụ thuộc dữ liệu do con người nhập (gõ bàn phím, quét mã vạch, v.v.). IoT chuyển trọng tâm sang \emph{thiết bị có thể tự cảm nhận thế giới thực}. Ông lập luận trong bài viết đăng trên \emph{RFID Journal} ngày \textbf{22/06/2009} rằng nếu máy tính biết về đồ vật nhờ dữ liệu thu thập tự động, chúng ta có thể theo dõi và kiểm đếm tốt hơn, giảm lãng phí, thất thoát và chi phí, nhận biết lúc cần thay, sửa, thu hồi hoặc khi vật phẩm mất độ tươi mới. Đây là nội dung diễn giải ý chính của trích dẫn trong slide, không phải nguyên văn tiếng Anh.\chonesourcepage{10}

Slide 11 bổ sung: năm \textbf{1999}, \textbf{Neil Gershenfeld} xuất bản cuốn \emph{When Things Start to Think} (slide ghi ``When Everything Started Thinking''); bối cảnh là viễn cảnh đồ vật cùng sử dụng Internet. Slide xem \textbf{máy ATM năm 1974} là một trong những đối tượng thông minh kết nối mạng đầu tiên, đồng thời nhắc máy bán nước ngọt có kết nối Internet tại Khoa Khoa học Máy tính, \textbf{Carnegie Mellon University} vào \textbf{1980}. Khoảng \textbf{2010}, thuật ngữ \emph{Industrial Internet of Things -- IIoT} được nhấn mạnh trong công nghiệp. Slide còn đề xuất rằng ``Network/Net'' có thể mô tả kết nối rộng hơn ``Internet'', liên hệ tới tên gọi \emph{Internet of Everything} hoặc \emph{Net of Everything}. \emph{Lưu ý biên soạn:} tên sách được ghi theo slide và đồng thời nêu tên xuất bản phổ biến để người học có thể đối chiếu.\chonesourcepage{11}

\section{Khái niệm, đối tượng và môi trường thông minh}\label{sec:concepts}
\subsection{IoT là mạng lưới các ``things'' kết nối}
Slide 13 diễn giải IoT như việc kết nối thiết bị (ví dụ thiết bị có nút bật/tắt) với Internet và với các thiết bị khác. Các thực thể trong mạng -- cả đồ vật và con người -- có thể thu thập/chia sẻ dữ liệu về cách chúng được sử dụng và về môi trường xung quanh. Cách diễn đạt ``thiết bị có nút bật/tắt'' là ví dụ khái quát từ nguồn, không phải điều kiện kỹ thuật bắt buộc đối với mọi nút IoT. Hình minh họa tại đây liệt kê những lĩnh vực như phương tiện, năng lượng, giao thông, quản lý môi trường, sinh hoạt và ứng dụng dịch vụ.\chonesourcepage{13}
\choneslidefigure{13}{Đồ vật, con người và các lĩnh vực được kết nối thành Internet vạn vật}

\subsection{Phạm vi của một IoT Thing}
Ví dụ slide 14 gồm: người mang thiết bị cấy ghép theo dõi tim mạch, động vật nông nghiệp được gắn bộ phát đáp/chip sinh học, ô tô có cảm biến cảnh báo áp suất lốp thấp. Nhiều đối tượng tự nhiên hoặc nhân tạo có thể được định danh trên mạng (slide lấy ví dụ địa chỉ IP) và gửi dữ liệu qua Internet. Các đồ vật IoT vì vậy đa dạng về \emph{kích thước, cấu trúc, mục đích, cơ chế kết nối và đối tượng được theo dõi}. \chonesourcepage{14}
\choneslidefigure{14}{Ví dụ thiết bị theo dõi tim, ô tô gắn cảm biến và các đối tượng IoT}

\subsection{Smart Object -- Đối tượng thông minh}
\textbf{Smart Object (SO)} là đối tượng vật lý được tích hợp các thành phần xử lý, lưu trữ, cảm biến và kết nối mạng. Một số SO còn có thể \emph{tác động ngược lại} môi trường qua \textbf{bộ chấp hành (actuator)}. Cách biểu diễn phần cứng tại slide 16 chỉ ra các khối: cảm biến (nhận tín hiệu như độ sáng, nhiệt độ, áp suất), actuator, bộ xử lý công suất thấp, khối truyền thông, nguồn điện; trong đó dữ liệu cảm biến được bộ xử lý thu thập, xử lý và truyền tới thiết bị khác, còn actuator thực thi lệnh để thay đổi trạng thái vật lý.\chonesourcepage{15--16}
\choneslidefigure{16}{Cấu trúc đối tượng thông minh với cảm biến, cơ cấu chấp hành, xử lý, liên lạc và nguồn}

Từ hình tròn tại slide 15 có thể đọc các ý: \emph{data processing} (xử lý dữ liệu), \emph{communication} (truyền thông), \emph{power and battery} (nguồn và pin), \emph{environmental sensing} (cảm nhận môi trường), \emph{actuation} (tác động môi trường). Đây là các năng lực liên kết quanh SO.\chonesourcepage{15}

\subsection{Xu hướng phát triển của Smart Object}
Theo slide 17, năm xu thế chính là: (1) kích thước thiết bị nhỏ dần; (2) điện năng tiêu thụ giảm; (3) năng lực xử lý tăng; (4) chất lượng, khả năng truyền thông được cải thiện; và (5) giao tiếp ngày càng được chuẩn hóa. Các xu thế phối hợp khiến đối tượng thông minh được sử dụng linh hoạt hơn trong thiết bị thông thường.\chonesourcepage{17}

\subsection{Smart Environment -- Môi trường thông minh}
\textbf{Smart Environment (SE)} là phạm vi không gian vật lý nơi các SO được triển khai, kết nối và tương tác. Nhiều SO trong một không gian như xe, nhà hoặc đô thị kết hợp thành \emph{smart car, smart home, smart city}. Slide 18 trích dẫn ý tưởng của \textbf{Mark Weiser (1999)}: cảm biến, cơ cấu chấp hành và thiết bị tính toán đan xen, gần như vô hình trong đồ vật hằng ngày, kết nối liên tục; cách nói này nhấn mạnh công nghệ được tích hợp vào chính môi trường sống thay vì một máy tính riêng lẻ. Hình minh họa đặt \emph{Smart Space} ở giữa, kết nối SO và con người.\chonesourcepage{18}
\choneslidefigure{18}{Môi trường thông minh gồm các đối tượng SO và con người trong cùng không gian}

\subsection{Ubiquitous Computing -- Điện toán khắp nơi}
\textbf{Ubiquitous Computing (Ubicomp)} là mô hình điện toán hiện diện rộng khắp: một người có thể đồng thời tương tác với nhiều máy tính nhỏ, kết nối lẫn nhau, nằm trong các vật dụng hoặc môi trường hằng ngày và có thể ít được nhận thấy như một máy tính độc lập. Slide 19 coi \textbf{Mark Weiser} là người đặt nền tảng khái niệm này; hình liên hệ điện toán khắp nơi với \emph{computer}, \emph{pervasive computing}, \emph{wearables} và \emph{smart devices}.\chonesourcepage{19}
\choneslidefigure{19}{Quan hệ giữa điện toán khắp nơi, thiết bị đeo và thiết bị thông minh}

\subsection{Các khái niệm liên quan: M2M, IIoT, Industry 4.0}
\textbf{M2M -- Machine-to-Machine Communication}: các máy hoặc thiết bị trao đổi thông tin với nhau mà không cần người trực tiếp điều khiển mỗi lần truyền. Hình ở slide 20 đối chiếu M2M kết nối trực tiếp theo cặp với một hệ IoT dùng đám mây/nền tảng để phối hợp nhiều thiết bị. M2M là khái niệm liên quan mật thiết nhưng không đồng nhất với toàn bộ hệ sinh thái IoT.\chonesourcepage{20}

\textbf{IIoT -- Industrial Internet of Things}: IoT được triển khai trong hệ thống công nghiệp, với thiết bị, công nhân, cảm biến, máy móc, hệ thống điều khiển, nhà máy và nền tảng dữ liệu được kết nối nhằm giám sát, phối hợp và tối ưu vận hành. Hình slide 21 biểu diễn hệ thống công nghiệp có các thiết bị, dây chuyền, đám mây và dữ liệu kết nối hai chiều.\chonesourcepage{21}

\textbf{Industry 4.0 -- Công nghiệp 4.0}: hình slide 22 liên hệ \emph{Industrial IoT, Network Centric, Industrial Apps, Quality, Safety, Automation, Monitoring Tools, Sensors}. Như vậy IIoT là một thành tố của bức tranh công nghiệp số rộng hơn, không phải tên khác hoàn toàn tương đương với Industry 4.0.\chonesourcepage{22}
\choneslidefigure{20}{Hình so sánh mạng truyền thông M2M và mô hình IoT}
\choneslidefigure{22}{Các yếu tố trong hệ sinh thái Industry 4.0}

\section{Kiến trúc hệ thống IoT}\label{sec:architecture}
\subsection{Bốn khối chức năng cơ bản}
Theo slide 23, kiến trúc tổng quát có bốn khối \textbf{Things -- Gateway/Network -- Cloud -- Services}; hình còn thể hiện \emph{User Interface} (giao diện người dùng) liên kết ở phía ứng dụng. Luồng cơ bản: thiết bị đo đạc và điều khiển $\leftrightarrow$ cổng mạng/gateway $\leftrightarrow$ đám mây $\leftrightarrow$ dịch vụ và giao diện. Dữ liệu, lệnh điều khiển có thể đi hai chiều.\chonesourcepage{23}
\choneslidefigure{23}{Quan hệ giữa thiết bị, gateway, cloud, dịch vụ và giao diện người dùng}

\subsection{Things -- Thiết bị/nút IoT}
Thiết bị gồm \textbf{cảm biến}, \textbf{bộ chấp hành} và \textbf{bộ điều khiển}: bộ vi xử lý hoặc vi điều khiển, \emph{PLC (Programmable Logic Controller)}, \emph{FPGA (Field-Programmable Gate Array)} cho tới máy tính nhúng. Nhiệm vụ là đo, thu thập thông tin thế giới thực, điều khiển cơ cấu chấp hành và truyền/nhận dữ liệu với các thiết bị khác. Dữ liệu tầng thiết bị là đầu vào nền tảng của toàn hệ thống. Công nghệ bán dẫn mở ra cảm biến thông minh rất nhỏ; ví dụ cảm biến \emph{nhiệt độ, áp suất, độ ẩm, chuyển động, cường độ ánh sáng}.\chonesourcepage{24}

\subsection{Gateway/Network -- Cổng mạng và tầng kết nối}
Gateway quản lý \textbf{luồng dữ liệu hai chiều} giữa các mạng có chuẩn truyền thông khác nhau, thực hiện \textbf{chuyển đổi giao thức} và hỗ trợ khả năng \textbf{tương tác liên thông (interoperability)}. Gateway có thể tổng hợp/lọc/tiền xử lý dữ liệu từ hàng nghìn cảm biến cục bộ rồi mới đưa sang giai đoạn kế tiếp; trong một số triển khai nó phải tương thích với \textbf{TCP/IP}. Cổng còn là tuyến phòng vệ giữa thiết bị và đám mây nhờ cơ chế bảo mật và mã hóa đường truyền, giảm rủi ro truy cập trái phép hoặc tấn công vào mạng thiết bị.\chonesourcepage{25}

\subsection{Cloud -- Nền tảng điện toán đám mây}
Cloud cung cấp hạ tầng, phần mềm để \textbf{thu thập, xử lý, quản lý và lưu trữ dữ liệu quy mô lớn theo thời gian thực}. Các máy chủ có hiệu năng cao, hệ thống phân tán, cơ sở dữ liệu và phân tích hỗ trợ dữ liệu từ lượng rất lớn thiết bị, sensor, gateway; đồng thời hỗ trợ quản lý lưu lượng, phân tích chính xác và dự báo. Trong mô tả tại slide, \textbf{hệ quản trị cơ sở dữ liệu phân tán} là một thành phần then chốt.\chonesourcepage{26}

\subsection{Services -- Dịch vụ và phân tích thông minh}
Tầng dịch vụ biến dữ liệu thô thu được từ thiết bị/cảm biến thành \textbf{thông tin, nhận thức và chỉ dẫn có thể hành động}. Phân tích dữ liệu được dùng để hiểu và quản lý hoạt động, chuẩn bị hành động đúng thời điểm, cải thiện hệ thống và khai thác cơ hội kinh doanh về sau. Slide sử dụng cách diễn đạt ``transform analog data''; về mặt quy trình cần hiểu là chuyển dữ liệu đầu vào của thế giới thực sang biểu diễn có thể phân tích, chứ không giới hạn mọi dữ liệu đầu vào ở dạng analog.\chonesourcepage{27}

\subsection{User Interface -- Giao diện người dùng}
Là phần người sử dụng trực tiếp quan sát và thao tác với hệ IoT, ví dụ màn hình trên điện thoại hoặc bảng điều khiển cảm ứng. Một giao diện tốt biến quy trình phức tạp thành thao tác dễ hiểu, hỗ trợ hiển thị trạng thái, cảnh báo và ra lệnh điều khiển. Slide 28 minh họa giao diện \emph{My Home} trên điện thoại.\chonesourcepage{28}

\subsection{Kiến trúc phân tầng: tổng quan 3, 4 và 5 tầng}
Sơ đồ slide 29 đối chiếu ba mô hình kiến trúc IoT; tên các tầng từ dưới lên được tổng hợp như sau:
\begin{table}[H]\centering\small\caption{Đối chiếu kiến trúc 3/4/5 tầng theo slide 29 và 34--39.}
\begin{tabularx}{\linewidth}{@{}l X@{}}\toprule
Mô hình & Các tầng theo thứ tự từ dưới lên\\\midrule
3 tầng & Perception (cảm nhận) $\to$ Network (mạng) $\to$ Application (ứng dụng)\\
4 tầng (SoA) & Perception/Sensors \& Actuators $\to$ Network $\to$ Data Processing $\to$ Application\\
5 tầng & Perception $\to$ Transport/Network $\to$ Processing/Middleware $\to$ Application $\to$ Business\\\bottomrule
\end{tabularx}\end{table}
\choneslidefigure{29}{So sánh cấu trúc kiến trúc IoT ba, bốn và năm tầng}

\subsection{Chi tiết mô hình 3 tầng}
\textbf{Tầng cảm nhận (Perception layer)} bao gồm ba nhóm thể hiện ở slide 31: \emph{thing/device hardware} -- phần cứng cốt lõi; \emph{IoT components} -- sensor, actuator, processor, cổng giao tiếp và anten; \emph{thing/device software} -- phần mềm nhúng chạy trên vật thể để quản lý và vận hành chức năng. Chức năng chính là nhận biết nút IoT và thu nhận dữ liệu từ môi trường.\chonesourcepage{30--31}

\textbf{Tầng mạng (Network layer)} trong slide 32 gồm \emph{network communication}: các giao thức cho phép thiết bị/vật thể giao tiếp với hạ tầng mạng lõi và/hoặc đám mây. Chức năng là truyền dữ liệu giữa đầu cuối và dịch vụ, đồng thời liên kết các mạng và thiết bị.\chonesourcepage{32}

\textbf{Tầng ứng dụng (Application layer)} trên slide 33 được minh họa thành \emph{năm nhóm phần mềm chức năng}, không phải thêm năm tầng độc lập:
\begin{enumerate}
\item \textbf{Thing/device communications and management}: cung cấp, giao tiếp, quản lý thiết bị.
\item \textbf{Application platform}: môi trường xây dựng và thực thi ứng dụng IoT.
\item \textbf{Process management}: định nghĩa, thực thi, giám sát quy trình liên quan người, hệ thống và thiết bị.
\item \textbf{Analytics and data management}: lưu trữ, xử lý và phân tích lượng lớn dữ liệu máy móc dạng chuỗi thời gian.
\item \textbf{IoT application}: điều phối sự tương tác giữa con người, hệ thống và đồ vật nhằm phục vụ một mục tiêu cụ thể.
\end{enumerate}
\choneslidefigure{31}{Ba bộ phận cấu thành lớp cảm nhận trong kiến trúc 3 tầng}
\choneslidefigure{33}{Năm nhóm chức năng trong hình minh họa của lớp ứng dụng}

\subsection{Kiến trúc hướng dịch vụ SoA -- Bốn tầng}
\textbf{Service-Oriented Architecture (SoA)} là kiến trúc dựa trên mô hình các thành phần dịch vụ: các khối chức năng được kết nối thông qua giao diện và giao thức đã xác định. Trong cách trình bày của slide 34--36, kiến trúc bốn tầng gồm:
\begin{enumerate}
\item \textbf{Perception / Sensors and Actuators}: nhận diện thiết bị, lấy dữ liệu từ cảm biến nhiều loại; tầng phải thích nghi với đặc tính hoạt động của từng cảm biến.
\item \textbf{Networking (Transmission)}: cầu nối truyền/nhận dữ liệu từ đối tượng vật lý, kết nối thiết bị và các mạng khác nhau.
\item \textbf{Data Processing}: tiếp nhận dữ liệu của tầng cảm nhận và kiểm tra dữ liệu được chuyển từ người dùng/nguồn hợp lệ, đồng thời hỗ trợ chống các mối đe dọa.
\item \textbf{Application}: xác định ứng dụng IoT, cung cấp giao diện giữa thiết bị cuối và mạng, cho phép các dịch vụ dựa trên dữ liệu cảm biến thu thập.
\end{enumerate}
\choneslidefigure{34}{Sơ đồ bốn tầng của kiến trúc hướng dịch vụ SoA}

\subsection{Kiến trúc năm tầng}
Các tầng từ dưới lên và nội dung chi tiết của hình slide 37:
\begin{description}[leftmargin=0pt,labelindent=0pt,style=nextline]
\item[1. Perception -- Cảm nhận:] đối tượng vật lý, \emph{RFID}, mã vạch, cảm biến hồng ngoại và các nguồn dữ liệu tương tự.
\item[2. Network/Transport -- Mạng/truyền tải:] truyền an toàn; \emph{3G, UMTS, Wi-Fi, Bluetooth, hồng ngoại, ZigBee}, v.v. Slide 38 nêu thêm \emph{LAN, RFID, NFC} và liên lạc không dây làm phương tiện chuyển dữ liệu từ tầng cảm nhận tới tầng xử lý và ngược lại.
\item[3. Middleware/Processing -- Trung gian/xử lý:] \emph{information processing}, \emph{ubiquitous computing}, \emph{service management}, \emph{database}, \emph{decision unit}; lưu trữ, phân tích dữ liệu lớn, quản lý dịch vụ bằng cơ sở dữ liệu, điện toán đám mây và xử lý Big Data.
\item[4. Application -- Ứng dụng:] các ứng dụng thông minh và quản lý ứng dụng, phục vụ mục đích sử dụng cụ thể.
\item[5. Business -- Kinh doanh:] \emph{system management}, \emph{business models}, \emph{flowchart}, \emph{graphs}; quản lý toàn hệ IoT, mô hình lợi nhuận, quản trị ứng dụng và quyền riêng tư người dùng.
\end{description}
Ở hình này tầng \emph{Network} được dùng tương ứng với chức năng \emph{Transport} trong đoạn văn của slide 38; vì thế hai tên được giải thích cùng nhau thay vì coi là sáu tầng.\chonesourcepage{37--39}
\choneslidefigure{37}{Kiến trúc IoT năm tầng với các thành phần cụ thể của mỗi tầng}

\subsection{Kiến trúc sáu tầng}
Slide 40 nêu sơ đồ với năm dải chức năng nhìn thấy: \textbf{Physical/device}, \textbf{Network}, \textbf{Data/database}, \textbf{Application/integration}, \textbf{Analytics/visualization}, đồng thời có dải \textbf{Security and management} chạy xuyên suốt. Tài liệu gốc gọi là ``six-layer architecture''; cách đọc nhất quán với hình là tính \emph{Security and Management} là lớp/chức năng thứ sáu liên tầng. Lớp bảo mật và quản trị áp dụng lên toàn bộ kiến trúc, không nhất thiết là một trạm xử lý dữ liệu độc lập trên đường đi.\chonesourcepage{40}
\choneslidefigure{40}{Mô hình sáu tầng theo slide gốc, bao gồm bảo mật và quản trị liên tầng}

\subsection{Kiến trúc bảy tầng}
Hình tại slide 41 mô tả bảy tầng theo thứ tự từ thấp lên cao:
\begin{enumerate}
\item \emph{Things}: thiết bị, cảm biến, bộ điều khiển.
\item \emph{Connectivity / Edge Computing}: kết nối, giao thức, mạng, M2M, Wi-Fi, viễn thông, phần cứng trung gian; được chú thích thuộc khu vực Fog.
\item \emph{Global Infrastructure}: hạ tầng Cloud công cộng, riêng, lai hoặc được quản lý.
\item \emph{Data Ingestion}: thu nhận, tập hợp và lưu trữ dữ liệu lớn từ đồ vật.
\item \emph{Data Analysis}: báo cáo, khai phá dữ liệu và học máy.
\item \emph{Applications}: ứng dụng tùy biến từ dữ liệu của đồ vật.
\item \emph{People and Process}: quyết định chuyển đổi dựa trên dữ liệu và ứng dụng IoT.
\end{enumerate}
Sơ đồ còn nhóm tầng 1--2 là \emph{Fog}, tầng 3 là \emph{Cloud}, tầng 4--5 là \emph{Big Data}, tầng 6--7 là \emph{Business Value}. Đây là một cách nhìn khác với kiến trúc ba tầng, cho thấy đường chuyển từ \emph{thiết bị vật lý} tới \emph{giá trị kinh doanh}.\chonesourcepage{41}
\choneslidefigure{41}{Kiến trúc bảy tầng, từ thiết bị đến giá trị kinh doanh}

\section{Các ứng dụng của Internet vạn vật}\label{sec:applications}
\subsection{Mười nhóm ứng dụng minh họa}
Hình ``IoT Applications and Usecases'' tại slide 42 liệt kê \textbf{10 nhóm}: \emph{Smart home} (nhà thông minh), \emph{Wearables} (thiết bị đeo), \emph{Smart City} (đô thị thông minh), \emph{Smart grids} (lưới điện thông minh), \emph{Industrial Internet} (Internet công nghiệp), \emph{Connected car} (xe kết nối), \emph{Connected Health} (y tế kết nối), \emph{Smart Retail} (bán lẻ thông minh), \emph{Smart Supply Chain} (chuỗi cung ứng thông minh) và \emph{Smart Farming} (nông nghiệp thông minh).\chonesourcepage{42}
\begin{longtable}{@{}p{.26\linewidth}p{.68\linewidth}@{}}\caption{Ứng dụng IoT xuất hiện trong hai hình slide 42--43.}\\\toprule
Lĩnh vực & Ý nghĩa/đối tượng điển hình trong hình\\\midrule\endfirsthead\toprule Lĩnh vực & Ý nghĩa/đối tượng điển hình\\\midrule\endhead
Nhà thông minh & Điều khiển đồ gia dụng, an ninh, tiện nghi và tự động hóa trong nhà.\\
Thiết bị đeo & Theo dõi trạng thái cơ thể, hỗ trợ cá nhân trên thiết bị đeo.\\
Đô thị thông minh & Giám sát và quản trị hạ tầng, chiếu sáng, giao thông và dịch vụ đô thị.\\
Lưới điện thông minh & Thiết bị đo thông minh, giám sát và điều phối lưới điện.\\
Internet công nghiệp & Máy móc, nhà xưởng, cảm biến và vận hành sản xuất.\\
Xe kết nối & Cảm biến trên xe, chẩn đoán, chia sẻ dữ liệu xe--dịch vụ.\\
Y tế kết nối & Thiết bị theo dõi và hỗ trợ chăm sóc sức khỏe.\\
Bán lẻ thông minh & Giám sát hàng hóa, tương tác mua sắm, quản lý bán lẻ.\\
Chuỗi cung ứng thông minh & Truy vết hàng hóa, vận tải, kho và logistics.\\
Nông nghiệp thông minh & Giám sát môi trường canh tác, cây trồng và vận hành nông trại.\\\bottomrule
\end{longtable}
Slide 43 bổ sung tranh ghép minh họa ngành bán lẻ, hệ thống điện, xe thông minh, các cảm biến và bối cảnh nông nghiệp; đây là hình minh họa khả năng ứng dụng, không phải bảng phân loại chính thức có ranh giới cứng.\chonesourcepage{43}
\choneslidefigure{42}{Mười nhóm ứng dụng và tình huống sử dụng IoT}
\choneslidefigure{43}{Ảnh ghép các ngữ cảnh triển khai IoT trong đời sống và sản xuất}

\subsection{Tình huống 1 -- Đồng hồ báo thức thông minh tại nhà}
Ở slide 44, người học thường phải dậy lúc \textbf{7:00} để đi học. Một hôm tuyến xe buýt bị hủy; do trời mưa nên không đặt được xe Grab và phải chuyển sang đi xe đạp, đòi hỏi dậy lúc \textbf{6:30} để không đi muộn. Đồng hồ báo thức có kết nối IoT có thể kết hợp \emph{trạng thái giao thông/vận tải}, \emph{dự báo thời tiết}, \emph{thời gian di chuyển thay thế} và \emph{giờ đến lớp} rồi tự cập nhật giờ thức dậy. Nếu được tích hợp sâu hơn, thiết bị có thể đồng bộ máy pha cà phê kết nối mạng để chuẩn bị đồ uống đúng giờ. Ví dụ cho thấy IoT không chỉ truyền số liệu, mà còn kết hợp \emph{ngữ cảnh, quyết định và điều khiển nhiều thiết bị}.\chonesourcepage{44}

\subsection{Tình huống 2 -- Chẩn đoán xe và dịch vụ bảo dưỡng}
Ở slide 45, đèn cảnh báo động cơ sáng khi người sử dụng đang lái xe. Các cảm biến liên quan truyền dữ liệu qua \emph{chuỗi chẩn đoán (diagnostic chain)} đến \emph{gateway trên xe}; gateway gửi thông tin quan trọng tới cơ sở của nhà sản xuất. Từ đó nhà sản xuất hoặc hệ dịch vụ có thể xác định mức độ khẩn cấp, đặt lịch sửa chữa, chuyển địa chỉ đại lý gần nhất và chuẩn bị sẵn phụ tùng khi xe đến. Quy trình minh họa liên hệ giữa \emph{thu nhận dữ liệu} $\to$ \emph{truyền dữ liệu} $\to$ \emph{phân tích} $\to$ \emph{dịch vụ và hành động}. Slide kết thúc bằng gợi ý người học tự liệt kê các đồ vật có thể áp dụng IoT thực tế.\chonesourcepage{45}

\subsection{Năm dự báo IoT cho năm 2025 -- dữ liệu theo hình nguồn}
Slide 46 chứa đồ họa có tiêu đề ``5 predictions for IoT in 2025'' (ghi nguồn trên ảnh: \emph{Machina Research}). Các số liệu dưới đây là \textbf{dự báo được trích từ slide, không phải số liệu thống kê kiểm chứng năm 2025/2026}:
\begin{enumerate}
\item \textbf{27 tỷ} kết nối IoT toàn cầu (chú thích ảnh: 6 tỷ năm 2015).
\item \textbf{3 nghìn tỷ USD} dự báo doanh thu (revenue forecast).
\item \textbf{2,2 tỷ} kết nối IoT qua mạng di động (chú thích ảnh: 334 triệu năm 2015).
\item \textbf{45\%} kết nối IoT di động có liên hệ tới xe kết nối.
\item \textbf{2 zettabyte} dữ liệu do IoT tạo ra, chủ yếu từ thiết bị người tiêu dùng.
\end{enumerate}
Các con số này không nên được sử dụng làm mô tả quy mô IoT hiện nay nếu không kiểm chứng lại nguồn.\chonesourcepage{46}
\choneslidefigure{46}{Infographic năm dự báo IoT cho năm 2025 trong tài liệu nguồn}

\section{Cơ hội, thách thức và tổng kết}\label{sec:challenges}
\subsection{Bốn hướng tác động và cơ hội}
Theo slide 47, các thay đổi cơ bản của IoT gồm \textbf{New Markets} -- hình thành thị trường mới; \textbf{Changed Value Creation} -- thay đổi cách tạo ra giá trị và mô hình sản phẩm/dịch vụ; \textbf{Increased Awareness for Information Spaces} -- nhận thức cao hơn về các không gian dữ liệu/thông tin hiện diện cùng môi trường vật lý; và \textbf{Social, Ethical, Legal, and Risk Aspects} -- các khía cạnh xã hội, đạo đức, pháp lý và rủi ro cần được cân nhắc. Hình đám mây từ khóa minh họa các chủ đề phát triển thị trường, dữ liệu, hành vi và công nghệ Internet.\chonesourcepage{47}

\subsection{Những thách thức kỹ thuật và triển khai}
Slide 48 đưa ra chín thách thức; mỗi mục dưới đây giữ nguyên phạm vi kiến thức nguồn và bổ sung diễn giải thuật ngữ tiếng Việt:
\begin{enumerate}
\item \textbf{Availability of the Internet Anywhere Anytime}: khả năng sẵn có của Internet mọi lúc, mọi nơi; bài toán vùng phủ sóng và độ liên tục kết nối.
\item \textbf{Social and Individual Acceptability}: mức chấp nhận của xã hội và cá nhân đối với thiết bị/kết nối, dữ liệu được thu thập và cách công nghệ được ứng dụng.
\item \textbf{Low-cost Smart Sensor Node Development}: phát triển nút cảm biến thông minh có giá thành thấp nhưng phù hợp yêu cầu sử dụng.
\item \textbf{Collect Energy}: thu thập/khai thác năng lượng (energy harvesting) để duy trì hoạt động thiết bị.
\item \textbf{Computing Capability}: năng lực tính toán tại nút hoặc hệ thống khi tài nguyên xử lý có giới hạn.
\item \textbf{Security Issues}: an ninh/bảo mật của thiết bị, mạng và dữ liệu.
\item \textbf{Scalability, Reliability}: khả năng mở rộng và độ tin cậy khi số lượng thiết bị tăng.
\item \textbf{Fault Tolerance}: khả năng chịu lỗi để hệ thống vẫn cung cấp chức năng cần thiết khi một thành phần gặp sự cố.
\item \textbf{Power Consumption of Nodes and Transceivers}: tiêu thụ năng lượng của nút và bộ thu phát truyền thông.
\end{enumerate}
\chonesourcepage{48}

\subsection{Kết luận của chương}
Slide 49 rút ra bốn kết luận chính. Thứ nhất, hệ IoT gồm \emph{cảm biến, các thiết bị kết nối và hạ tầng đám mây lưu trữ lớn}. Thứ hai, IoT phát triển từ sự hội tụ của \emph{công nghệ không dây, công nghệ vi cơ điện (micromechanical) và Internet}. Thứ ba, các công nghệ/giao thức mới liên tục tham gia hệ sinh thái, giúp ứng dụng dễ tiếp cận hơn, có chi phí thấp hơn, sử dụng năng lượng hiệu quả hơn và an toàn hơn. Cuối cùng, nhu cầu ở nhiều lĩnh vực tạo động lực để IoT tiếp tục mở rộng trong tương lai. Slide 50 là lời kết ``Thank you for attention!''.\chonesourcepage{49--50}

\section*{Bảng kiểm đối chiếu nguồn theo trang slide}
\addcontentsline{toc}{section}{Bảng kiểm đối chiếu nguồn theo trang slide}
\small
\begin{longtable}{@{}p{.14\linewidth}p{.79\linewidth}@{}}\toprule
Slide & Nội dung đã ghi nhận trong chương\\\midrule\endfirsthead\toprule Slide & Nội dung đã ghi nhận\\\midrule\endhead
1--2 & Trang bìa, giảng viên, tên chương và sáu nhóm nội dung.\\
3--4 & Mã học phần, mục tiêu, giao thức nêu trong học phần, quy định học và lịch lớp 2022.\\
5, 12 & Định nghĩa IoT (lặp lại) và hình cây thiết bị kết nối.\\
6--8 & Bốn giai đoạn tiến hóa Internet, bảng định nghĩa và sơ đồ Pre-Internet đến IoT.\\
9--11 & Kevin Ashton, RFID, P\&G, bài viết 2009, Neil Gershenfeld, ATM, máy bán hàng 1980 và IIoT.\\
13--14 & Khái niệm mạng IoT, ví dụ đối tượng gắn chip/cảm biến và các dạng thiết bị.\\
15--17 & Cấu trúc Smart Object và năm xu hướng phát triển.\\
18--19 & Smart Environment và Ubiquitous Computing/Mark Weiser.\\
20--22 & M2M, IIoT và các thành phần Industry 4.0.\\
23--28 & Kiến trúc chức năng: Things, Gateway, Cloud, Services, User Interface.\\
29--33 & So sánh kiến trúc 3/4/5 tầng; chi tiết ba tầng và năm nhóm chức năng lớp ứng dụng.\\
34--36 & SoA bốn tầng: Application, Data Processing, Network, Perception.\\
37--39 & Kiến trúc năm tầng, bộ phận mỗi tầng và vai trò Transport, Middleware, Business.\\
40--41 & Kiến trúc sáu và bảy tầng, bảo mật xuyên lớp, phân nhóm Fog/Cloud/Big Data/Business.\\
42--43 & 10 nhóm ứng dụng và tranh ghép tình huống IoT.\\
44--45 & Hai tình huống: báo thức thông minh, chẩn đoán xe kết nối.\\
46 & Năm số liệu dự báo cho 2025 (giữ ngữ cảnh dự báo).\\
47--48 & Cơ hội và chín thách thức nghiên cứu/triển khai IoT.\\
49--50 & Bốn kết luận cuối chương và lời cảm ơn.\\\bottomrule
\end{longtable}
\normalsize
\noindent\textbf{Lưu ý về hình minh họa và quyền nguồn:} Hình ảnh trong bản PDF biên soạn này được tái sử dụng từ file slide mà người dùng cung cấp để phục vụ đối chiếu học tập. Nội dung gốc thuộc tác giả/đơn vị tài liệu đã nêu. Ở bản đưa lên GitHub, có thể tự tạo lại ảnh phục vụ biên dịch bằng \texttt{scripts/extract\_slides.py} từ slide gốc, thay vì lưu trữ hoặc phân phối lại slide gốc trong repo công khai.
% Khi có tài liệu hai chương tiếp theo mới bật hai dòng sau:
% % =============================================================
% CHƯƠNG 2 — PHẦN CỨNG IoT
% Module độc lập: nhập vào tài liệu chính bằng % =============================================================
% CHƯƠNG 2 — PHẦN CỨNG IoT
% Module độc lập: nhập vào tài liệu chính bằng % =============================================================
% CHƯƠNG 2 — PHẦN CỨNG IoT
% Module độc lập: nhập vào tài liệu chính bằng \input{chapter2}
% Nguồn: Chapter 2_IoT Hardware(2).pdf (79 trang)
% Không chứa nội dung Chương 1 hay Chương 3.
% Sử dụng LuaLaTeX; main.tex khai báo fontspec, babel, pdfpages,
% graphicx, tikz, tcolorbox, listings và các macro cho hình.
% =============================================================
\chapter{Phần cứng IoT (IoT Hardware)}\label{chap:iot-hardware}
\begin{concept}
\textbf{Phạm vi Chương 2:} (1) Cảm biến và phân loại cảm biến; (2) cơ cấu chấp hành; (3) nút cảm biến và sơ đồ phần cứng; (4) vi điều khiển ATmega/Arduino cùng kiến thức lập trình; (5) ESP8266 và ESP32; (6) Intel Edison; (7) điện toán biên và Raspberry Pi; (8) so sánh nền tảng. Nội dung được diễn giải sang tiếng Việt và đánh dấu số trang slide gốc nhằm dễ kiểm chứng. Phụ lục cuối chương giữ nguyên \textbf{toàn bộ 79 trang slide}, bao gồm ảnh, sơ đồ chân, ảnh màn hình IDE và bảng biểu.
\end{concept}
\section{Bối cảnh: kiến trúc phần cứng của hệ thống IoT}
Hình ở trang 3 trình bày chuỗi \emph{Things / IoT Devices} (các thiết bị vật lý), \emph{Edge} (điện toán biên), \emph{Network Connectivity} (liên kết mạng) và \emph{Cloud} (đám mây). Về phía phần cứng, thiết bị cảm nhận hoặc tác động lên môi trường; một bo mạch vi điều khiển hay máy tính nhúng thực hiện xử lý; phần giao tiếp mạng trao đổi dữ liệu với biên và đám mây. Các minh họa trong slide đưa ra bo mạch và gateway thuộc nhiều loại khác nhau; không nhất thiết mọi hệ thống phải có tất cả các tầng riêng biệt.
\slidefig{3}{Kiến trúc thiết bị IoT, biên và đám mây}
\section{Cảm biến (Sensors)}
\subsection{Khái niệm và vai trò}
Cảm biến là phần tử cơ bản, thiết yếu của hệ thống IoT, tạo cầu nối giữa thế giới thực và hệ điện tử. Nó đo một đại lượng vật lý, hóa học hoặc sinh học rồi biến đổi thành một tín hiệu mà con người hay thiết bị điện tử có thể đọc. Tín hiệu -- thường là tín hiệu điện -- sau đó được phân tích, xử lý, lưu trữ và truyền đi bằng phương tiện có dây hoặc không dây.

Nếu gọi \(m\) là đại lượng đo và \(s\) là tín hiệu điện đầu ra, quan hệ chuyển đổi được viết:
\begin{equation}\boxed{s=f(m).}\end{equation}
Vì \(f\) mô tả đặc tính chuyển đổi của cảm biến, một phép đo cần quan tâm cả đại lượng đầu vào lẫn cách ánh xạ sang tín hiệu điện; slide không đưa ra hàm hiệu chuẩn cụ thể cho từng cảm biến.\slidesource{4}
\subsection{Các loại cảm biến minh họa}
Hai trang ảnh giới thiệu những ví dụ sau; các tên được chuẩn hóa tiếng Việt và giữ thuật ngữ gốc để thuận tiện tìm linh kiện.
\begin{longtable}{p{.31\linewidth}p{.60\linewidth}}
\toprule\textbf{Tên trong slide} & \textbf{Đại lượng / chức năng tương ứng}\\\midrule\endhead
Temperature sensor & Cảm biến nhiệt độ.\\
Soil moisture sensor & Cảm biến độ ẩm đất.\\
Humidity sensor & Cảm biến độ ẩm môi trường.\\
Infrared sensor; IR optical sensor & Cảm biến hồng ngoại; cảm biến quang hồng ngoại.\\
Gas sensor; alcohol sensor & Cảm biến khí; cảm biến hơi cồn.\\
CO\(_2\) sensor & Cảm biến nồng độ khí carbon dioxide.\\
Digital push button & Nút nhấn số, cung cấp trạng thái đóng/ngắt để làm tín hiệu vào.\\
Ultrasonic sensor & Cảm biến siêu âm.\\
LDR sensor & Quang trở, thay đổi điện trở theo ánh sáng.\\
Gyroscope sensor & Cảm biến tốc độ góc / con quay hồi chuyển.\\
Rain sensor & Cảm biến mưa.\\
Photodiode & Điốt quang.\\
IR proximity sensor; proximity sensor & Cảm biến tiệm cận hồng ngoại; cảm biến tiệm cận nói chung.\\
PIR sensor & Cảm biến hồng ngoại thụ động phát hiện biến thiên bức xạ hồng ngoại.\\\bottomrule
\end{longtable}
\slidesource{5--6}
\slidefig{5}{Ví dụ cảm biến nhiệt độ, độ ẩm, khí, hồng ngoại và nút nhấn}
\slidefig{6}{Bộ sưu tập cảm biến cồn, siêu âm, quang, mưa, tiệm cận, PIR và con quay hồi chuyển}
\subsection{Phân loại theo cách cảm nhận: chủ động và thụ động}
\textbf{Cảm biến tích cực / chủ động (active sensor)} theo cách dùng của bài giảng là loại cần nguồn năng lượng hoặc tín hiệu kích thích từ bên ngoài để thực hiện phép đo. Radar phát sóng vô tuyến, sonar/siêu âm phát sóng âm đến vật thể và nhận phản xạ để xác định khoảng cách. Bài giảng còn liệt kê camera video như ví dụ cảm biến tích cực.\slidesource{7}

\textbf{Cảm biến thụ động (passive sensor)} theo cách trình bày của bài giảng không cần cấp năng lượng kích thích cho chính cơ chế chuyển đổi. \emph{Cặp nhiệt điện} sinh điện áp nhiệt điện khi có chênh lệch nhiệt độ; \emph{cảm biến áp suất} có thể sử dụng màng hoặc vật liệu áp điện tạo ra dịch chuyển/điện áp do tác dụng của áp suất.\slidesource{8}
\begin{noteiot}
Khái niệm chủ động/thụ động ở đây dựa theo cách phân loại trong slide. Khi thiết kế mạch thực tế phải phân biệt \emph{năng lượng kích thích phần tử đo} và \emph{nguồn nuôi toàn bộ mô-đun}; mô-đun có vi mạch xử lý có thể cần cấp nguồn dù phần tử nhạy có cơ chế thụ động.
\end{noteiot}
\subsection{Phân loại theo đại lượng và cơ chế biến đổi}
Ngoài chủ động/thụ động, bài giảng nêu cách phân nhóm theo đối tượng và phương tiện phát hiện (điện, sinh học, hóa học, phóng xạ, v.v.) hoặc hiện tượng chuyển đổi (quang điện, nhiệt điện, sinh học, hóa học, phóng xạ). Các ví dụ theo đại lượng gồm vị trí, nhiệt độ, áp suất, tốc độ, độ ẩm và màu sắc. Đây là các trục phân loại khác nhau, có thể đồng thời mô tả một cảm biến.\slidesource{7--9}
\subsection{Cảm biến tương tự và cảm biến số}
\textbf{Cảm biến tương tự (analog sensor)} tạo tín hiệu có thể biến thiên liên tục trong miền giá trị hoạt động; ví dụ trong slide gồm gia tốc kế, cảm biến ánh sáng, âm thanh, áp suất, nhiệt độ. \textbf{Cảm biến số (digital sensor)} cung cấp giá trị rời rạc hoặc dữ liệu đã số hóa. Bài giảng mô tả cảm biến số được kết hợp với \emph{bộ chuyển đổi tương tự--số} (ADC), còn cảm biến tương tự không tích hợp ADC ở ngõ ra; ví dụ cảm biến gia tốc số, nhiệt độ số DS1620 và HYT~271, cảm biến độ ẩm số. Slide lưu ý nhóm cảm biến số có thể bị giới hạn dải tính toán/đo, đồng thời đưa ví dụ về cách diễn giải đầu ra nhiệt độ.

Phân biệt thêm \emph{tín hiệu số hai mức} (logic HIGH/LOW, ví dụ nút nhấn) với \emph{dữ liệu số nhiều bit} (ví dụ nhiệt độ đọc qua giao tiếp số). Trên slide 10, hình tín hiệu số minh họa các xung nhị phân mức 0/1, còn hình tương tự biểu diễn sóng biến thiên liên tục.\slidesource{9--12}
\slidefig{10}{So sánh tín hiệu tương tự liên tục và tín hiệu nhị phân}
\section{Cơ cấu chấp hành (Actuators)}
\subsection{Định nghĩa và ví dụ}
Cơ cấu chấp hành nhận tín hiệu điều khiển -- chẳng hạn tín hiệu điện từ một bộ điều khiển -- và biến đổi thành tác động đối với vật lý bên ngoài: chuyển động cơ học, biến thiên nhiệt độ, ánh sáng hoặc các dạng năng lượng khác. Bài giảng đưa ví dụ \emph{động cơ bước}: động cơ quay từng góc xác định tương ứng chuỗi xung điều khiển; động cơ bước thích hợp khi yêu cầu điều khiển chính xác vị trí, như cánh tay rô-bốt. Các hình còn cho thấy động cơ DC, màn hình LCD, LED, cơ cấu tuyến tính, ma trận LED và còi/buzzer, qua đó thể hiện cơ cấu chấp hành không chỉ giới hạn ở động cơ quay.\slidesource{13--14}
\slidefig{14}{Ví dụ các cơ cấu chấp hành điện--cơ, hiển thị và âm thanh}
\subsection{So sánh cảm biến và cơ cấu chấp hành}
\begin{center}\begin{tabularx}{\linewidth}{>{\bfseries}p{.22\linewidth}XX}\toprule & Cảm biến & Cơ cấu chấp hành\\\midrule
Chức năng & Chuyển đại lượng vật lý thành tín hiệu đo/điện & Chuyển lệnh hoặc tín hiệu điện thành tác động vật lý\\
Vị trí trong điều khiển & Đầu vào của bộ điều khiển & Đầu ra của bộ điều khiển\\
Ví dụ & Công tắc, micro, cảm biến quay & LED, loa/buzzer, bộ điều khiển động cơ, laser\\\bottomrule\end{tabularx}\end{center}
Một vòng lặp điều khiển điển hình: \emph{môi trường} \(\rightarrow\) \emph{cảm biến} \(\rightarrow\) \emph{bộ điều khiển} \(\rightarrow\) \emph{cơ cấu chấp hành} \(\rightarrow\) \emph{môi trường}.\slidesource{15}
\section{Nút cảm biến IoT (IoT Sensor Node)}
Trong trang 16, bài giảng minh họa cách ghép nhóm \emph{Sensors + Boards + Connectivity} thành \emph{IoT Sensor Nodes}; có ví dụ cảm biến, các bo xử lý, mô-đun truyền thông và bo lắp ráp cuối cùng. Trang 17 đưa sơ đồ khối chi tiết: cảm biến thông minh, cảm biến số, cảm biến video/môi trường; chuyển đổi ADC; MCU; các bộ nhớ flash/RAM; cấp nguồn, DC/DC, quản lý pin; mô-đun truyền thông không dây (Wi-Fi, Bluetooth, IEEE 802.15.4) hoặc có dây. Sơ đồ chỉ ra MCU đóng vai trò phối hợp lấy mẫu, giao tiếp, xử lý và truyền dữ liệu.
\begin{center}\begin{tikzpicture}[node distance=8mm and 7mm, box/.style={draw=navy,rounded corners,fill=blue!4,align=center,minimum height=11mm,minimum width=21mm,font=\footnotesize}, >={Latex}]
\node[box] (s) {Cảm biến};\node[box,right=of s] (adc) {Điều hòa\ tín hiệu / ADC};\node[box,right=of adc] (m) {MCU / CPU};\node[box,right=of m] (n) {Kết nối\ mạng};\node[box,below=of m] (mem) {RAM / Flash};\node[box,below=of adc] (p) {Nguồn, pin,\ DC/DC};
\draw[->] (s)--(adc);\draw[->](adc)--(m);\draw[<->](m)--(n);\draw[<->](m)--(mem);\draw[->](p)--(adc);\draw[->](p.east) -| (m.south);
\end{tikzpicture}\end{center}
\slidefig{17}{Sơ đồ khối nguyên bản của một nút cảm biến IoT}
\section{Vi điều khiển và ATmega328P}
\subsection{Khái niệm microcontroller}
\emph{Microcontroller} (MCU) là một vi mạch tích hợp (IC), có thể hiểu như máy tính thu nhỏ gồm bộ xử lý, bộ nhớ và khối vào/ra. MCU có thể lập trình và được thiết kế cho các ứng dụng nhúng. Trang 18 minh họa các vi mạch Intel P8051, Hitachi HD68000, Motorola MC68HC000LC8. Cần phân biệt \emph{chip MCU} với \emph{bo phát triển} chứa MCU cùng mạch nguồn, giao tiếp và chân cắm.\slidesource{18}
\subsection{Đặc điểm của ATmega328P}
\begin{longtable}{p{.46\linewidth}p{.45\linewidth}}\toprule\textbf{Thông số được nêu trong slide} & \textbf{Giá trị}\\\midrule\endhead
Họ vi điều khiển & AVR 8 bit, hiệu năng cao, công suất thấp\\
Flash chương trình & 32 KB\\
SRAM & 2 KB\\
EEPROM & 1 KB\\
Đường I/O có thể lập trình & 23\\
Điện áp hoạt động & 1,8--5,5 V\\
Dải nhiệt độ & $-40$ đến $85\,^{\circ}\mathrm C$\\
Dòng chế độ Active & 0,2 mA (điều kiện cụ thể không được slide nêu)\\
Dòng chế độ Power-down & 0,1 $\mu$A (điều kiện cụ thể không được slide nêu)\\
Dòng chế độ Power-save & 0,75 $\mu$A (điều kiện cụ thể không được slide nêu)\\\bottomrule\end{longtable}
Hình trang 19 trình bày sơ đồ chân ATmega328P dạng DIP-28, với các nhóm chân: cấp nguồn VCC/AVCC/GND, RESET, dao động XTAL, chân PORTB/PORTC/PORTD, ADC, UART, SPI, I\(^{2}\)C và các chức năng ghép kênh. Trang 21 cho ví dụ mạch lập trình/chuyển chương trình vào chip, đầu nối và bảng pin của giao thức ISP/SPI.\slidesource{19--21}
\slidefig{19}{Sơ đồ 28 chân ATmega328P và các chức năng ghép kênh}
\section{Nền tảng Arduino}
\subsection{Khái quát và lịch sử}
Arduino là nền tảng điện tử \emph{mã nguồn mở} bao gồm cả phần cứng và phần mềm, hướng tới tính dễ dùng. Theo bài giảng, dự án hình thành năm 2003 tại Interaction Design Institute Ivrea (Ý) cho sinh viên; tên gọi bắt nguồn từ quán bar nơi nhóm sáng lập gặp nhau, quán được đặt theo Arduin of Ivrea. Phần cứng ban đầu do Smart Projects tại Ý sản xuất; một số bo mang thương hiệu Arduino được thiết kế bởi các đơn vị khác như SparkFun và Adafruit.\slidesource{22}

Ưu điểm được liệt kê gồm giá thành thấp, đa nền tảng, môi trường lập trình đơn giản, phần mềm mở và phần cứng mở. Các hình giới thiệu họ bo Arduino có nhiều kiểu dáng/kích thước, không phải một bộ xử lý đơn nhất.\slidesource{23--24,31}
\subsection{Arduino Uno và sơ đồ chức năng}
Arduino Uno sử dụng \emph{ATmega328P}, được Arduino.cc phát triển; có 14 chân I/O số và 6 đầu vào analog. Bo có thể lập trình trong Arduino IDE qua cáp USB Type-B; bootloader được nạp sẵn giúp tải chương trình mà không cần thiết bị nạp phần cứng ngoài.

Mười khu vực được chỉ ra trên ảnh bo Uno: (1) các chân digital; (2) LED trên pin 13; (3) LED báo nguồn; (4) ATmega; (5) chân analog; (6) GND và 5V; (7) đầu nối nguồn; (8) LED TX/RX; (9) cổng USB; (10) nút RESET. Sơ đồ pinout ở trang 27 minh họa quan hệ \emph{chân vật lý ATmega328P -- ký hiệu Port -- chức năng Arduino -- đầu vào ADC và các chức năng ngoại vi}.\slidesource{25--27}
\slidefig{26}{Sơ đồ mười khu vực trên bo Arduino Uno}
\slidefig{27}{Bản đồ chân ATmega328P trong hệ Arduino}
\subsection{Các mẫu bo: Leonardo, Mega, Nano}
\begin{longtable}{>{\raggedright\arraybackslash}p{.28\linewidth}p{.20\linewidth}p{.22\linewidth}p{.20\linewidth}}\toprule\textbf{Thông số} & \textbf{Leonardo} & \textbf{Mega} & \textbf{Nano}\\\midrule\endhead
Vi điều khiển & ATmega32U4 & ATmega2560 & ATmega328\\
Điện áp hoạt động & 5 V & 5 V & 5 V\\
Điện áp vào khuyến nghị & 7--12 V & 7--12 V & 7--12 V\\
Giới hạn điện áp vào & 6--20 V & 6--20 V & Không ghi trong slide\\
I/O số & 23 & 54 & 22 (6 PWM)\\
Kênh analog & 12 & 16 & 8\\
Flash & 32 KB, trong đó 4 KB cho bootloader & 256 KB & 32 KB\\
SRAM & 2,5 KB & 8 KB & 2 KB\\
EEPROM & 1 KB & 4 KB & 1 KB\\
Dòng I/O nêu trong slide & 40 mA & 40 mA & Không ghi\\
Dòng chân 3,3 V nêu trong slide & 50 mA & 50 mA & Không ghi\\
Kích thước & 68,6 $\times$ 53,3 mm & Không ghi & 18 $\times$ 45 mm\\
Khối lượng & 20 g & Không ghi & 7 g\\
Thông số bổ sung & -- & -- & 19 mA (tiêu thụ theo slide)\\\bottomrule\end{longtable}
Các mức dòng trong bảng được tái hiện theo slide để phục vụ đối chiếu, không thay thế cho giới hạn thiết kế điện và khuyến cáo của datasheet.\slidesource{28--30}
\subsection{Các lựa chọn ngoài Arduino và shield mở rộng}
Trang 32--34 trình bày \emph{LaunchPad MSP430, Particle Photon, NodeMCU ESP8266, SparkFun Thing Plus, Adafruit Feather HUZZAH, BeagleBoard/PocketBeagle, Seeeduino Nano, SparkFun RedBoard Artemis, STM32F3 Discovery}. Đây là các nền tảng hoặc bo có thể dùng cho nguyên mẫu IoT, không nhất thiết tương thích về phần cứng hay công cụ lập trình.

\emph{Arduino Shield} là bo cắm chồng lên PCB Arduino để bổ sung chức năng, được thiết kế theo hệ chân và triết lý dễ lắp, dễ tháo, dễ sản xuất chi phí thấp. Hình minh họa cho thấy các shield hiển thị, kết nối, nguyên mẫu và mở rộng giao tiếp. Bộ \emph{Arduino Starter Kit} trong trang 36 tập hợp bo mạch, breadboard, linh kiện rời, dây cắm và phụ kiện điện tử.\slidesource{32--36}
\slidefig{34}{Các bo phát triển thay thế được bài giảng liệt kê}
\slidefig{35}{Các dạng Arduino Shield mở rộng}
\subsection{Breadboard và nguyên mẫu}
\emph{Breadboard} là bảng cắm mạch thử nghiệm không cần hàn, thích hợp kiểm chứng nhanh mạch điện và thiết kế trước khi chế tạo PCB. Trang 37 minh họa Arduino nối một LED và điện trở trên breadboard; cần kiểm tra hướng kết nối, hàng/cột chân nối thông và điện trở giới hạn dòng.\slidesource{37}
\slidefig{37}{Mô hình thử nghiệm Arduino với breadboard}
\section{Môi trường và ngôn ngữ lập trình Arduino}
\subsection{Arduino Sketch và Arduino Desktop IDE}
\emph{Sketch} là tên Arduino dùng cho một chương trình, tức đơn vị mã nguồn được biên dịch, nạp và chạy trên bo. Ngôn ngữ Arduino dựa trên các hàm C/C++. Trang 38 mô tả một sketch mẫu và cách tổ chức chương trình thành \texttt{setup()} chạy khởi tạo và \texttt{loop()} thực thi lặp.

Arduino Desktop IDE hỗ trợ Windows, macOS, Linux. Trang 39 dẫn các hướng dẫn cài đặt của Arduino cho từng hệ điều hành, kèm lưu ý cách cài bằng \texttt{sudo apt-get install arduino} trên Linux không được bài giảng khuyến nghị. Trang 40 giới thiệu giao diện IDE, nơi viết sketch, xác minh/biên dịch, chọn bo/cổng giao tiếp, nạp và theo dõi cửa sổ Serial Monitor.\slidesource{38--40}
\slidefig{40}{Giao diện Arduino IDE trong slide gốc}
\subsection{Chú thích chương trình (comments)}
Chú thích được trình biên dịch bỏ qua, dùng để giải thích ý nghĩa mã lệnh. Trang 41 minh họa hai dạng:
\begin{lstlisting}[language=C++]
/* Day la chu thich nhieu dong.
   Cac dong nay duoc bo qua. */
// Day la chu thich mot dong.
\end{lstlisting}
\noindent Lưu ý: chú thích nhiều dòng theo cú pháp \texttt{/* ... */} không được viết lồng nhau; dấu \texttt{//} dùng cho chú thích một dòng.
\subsection{Cấu trúc điều khiển: for, if, while, do--while}
\paragraph{Vòng lặp for (slide 42).} Lặp khối câu lệnh giữa ngoặc nhọn với bộ đếm thường gồm ba biểu thức khởi tạo, điều kiện, bước cập nhật. Biến đếm có thể tăng hoặc giảm.
\begin{lstlisting}[language=C++]
for (int i = 0; i <= 10; i++) {
  Serial.write(7);
}
\end{lstlisting}
\paragraph{Rẽ nhánh if (slide 43).} Kiểm tra biểu thức điều kiện; nếu đúng thì thực hiện nhánh \texttt{if}, nếu sai có thể thực hiện \texttt{else}. Bài giảng minh họa kiểm tra nhiệt độ so với ngưỡng \(70\), đồng thời liệt kê các toán tử \texttt{==}, \texttt{!=}, \texttt{<}, \texttt{<=}, \texttt{>}, \texttt{>=}. Cần phân biệt \texttt{=} (gán) và \texttt{==} (so sánh bằng).
\begin{lstlisting}[language=C++]
if (temperature > 70) {
  // Xu ly khi cao hon nguong
} else {
  // Xu ly khi khong cao hon nguong
}
\end{lstlisting}
\paragraph{Vòng lặp while (slide 44).} Lặp khi điều kiện vẫn đúng; kiểm tra điều kiện trước vòng lặp. \texttt{do\{...\}while(...);} thực hiện thân lặp ít nhất một lần rồi mới kiểm tra điều kiện.
\begin{lstlisting}[language=C++]
int var = 0;
while (var < 200) { var++; }

do { var++; } while (var < 400);
\end{lstlisting}
\slidesource{41--44}
\subsection{Hàm đọc và ghi tín hiệu số}
\texttt{pinMode(pin, mode)} thiết lập chân làm đầu vào hoặc đầu ra (trong slide: \texttt{INPUT}, \texttt{OUTPUT}). \texttt{digitalRead(pin)} đọc trạng thái logic \texttt{HIGH}/\texttt{LOW}. \texttt{digitalWrite(pin, value)} xuất logic \texttt{HIGH} hoặc \texttt{LOW} ra một chân đã được cấu hình phù hợp. Tham số \texttt{pin} là số hiệu chân digital cần tác động, \texttt{value} là mức logic yêu cầu.\slidesource{45}
\begin{lstlisting}[language=C++]
const int LED_PIN = 13;
void setup() { pinMode(LED_PIN, OUTPUT); }
void loop() {
  digitalWrite(LED_PIN, HIGH);
  delay(1000);
  digitalWrite(LED_PIN, LOW);
  delay(1000);
}
\end{lstlisting}
\subsection{Tín hiệu tương tự và PWM}
\texttt{analogRead(pin)} đọc đầu vào ADC (slide mô tả bộ chuyển đổi 10 bit). \texttt{analogWrite(pin, value)} xuất \emph{điều chế độ rộng xung} PWM lên chân hỗ trợ; giá trị duty trong ví dụ Arduino nằm trong khoảng 0--255, tương ứng thời gian xung ở mức cao trong mỗi chu kỳ, chứ không phải một đầu ra điện áp DAC lý tưởng. Hình trang 46 biểu diễn dạng sóng PWM có độ rộng khác nhau.\slidesource{46}
\slidefig{46}{Sơ đồ minh họa nguyên lý PWM và các giá trị analog của Arduino}
\subsection{Các hàm toán học thông dụng}
\begin{longtable}{p{.26\linewidth}p{.64\linewidth}}\toprule\textbf{Hàm} & \textbf{Tác dụng theo bài giảng}\\\midrule\endhead
\texttt{abs(x)} & Trả về trị tuyệt đối.\\
\texttt{constrain(x,a,b)} & Giới hạn \(x\) trong đoạn từ \(a\) đến \(b\).\\
\texttt{max(x,y)} & Chọn giá trị lớn nhất.\\
\texttt{min(x,y)} & Chọn giá trị nhỏ nhất.\\
\texttt{pow(base,exponent)} & Lũy thừa \(\mathrm{base}^{\mathrm{exponent}}\).\\
\texttt{sq(x)} & Bình phương \(x\).\\
\texttt{sqrt(x)} & Căn bậc hai \(\sqrt{x}\).\\\bottomrule\end{longtable}
\slidesource{47}
\subsection{Trì hoãn và chỉ thị tiền xử lý}
\texttt{delay(milliseconds)} tạm dừng chương trình trong khoảng thời gian tính theo mili-giây, \(1000\,\mathrm{ms}=1\,\mathrm{s}\). \texttt{\#define} đặt tên hằng hoặc macro trước khi chương trình được biên dịch, ví dụ \texttt{\#define ledPin 3}. \texttt{\#include} thêm tiêu đề thư viện, ví dụ \texttt{\#include <SPI.h>}; nhờ đó sketch truy cập các hàm, lớp và định nghĩa đã được thư viện cung cấp.\slidesource{48--49}
\begin{lstlisting}[language=C++]
#define ledPin 3
#include <SPI.h>
\end{lstlisting}
\subsection{Hai hàm chính setup() và loop()}
\texttt{setup()} được gọi một lần mỗi khi khởi động hoặc reset để thiết lập chân, Serial, thư viện và các trạng thái ban đầu. \texttt{loop()} được gọi lặp liên tục sau khi \texttt{setup()} hoàn tất, là nơi đặt chương trình chính xử lý tín hiệu và điều khiển thiết bị. Slide 50 dùng \texttt{pinMode(8, INPUT)} minh họa thao tác khởi tạo; slide 51 minh họa dùng \texttt{digitalRead(8)} và ghi chữ qua Serial theo trạng thái nút nhấn.\slidesource{50--51}
\begin{lstlisting}[language=C++]
void setup() {
  pinMode(8, INPUT);
  Serial.begin(9600);
}
void loop() {
  if (digitalRead(8) == HIGH) {
    Serial.write('H');
  } else {
    Serial.write('L');
  }
}
\end{lstlisting}
\subsection{Thêm thư viện và giao tiếp Serial}
Trang 52 minh họa thao tác \emph{Sketch > Include Library} trong Arduino IDE. Các thư viện có thể đi kèm IDE hoặc được thêm bởi người dùng. Cổng COM là tên truyền thống của cổng giao tiếp nối tiếp trên máy tính tương thích IBM PC, ngày nay còn bao gồm cổng nối tiếp ảo thông qua Bluetooth hoặc cầu USB-to-UART. Arduino thường cung cấp ít nhất một cổng nối tiếp, một số bo có nhiều cổng. Trên Arduino Uno, \textbf{pin 0 là RX và pin 1 là TX}; người lập trình thường dùng \texttt{Serial.begin(speed)} khởi tạo với tốc độ baud và \texttt{Serial.write(value)} ghi dữ liệu ra cổng; baud là số ký hiệu truyền mỗi giây (bài giảng mô tả gần đúng theo bit/giây).\slidesource{52--53}
\slidefig{52}{Thao tác thêm thư viện trong Arduino IDE}
\section{Nền tảng Wi-Fi ESP8266}
\subsection{Tổng quan ESP8266EX và khả năng phần mềm}
ESP8266EX là chip Wi-Fi chi phí thấp của Espressif Systems, tích hợp giao thức TCP/IP và có thể ghép vào các thiết kế dùng vi điều khiển. Phần lõi hỗ trợ Arduino của ESP8266 cung cấp thư viện Wi-Fi cho TCP, UDP, HTTP, mDNS, SSDP, máy chủ DNS; cho phép cập nhật firmware qua mạng (OTA), hệ thống tệp trên flash hoặc SD, điều khiển servo và trao đổi SPI/I\(^{2}\)C. Slide 55 liệt kê các biến thể mô-đun ESP-01, ESP-02, ESP-03, ESP-04, ESP-05, ESP-06, ESP-07, ESP-08, ESP-09, ESP-10, ESP-11, ESP-12, ESP-13.\slidesource{54--55}
\slidefig{55}{Tập hợp biến thể mô-đun ESP8266}
\subsection{ESP-01 và ESP-07}
\begin{longtable}{p{.34\linewidth}p{.57\linewidth}}\toprule\textbf{Mô-đun} & \textbf{Đặc điểm được cung cấp trong slide}\\\midrule\endhead
ESP-01 & Kích thước $24{,}75\times14{,}5$ mm; nguồn 3,3 V; anten mạch in tích hợp; LED nguồn và TXD; chế độ AP, STA, AP+STA; có thể điều khiển bằng lệnh AT; khoảng cách chân 2,54 mm.\\
ESP-07 & Nguồn 3,3 V; anten gốm tích hợp và đầu nối anten ngoài U.FL; Wi-Fi 802.11b/g/n băng 2,4 GHz với WPA/WPA2; TCP/UDP; hỗ trợ lệnh AT; khoảng cách chân 2 mm.\\\bottomrule\end{longtable}
\slidesource{56--57}
\slidefig{56}{Sơ đồ chân, module ESP-01 và đặc điểm ngoại hình}
\slidefig{57}{Sơ đồ chân ESP-07 trong bài giảng}
\subsection{SDK của Espressif và môi trường thay thế}
Bài giảng nêu năm 2014 Espressif giới thiệu SDK cho phép ESP8266 chạy như vi điều khiển độc lập thay vì chỉ là mô-đun Wi-Fi của MCU khác. Hai nhánh SDK được nhắc đến là \emph{ESP8266 NonOS SDK} và \emph{ESP8266 RTOS SDK}. Các môi trường hoặc SDK thay thế trên slide gồm ESP-Open-SDK, Unofficial Development Kit, Arduino, ESP8266 BASIC, ESP Easy và ESPHome.\slidesource{58}
\subsection{Cấu hình ESP8266 trên Arduino IDE}
Ba trang hướng dẫn tuần tự: \textbf{(1)} vào \emph{File > Preferences} và thêm URL gói bo \url{http://arduino.esp8266.com/stable/package_esp8266com_index.json}; \textbf{(2)} vào \emph{Tools > Board > Boards Manager}, tìm ``esp8266'' và chọn \emph{Install}; \textbf{(3)} vào \emph{Tools > Board} chọn \emph{Generic ESP8266 Module} nếu đúng loại mô-đun cần sử dụng. Các ảnh minh họa trong slide là giao diện IDE tại thời điểm biên soạn; vị trí menu có thể khác ở bản IDE mới.\slidesource{59--61}
\slidefig{59}{Thêm liên kết JSON của gói bo ESP8266}
\slidefig{60}{Cài gói hỗ trợ ESP8266 qua Boards Manager}
\slidefig{61}{Chọn bo Generic ESP8266 Module}
\section{Nền tảng ESP32}
ESP32 là họ \emph{system-on-chip} giá tương đối thấp do Espressif phát triển, là bước tiến so với họ ESP8266. Slide giới thiệu khả năng đa lõi, đồng xử lý tiết kiệm năng lượng và hỗ trợ kết nối Wi-Fi, Bluetooth. Những thiết bị minh họa gồm ESP32 Development Board, ESP32-CAM, ESP32-WROOM-32U/-32D, DOIT DevKit V1, ESP32 DevKit, ESP32S NodeMCU, ESP32 Thing, Wemos Lolin32, Wemos OLED và HUZZAH32.

Các \textbf{ưu điểm theo slide}: CPU dual-core khoảng 160--240 MHz; Wi-Fi/Bluetooth tiện cho điều khiển và giám sát; số chân GPIO nhiều hơn ESP8266; slide ghi tốc độ tối đa 150 Mb/s. \textbf{Nhược điểm theo slide}: giá cao hơn ESP8266. Thông số là mô tả cho các phiên bản ESP32 được trình bày trong bài giảng, không đại diện cho mọi biến thể ESP32. Đặc biệt, ảnh trang 64 minh họa pinout bo DOIT ESP32 DevKit V1 loại 30 chân, gồm các chân cấp nguồn, EN, chân UART, ADC, digital GPIO, SPI và I\(^{2}\)C tùy cấu hình.\slidesource{62--65}
\slidefig{63}{Một số bo ESP32 được giới thiệu trong slide}
\slidefig{64}{Sơ đồ chân bo DOIT ESP32 DevKit V1 (30 GPIOs theo tiêu đề slide)}
\section{Nền tảng Intel Edison}
\subsection{Cấu hình mô-đun Edison}
Intel Edison trong bài giảng có kích thước $35{,}5\times25\times3{,}9$ mm, linh kiện gắn ở hai mặt. SoC Intel Atom Tangier Z34XX tiến trình 22 nm tích hợp hai lõi Silvermont 500 MHz và một lõi Intel Quark ở 100 MHz (đơn vị trong slide không ghi đầy đủ tại dòng này); RAM 1 GB và eMMC 4 GB; có Wi-Fi, Bluetooth 4, điều khiển USB. Kết nối ra ngoài qua đầu nối Hirose DF40 mật độ cao 70 chân (USB, SD, UART, GPIO). Các môi trường lập trình nêu trong slide: Arduino IDE, Eclipse (C, C++, Python), Intel XDK (NodeJS/HTML5).\slidesource{66}
\slidefig{66}{Mô-đun Intel Edison và cấu hình phần cứng}
\subsection{Bo breakout và bo dạng Arduino cho Edison}
\textbf{Breakout board}: xuất các mức I/O gốc 1,8 V của Edison ra các điểm hàn dạng lưới 0,1 inch; có USB OTG Micro Type-AB, công tắc nguồn OTG, bộ sạc pin, cầu USB--UART qua Micro-USB, đầu cấp nguồn DC 7--15 V. Điện áp I/O 1,8 V là thông tin quan trọng khi giao tiếp với ngoại vi.

\textbf{Arduino board for Edison Module}: 20 I/O số (6 PWM), 6 đầu vào analog, 1 UART, 1 I\(^{2}\)C, đầu ICSP 6 chân cho SPI, cổng Micro-USB device hoặc cổng USB host Type-A chọn qua công tắc cơ, cổng Micro-USB UART, khe SD và đầu nguồn DC 7--15 V.\slidesource{67--68}
\slidefig{67}{Intel Edison Breakout Board}
\slidefig{68}{Bo tương thích bố cục Arduino cho Edison}
\section{Điện toán biên (Edge Computing)}
Slide 69 trình bày sự khác nhau về vị trí xử lý giữa đám mây, trung tâm dữ liệu doanh nghiệp và biên gần thiết bị; trang 70 nhấn mạnh \emph{xử lý thông tin nhạy cảm với độ trễ}, \emph{quyền riêng tư và an ninh}, \emph{băng thông và chất lượng kết nối}. Trang 71 mô tả dữ liệu do thiết bị IoT tạo ra được xử lý một phần ở máy hoặc gateway gần nguồn, trong khi đám mây và trung tâm dữ liệu tiếp tục thực hiện các công việc ở quy mô lớn. Khác với gửi tất cả dữ liệu thô lên cloud, cách bố trí này có thể giảm lượng dữ liệu truyền và cho phép đưa ra quyết định cục bộ.

Các \textbf{thiết bị biên} minh họa trang 72 gồm Raspberry Pi, NVIDIA Jetson Nano và NVIDIA Jetson Xavier. Liên kết tham khảo trên slide: \url{https://www.networkworld.com/article/3224893/what-is-edge-computing-and-how-it-s-changing-the-network.html}.\slidesource{69--72}
\slidefig{69}{Quan hệ giữa cloud, edge gateway và trung tâm dữ liệu}
\slidefig{71}{Sơ đồ hoạt động của điện toán biên}
\slidefig{72}{Các máy tính biên: Raspberry Pi và NVIDIA Jetson}
\section{Máy tính đơn bo Raspberry Pi}
\subsection{Khái quát}
Raspberry Pi là dòng \emph{single-board computers} (SBC), tức máy tính tích hợp phần lớn thành phần trên một bo mạch nhỏ gọn. Theo số liệu lịch sử \textbf{được trích trong slide}, đến tháng 2/2015 hơn 5 triệu máy Raspberry Pi đã được bán, trở thành máy tính Anh bán chạy nhất thời điểm đó. Dòng bo do Raspberry Pi Foundation phát triển nhằm khuyến khích học khoa học máy tính cơ bản trong trường học và tại những quốc gia đang phát triển. Slide nêu các bo có Wi-Fi/Bluetooth tích hợp; khả năng này tùy theo phiên bản thực tế. Hình trang 73 minh họa các dự án robot và bộ phụ kiện Raspberry Pi.\slidesource{73--74}
\slidefig{73}{Hệ sinh thái thiết bị và ứng dụng Raspberry Pi}
\subsection{Bộ kit nhập môn}
Trang 75 liệt kê các thành phần: bo Raspberry Pi, thẻ microSD dung lượng \textbf{8 GB trở lên} theo slide, bộ cấp nguồn, bàn phím/chuột USB, màn hình có ngõ HDMI, cáp HDMI và hộp vỏ bảo vệ. Thẻ microSD làm bộ nhớ khởi động và lưu hệ điều hành đối với các cấu hình được minh họa trong bài giảng.\slidesource{75}
\slidefig{75}{Bộ kit Raspberry Pi gồm bo, microSD, nguồn và ngoại vi}
\subsection{Cấu tạo và chân Raspberry Pi 4}
Trang 76 chỉ các phần \textbf{USB-C power}, \textbf{micro-HDMI} (hai đầu ra), \textbf{Ethernet Gigabit}, cổng \textbf{USB 3.0/2.0}, header GPIO 40 chân và lựa chọn RAM 2/4/8 GB trong hình. Trang 77 chứa bảng pinout 40 chân: nguồn 3,3 V, 5 V, GND, các chân GPIO và các chân chức năng SPI/UART/I\(^{2}\)C. Khi đấu nối phải đối chiếu số thứ tự chân vật lý (physical pin) và tên GPIO/BCM, tránh cấp 5 V vào tín hiệu GPIO không chịu 5 V.\slidesource{76--77}
\slidefig{76}{Các cổng kết nối và khu vực chính của Raspberry Pi 4}
\slidefig{77}{Sơ đồ chân header 40 pin Raspberry Pi 4}
\section{So sánh Arduino Uno và Raspberry Pi 4}
\begin{longtable}{p{.38\linewidth}p{.26\linewidth}p{.26\linewidth}}\toprule\textbf{Tiêu chí trong bài giảng} & \textbf{Arduino Uno} & \textbf{Raspberry Pi 4}\\\midrule\endhead
Loại nền tảng & Vi điều khiển (microcontroller) & Máy vi tính (microcomputer)\\
Hệ điều hành & Không chạy hệ điều hành đa nhiệm thông dụng & Có thể chạy hệ điều hành\\
Đa nhiệm & Slide ghi ``No multitasking'' & Có đa nhiệm\\
Giá thành tương đối & Rẻ & Đắt hơn\\
Tài nguyên & Xung 16 MHz, SRAM 2 KB & Xung khoảng 1,5 GHz, RAM 2/4/8 GB\\
Ứng dụng thời gian thực & Được slide đánh giá là thích hợp & Slide đánh giá không lý tưởng cho thời gian thực cứng\\
Kết nối mạng tích hợp & Slide ghi không có kết nối mạng & Có khả năng kết nối mạng\\
ADC tích hợp trên bo & Có & Slide ghi không có bộ chuyển đổi A/D\\\bottomrule\end{longtable}
Dòng so sánh phản ánh \textbf{đúng hai bo được nêu trong slide 78}, không phải khẳng định cho toàn bộ họ Arduino và Raspberry Pi. Một số chức năng có thể mở rộng bằng shield, HAT hoặc IC ngoài; với Raspberry Pi, yêu cầu đo analog thường cần thêm ADC bên ngoài.\slidesource{78}
\slidefig{78}{Bảng so sánh Arduino Uno và Raspberry Pi 4 nguyên bản}
\section{Bảng tra cứu tương ứng với toàn bộ 79 slide}
\begin{longtable}{p{.15\linewidth}p{.76\linewidth}}\toprule\textbf{Trang PDF} & \textbf{Chủ đề và hình minh họa được lưu trong phụ lục gốc}\\\midrule\endhead
1--3 & Bìa giảng viên, mục lục chương và kiến trúc Cloud--Edge--IoT Devices.\\
4--6 & Khái niệm cảm biến, hàm chuyển đổi, ảnh mô-đun cảm biến đa dạng.\\
7--12 & Chủ động/thụ động; phân loại theo hiện tượng và analog/digital; dạng sóng; cảm biến tương tự và cảm biến số.\\
13--15 & Cơ cấu chấp hành, hình các thiết bị đầu ra, so sánh cảm biến/cơ cấu chấp hành.\\
16--17 & Phân rã phần cứng nút IoT và sơ đồ khối IoT Sensor Node.\\
18--21 & Khái niệm MCU, chip ATmega328P, cấu hình và bộ nạp.\\
22--27 & Lịch sử/ưu điểm Arduino, bo Uno, các vùng chức năng và pinout ATmega328P.\\
28--31 & Arduino Leonardo, Mega, Nano và ảnh các biến thể Arduino.\\
32--35 & Bo thay thế Arduino và shield mở rộng.\\
36--40 & Starter Kit, breadboard, Arduino Sketch, cài đặt và giao diện Arduino IDE.\\
41--44 & Comment và cấu trúc \texttt{for}, \texttt{if}, \texttt{while}, \texttt{do--while}.\\
45--48 & Digital I/O, Analog I/O và PWM, hàm toán học và delay.\\
49--53 & \texttt{\#define}, \texttt{\#include}, \texttt{setup()}, \texttt{loop()}, thư viện và Serial.\\
54--58 & ESP8266, các biến thể mô-đun, đặc điểm ESP-01, ESP-07, SDK.\\
59--61 & Quy trình cài đặt board ESP8266 trong Arduino IDE với ảnh các bước.\\
62--65 & ESP32, các bo ESP32, sơ đồ chân DOIT và ưu/nhược điểm.\\
66--68 & Intel Edison, bo breakout và bo giao tiếp phong cách Arduino.\\
69--72 & Điện toán biên, mô hình hoạt động và các bo edge devices.\\
73--77 & Raspberry Pi, bộ kit, sơ đồ các cổng và pinout Pi 4.\\
78--79 & So sánh Uno và Pi 4, lời kết.\\\bottomrule\end{longtable}
\section{Kết luận Chương 2}
Phần cứng IoT được tổ chức xung quanh \emph{cảm biến} (thu nhận dữ liệu), \emph{cơ cấu chấp hành} (tạo tác động), \emph{bộ điều khiển/tính toán} (quản lý, xử lý, điều khiển), \emph{giao diện I/O}, \emph{khối nguồn} và \emph{khối truyền thông}. Với các tác vụ điều khiển nhúng, Arduino/ATmega và ESP8266/ESP32 là những ví dụ tiêu biểu của MCU. Với các tác vụ cần hệ điều hành, đa nhiệm, xử lý ở biên và giao tiếp phần mềm đa dạng, Raspberry Pi hoặc một bo công nghiệp có thể phù hợp hơn. Intel Edison được bài giảng dùng làm minh họa một kiến trúc mô-đun máy tính nhúng kết hợp nhiều giao diện ngoại vi.

\textbf{Điểm cần ghi nhớ:} lựa chọn không chỉ theo xung CPU, mà theo số chân, điện áp I/O, ADC, dung lượng RAM/Flash, giao tiếp không dây, nguồn, phần mềm công cụ và độ trễ xử lý mong muốn. Các ví dụ, bảng chân và ảnh cấu trúc ở phụ lục giúp đối chiếu thông tin thực tế của chính từng bo, không phải thông số chuẩn cho toàn bộ họ chip.
\ifdefined\iotstandalone
\clearpage
\appendix
\chapter{Phụ lục: 79 trang slide Chương 2 nguyên bản}
Phụ lục đính kèm \textbf{đầy đủ 79 trang} tài liệu do người học cung cấp, nhằm bảo toàn các ảnh chụp phần cứng, sơ đồ chân, bảng so sánh, ví dụ mã và ảnh giao diện mà phần diễn giải bằng chữ không thể thay thế hoàn toàn. Bố cục slide gốc không bị chỉnh sửa. Khi ghép các chương, có thể bỏ qua phụ lục bằng cách không sao chép đoạn \texttt{\textbackslash includepdf} bên dưới.
\IfFileExists{assets/chapter2_original_slides.pdf}{%
  \includepdf[pages=1-79,fitpaper=false,offset=0 0,pagecommand={},width=.95\paperwidth]{assets/chapter2_original_slides.pdf}%
}{\begin{noteiot}Không có tệp \texttt{assets/chapter2\_original\_slides.pdf} trong thư mục dự án. Module \texttt{chapter2.tex} vẫn chứa các phần diễn giải và bảng tra cứu; để hiện ảnh slide nguyên bản hãy đặt PDF nguồn tại vị trí này rồi biên dịch lại.\end{noteiot}}
\fi % Phụ lục gốc chỉ xuất hiện khi biên dịch riêng Chương 2.

% Nguồn: Chapter 2_IoT Hardware(2).pdf (79 trang)
% Không chứa nội dung Chương 1 hay Chương 3.
% Sử dụng LuaLaTeX; main.tex khai báo fontspec, babel, pdfpages,
% graphicx, tikz, tcolorbox, listings và các macro cho hình.
% =============================================================
\chapter{Phần cứng IoT (IoT Hardware)}\label{chap:iot-hardware}
\begin{concept}
\textbf{Phạm vi Chương 2:} (1) Cảm biến và phân loại cảm biến; (2) cơ cấu chấp hành; (3) nút cảm biến và sơ đồ phần cứng; (4) vi điều khiển ATmega/Arduino cùng kiến thức lập trình; (5) ESP8266 và ESP32; (6) Intel Edison; (7) điện toán biên và Raspberry Pi; (8) so sánh nền tảng. Nội dung được diễn giải sang tiếng Việt và đánh dấu số trang slide gốc nhằm dễ kiểm chứng. Phụ lục cuối chương giữ nguyên \textbf{toàn bộ 79 trang slide}, bao gồm ảnh, sơ đồ chân, ảnh màn hình IDE và bảng biểu.
\end{concept}
\section{Bối cảnh: kiến trúc phần cứng của hệ thống IoT}
Hình ở trang 3 trình bày chuỗi \emph{Things / IoT Devices} (các thiết bị vật lý), \emph{Edge} (điện toán biên), \emph{Network Connectivity} (liên kết mạng) và \emph{Cloud} (đám mây). Về phía phần cứng, thiết bị cảm nhận hoặc tác động lên môi trường; một bo mạch vi điều khiển hay máy tính nhúng thực hiện xử lý; phần giao tiếp mạng trao đổi dữ liệu với biên và đám mây. Các minh họa trong slide đưa ra bo mạch và gateway thuộc nhiều loại khác nhau; không nhất thiết mọi hệ thống phải có tất cả các tầng riêng biệt.
\slidefig{3}{Kiến trúc thiết bị IoT, biên và đám mây}
\section{Cảm biến (Sensors)}
\subsection{Khái niệm và vai trò}
Cảm biến là phần tử cơ bản, thiết yếu của hệ thống IoT, tạo cầu nối giữa thế giới thực và hệ điện tử. Nó đo một đại lượng vật lý, hóa học hoặc sinh học rồi biến đổi thành một tín hiệu mà con người hay thiết bị điện tử có thể đọc. Tín hiệu -- thường là tín hiệu điện -- sau đó được phân tích, xử lý, lưu trữ và truyền đi bằng phương tiện có dây hoặc không dây.

Nếu gọi \(m\) là đại lượng đo và \(s\) là tín hiệu điện đầu ra, quan hệ chuyển đổi được viết:
\begin{equation}\boxed{s=f(m).}\end{equation}
Vì \(f\) mô tả đặc tính chuyển đổi của cảm biến, một phép đo cần quan tâm cả đại lượng đầu vào lẫn cách ánh xạ sang tín hiệu điện; slide không đưa ra hàm hiệu chuẩn cụ thể cho từng cảm biến.\slidesource{4}
\subsection{Các loại cảm biến minh họa}
Hai trang ảnh giới thiệu những ví dụ sau; các tên được chuẩn hóa tiếng Việt và giữ thuật ngữ gốc để thuận tiện tìm linh kiện.
\begin{longtable}{p{.31\linewidth}p{.60\linewidth}}
\toprule\textbf{Tên trong slide} & \textbf{Đại lượng / chức năng tương ứng}\\\midrule\endhead
Temperature sensor & Cảm biến nhiệt độ.\\
Soil moisture sensor & Cảm biến độ ẩm đất.\\
Humidity sensor & Cảm biến độ ẩm môi trường.\\
Infrared sensor; IR optical sensor & Cảm biến hồng ngoại; cảm biến quang hồng ngoại.\\
Gas sensor; alcohol sensor & Cảm biến khí; cảm biến hơi cồn.\\
CO\(_2\) sensor & Cảm biến nồng độ khí carbon dioxide.\\
Digital push button & Nút nhấn số, cung cấp trạng thái đóng/ngắt để làm tín hiệu vào.\\
Ultrasonic sensor & Cảm biến siêu âm.\\
LDR sensor & Quang trở, thay đổi điện trở theo ánh sáng.\\
Gyroscope sensor & Cảm biến tốc độ góc / con quay hồi chuyển.\\
Rain sensor & Cảm biến mưa.\\
Photodiode & Điốt quang.\\
IR proximity sensor; proximity sensor & Cảm biến tiệm cận hồng ngoại; cảm biến tiệm cận nói chung.\\
PIR sensor & Cảm biến hồng ngoại thụ động phát hiện biến thiên bức xạ hồng ngoại.\\\bottomrule
\end{longtable}
\slidesource{5--6}
\slidefig{5}{Ví dụ cảm biến nhiệt độ, độ ẩm, khí, hồng ngoại và nút nhấn}
\slidefig{6}{Bộ sưu tập cảm biến cồn, siêu âm, quang, mưa, tiệm cận, PIR và con quay hồi chuyển}
\subsection{Phân loại theo cách cảm nhận: chủ động và thụ động}
\textbf{Cảm biến tích cực / chủ động (active sensor)} theo cách dùng của bài giảng là loại cần nguồn năng lượng hoặc tín hiệu kích thích từ bên ngoài để thực hiện phép đo. Radar phát sóng vô tuyến, sonar/siêu âm phát sóng âm đến vật thể và nhận phản xạ để xác định khoảng cách. Bài giảng còn liệt kê camera video như ví dụ cảm biến tích cực.\slidesource{7}

\textbf{Cảm biến thụ động (passive sensor)} theo cách trình bày của bài giảng không cần cấp năng lượng kích thích cho chính cơ chế chuyển đổi. \emph{Cặp nhiệt điện} sinh điện áp nhiệt điện khi có chênh lệch nhiệt độ; \emph{cảm biến áp suất} có thể sử dụng màng hoặc vật liệu áp điện tạo ra dịch chuyển/điện áp do tác dụng của áp suất.\slidesource{8}
\begin{noteiot}
Khái niệm chủ động/thụ động ở đây dựa theo cách phân loại trong slide. Khi thiết kế mạch thực tế phải phân biệt \emph{năng lượng kích thích phần tử đo} và \emph{nguồn nuôi toàn bộ mô-đun}; mô-đun có vi mạch xử lý có thể cần cấp nguồn dù phần tử nhạy có cơ chế thụ động.
\end{noteiot}
\subsection{Phân loại theo đại lượng và cơ chế biến đổi}
Ngoài chủ động/thụ động, bài giảng nêu cách phân nhóm theo đối tượng và phương tiện phát hiện (điện, sinh học, hóa học, phóng xạ, v.v.) hoặc hiện tượng chuyển đổi (quang điện, nhiệt điện, sinh học, hóa học, phóng xạ). Các ví dụ theo đại lượng gồm vị trí, nhiệt độ, áp suất, tốc độ, độ ẩm và màu sắc. Đây là các trục phân loại khác nhau, có thể đồng thời mô tả một cảm biến.\slidesource{7--9}
\subsection{Cảm biến tương tự và cảm biến số}
\textbf{Cảm biến tương tự (analog sensor)} tạo tín hiệu có thể biến thiên liên tục trong miền giá trị hoạt động; ví dụ trong slide gồm gia tốc kế, cảm biến ánh sáng, âm thanh, áp suất, nhiệt độ. \textbf{Cảm biến số (digital sensor)} cung cấp giá trị rời rạc hoặc dữ liệu đã số hóa. Bài giảng mô tả cảm biến số được kết hợp với \emph{bộ chuyển đổi tương tự--số} (ADC), còn cảm biến tương tự không tích hợp ADC ở ngõ ra; ví dụ cảm biến gia tốc số, nhiệt độ số DS1620 và HYT~271, cảm biến độ ẩm số. Slide lưu ý nhóm cảm biến số có thể bị giới hạn dải tính toán/đo, đồng thời đưa ví dụ về cách diễn giải đầu ra nhiệt độ.

Phân biệt thêm \emph{tín hiệu số hai mức} (logic HIGH/LOW, ví dụ nút nhấn) với \emph{dữ liệu số nhiều bit} (ví dụ nhiệt độ đọc qua giao tiếp số). Trên slide 10, hình tín hiệu số minh họa các xung nhị phân mức 0/1, còn hình tương tự biểu diễn sóng biến thiên liên tục.\slidesource{9--12}
\slidefig{10}{So sánh tín hiệu tương tự liên tục và tín hiệu nhị phân}
\section{Cơ cấu chấp hành (Actuators)}
\subsection{Định nghĩa và ví dụ}
Cơ cấu chấp hành nhận tín hiệu điều khiển -- chẳng hạn tín hiệu điện từ một bộ điều khiển -- và biến đổi thành tác động đối với vật lý bên ngoài: chuyển động cơ học, biến thiên nhiệt độ, ánh sáng hoặc các dạng năng lượng khác. Bài giảng đưa ví dụ \emph{động cơ bước}: động cơ quay từng góc xác định tương ứng chuỗi xung điều khiển; động cơ bước thích hợp khi yêu cầu điều khiển chính xác vị trí, như cánh tay rô-bốt. Các hình còn cho thấy động cơ DC, màn hình LCD, LED, cơ cấu tuyến tính, ma trận LED và còi/buzzer, qua đó thể hiện cơ cấu chấp hành không chỉ giới hạn ở động cơ quay.\slidesource{13--14}
\slidefig{14}{Ví dụ các cơ cấu chấp hành điện--cơ, hiển thị và âm thanh}
\subsection{So sánh cảm biến và cơ cấu chấp hành}
\begin{center}\begin{tabularx}{\linewidth}{>{\bfseries}p{.22\linewidth}XX}\toprule & Cảm biến & Cơ cấu chấp hành\\\midrule
Chức năng & Chuyển đại lượng vật lý thành tín hiệu đo/điện & Chuyển lệnh hoặc tín hiệu điện thành tác động vật lý\\
Vị trí trong điều khiển & Đầu vào của bộ điều khiển & Đầu ra của bộ điều khiển\\
Ví dụ & Công tắc, micro, cảm biến quay & LED, loa/buzzer, bộ điều khiển động cơ, laser\\\bottomrule\end{tabularx}\end{center}
Một vòng lặp điều khiển điển hình: \emph{môi trường} \(\rightarrow\) \emph{cảm biến} \(\rightarrow\) \emph{bộ điều khiển} \(\rightarrow\) \emph{cơ cấu chấp hành} \(\rightarrow\) \emph{môi trường}.\slidesource{15}
\section{Nút cảm biến IoT (IoT Sensor Node)}
Trong trang 16, bài giảng minh họa cách ghép nhóm \emph{Sensors + Boards + Connectivity} thành \emph{IoT Sensor Nodes}; có ví dụ cảm biến, các bo xử lý, mô-đun truyền thông và bo lắp ráp cuối cùng. Trang 17 đưa sơ đồ khối chi tiết: cảm biến thông minh, cảm biến số, cảm biến video/môi trường; chuyển đổi ADC; MCU; các bộ nhớ flash/RAM; cấp nguồn, DC/DC, quản lý pin; mô-đun truyền thông không dây (Wi-Fi, Bluetooth, IEEE 802.15.4) hoặc có dây. Sơ đồ chỉ ra MCU đóng vai trò phối hợp lấy mẫu, giao tiếp, xử lý và truyền dữ liệu.
\begin{center}\begin{tikzpicture}[node distance=8mm and 7mm, box/.style={draw=navy,rounded corners,fill=blue!4,align=center,minimum height=11mm,minimum width=21mm,font=\footnotesize}, >={Latex}]
\node[box] (s) {Cảm biến};\node[box,right=of s] (adc) {Điều hòa\ tín hiệu / ADC};\node[box,right=of adc] (m) {MCU / CPU};\node[box,right=of m] (n) {Kết nối\ mạng};\node[box,below=of m] (mem) {RAM / Flash};\node[box,below=of adc] (p) {Nguồn, pin,\ DC/DC};
\draw[->] (s)--(adc);\draw[->](adc)--(m);\draw[<->](m)--(n);\draw[<->](m)--(mem);\draw[->](p)--(adc);\draw[->](p.east) -| (m.south);
\end{tikzpicture}\end{center}
\slidefig{17}{Sơ đồ khối nguyên bản của một nút cảm biến IoT}
\section{Vi điều khiển và ATmega328P}
\subsection{Khái niệm microcontroller}
\emph{Microcontroller} (MCU) là một vi mạch tích hợp (IC), có thể hiểu như máy tính thu nhỏ gồm bộ xử lý, bộ nhớ và khối vào/ra. MCU có thể lập trình và được thiết kế cho các ứng dụng nhúng. Trang 18 minh họa các vi mạch Intel P8051, Hitachi HD68000, Motorola MC68HC000LC8. Cần phân biệt \emph{chip MCU} với \emph{bo phát triển} chứa MCU cùng mạch nguồn, giao tiếp và chân cắm.\slidesource{18}
\subsection{Đặc điểm của ATmega328P}
\begin{longtable}{p{.46\linewidth}p{.45\linewidth}}\toprule\textbf{Thông số được nêu trong slide} & \textbf{Giá trị}\\\midrule\endhead
Họ vi điều khiển & AVR 8 bit, hiệu năng cao, công suất thấp\\
Flash chương trình & 32 KB\\
SRAM & 2 KB\\
EEPROM & 1 KB\\
Đường I/O có thể lập trình & 23\\
Điện áp hoạt động & 1,8--5,5 V\\
Dải nhiệt độ & $-40$ đến $85\,^{\circ}\mathrm C$\\
Dòng chế độ Active & 0,2 mA (điều kiện cụ thể không được slide nêu)\\
Dòng chế độ Power-down & 0,1 $\mu$A (điều kiện cụ thể không được slide nêu)\\
Dòng chế độ Power-save & 0,75 $\mu$A (điều kiện cụ thể không được slide nêu)\\\bottomrule\end{longtable}
Hình trang 19 trình bày sơ đồ chân ATmega328P dạng DIP-28, với các nhóm chân: cấp nguồn VCC/AVCC/GND, RESET, dao động XTAL, chân PORTB/PORTC/PORTD, ADC, UART, SPI, I\(^{2}\)C và các chức năng ghép kênh. Trang 21 cho ví dụ mạch lập trình/chuyển chương trình vào chip, đầu nối và bảng pin của giao thức ISP/SPI.\slidesource{19--21}
\slidefig{19}{Sơ đồ 28 chân ATmega328P và các chức năng ghép kênh}
\section{Nền tảng Arduino}
\subsection{Khái quát và lịch sử}
Arduino là nền tảng điện tử \emph{mã nguồn mở} bao gồm cả phần cứng và phần mềm, hướng tới tính dễ dùng. Theo bài giảng, dự án hình thành năm 2003 tại Interaction Design Institute Ivrea (Ý) cho sinh viên; tên gọi bắt nguồn từ quán bar nơi nhóm sáng lập gặp nhau, quán được đặt theo Arduin of Ivrea. Phần cứng ban đầu do Smart Projects tại Ý sản xuất; một số bo mang thương hiệu Arduino được thiết kế bởi các đơn vị khác như SparkFun và Adafruit.\slidesource{22}

Ưu điểm được liệt kê gồm giá thành thấp, đa nền tảng, môi trường lập trình đơn giản, phần mềm mở và phần cứng mở. Các hình giới thiệu họ bo Arduino có nhiều kiểu dáng/kích thước, không phải một bộ xử lý đơn nhất.\slidesource{23--24,31}
\subsection{Arduino Uno và sơ đồ chức năng}
Arduino Uno sử dụng \emph{ATmega328P}, được Arduino.cc phát triển; có 14 chân I/O số và 6 đầu vào analog. Bo có thể lập trình trong Arduino IDE qua cáp USB Type-B; bootloader được nạp sẵn giúp tải chương trình mà không cần thiết bị nạp phần cứng ngoài.

Mười khu vực được chỉ ra trên ảnh bo Uno: (1) các chân digital; (2) LED trên pin 13; (3) LED báo nguồn; (4) ATmega; (5) chân analog; (6) GND và 5V; (7) đầu nối nguồn; (8) LED TX/RX; (9) cổng USB; (10) nút RESET. Sơ đồ pinout ở trang 27 minh họa quan hệ \emph{chân vật lý ATmega328P -- ký hiệu Port -- chức năng Arduino -- đầu vào ADC và các chức năng ngoại vi}.\slidesource{25--27}
\slidefig{26}{Sơ đồ mười khu vực trên bo Arduino Uno}
\slidefig{27}{Bản đồ chân ATmega328P trong hệ Arduino}
\subsection{Các mẫu bo: Leonardo, Mega, Nano}
\begin{longtable}{>{\raggedright\arraybackslash}p{.28\linewidth}p{.20\linewidth}p{.22\linewidth}p{.20\linewidth}}\toprule\textbf{Thông số} & \textbf{Leonardo} & \textbf{Mega} & \textbf{Nano}\\\midrule\endhead
Vi điều khiển & ATmega32U4 & ATmega2560 & ATmega328\\
Điện áp hoạt động & 5 V & 5 V & 5 V\\
Điện áp vào khuyến nghị & 7--12 V & 7--12 V & 7--12 V\\
Giới hạn điện áp vào & 6--20 V & 6--20 V & Không ghi trong slide\\
I/O số & 23 & 54 & 22 (6 PWM)\\
Kênh analog & 12 & 16 & 8\\
Flash & 32 KB, trong đó 4 KB cho bootloader & 256 KB & 32 KB\\
SRAM & 2,5 KB & 8 KB & 2 KB\\
EEPROM & 1 KB & 4 KB & 1 KB\\
Dòng I/O nêu trong slide & 40 mA & 40 mA & Không ghi\\
Dòng chân 3,3 V nêu trong slide & 50 mA & 50 mA & Không ghi\\
Kích thước & 68,6 $\times$ 53,3 mm & Không ghi & 18 $\times$ 45 mm\\
Khối lượng & 20 g & Không ghi & 7 g\\
Thông số bổ sung & -- & -- & 19 mA (tiêu thụ theo slide)\\\bottomrule\end{longtable}
Các mức dòng trong bảng được tái hiện theo slide để phục vụ đối chiếu, không thay thế cho giới hạn thiết kế điện và khuyến cáo của datasheet.\slidesource{28--30}
\subsection{Các lựa chọn ngoài Arduino và shield mở rộng}
Trang 32--34 trình bày \emph{LaunchPad MSP430, Particle Photon, NodeMCU ESP8266, SparkFun Thing Plus, Adafruit Feather HUZZAH, BeagleBoard/PocketBeagle, Seeeduino Nano, SparkFun RedBoard Artemis, STM32F3 Discovery}. Đây là các nền tảng hoặc bo có thể dùng cho nguyên mẫu IoT, không nhất thiết tương thích về phần cứng hay công cụ lập trình.

\emph{Arduino Shield} là bo cắm chồng lên PCB Arduino để bổ sung chức năng, được thiết kế theo hệ chân và triết lý dễ lắp, dễ tháo, dễ sản xuất chi phí thấp. Hình minh họa cho thấy các shield hiển thị, kết nối, nguyên mẫu và mở rộng giao tiếp. Bộ \emph{Arduino Starter Kit} trong trang 36 tập hợp bo mạch, breadboard, linh kiện rời, dây cắm và phụ kiện điện tử.\slidesource{32--36}
\slidefig{34}{Các bo phát triển thay thế được bài giảng liệt kê}
\slidefig{35}{Các dạng Arduino Shield mở rộng}
\subsection{Breadboard và nguyên mẫu}
\emph{Breadboard} là bảng cắm mạch thử nghiệm không cần hàn, thích hợp kiểm chứng nhanh mạch điện và thiết kế trước khi chế tạo PCB. Trang 37 minh họa Arduino nối một LED và điện trở trên breadboard; cần kiểm tra hướng kết nối, hàng/cột chân nối thông và điện trở giới hạn dòng.\slidesource{37}
\slidefig{37}{Mô hình thử nghiệm Arduino với breadboard}
\section{Môi trường và ngôn ngữ lập trình Arduino}
\subsection{Arduino Sketch và Arduino Desktop IDE}
\emph{Sketch} là tên Arduino dùng cho một chương trình, tức đơn vị mã nguồn được biên dịch, nạp và chạy trên bo. Ngôn ngữ Arduino dựa trên các hàm C/C++. Trang 38 mô tả một sketch mẫu và cách tổ chức chương trình thành \texttt{setup()} chạy khởi tạo và \texttt{loop()} thực thi lặp.

Arduino Desktop IDE hỗ trợ Windows, macOS, Linux. Trang 39 dẫn các hướng dẫn cài đặt của Arduino cho từng hệ điều hành, kèm lưu ý cách cài bằng \texttt{sudo apt-get install arduino} trên Linux không được bài giảng khuyến nghị. Trang 40 giới thiệu giao diện IDE, nơi viết sketch, xác minh/biên dịch, chọn bo/cổng giao tiếp, nạp và theo dõi cửa sổ Serial Monitor.\slidesource{38--40}
\slidefig{40}{Giao diện Arduino IDE trong slide gốc}
\subsection{Chú thích chương trình (comments)}
Chú thích được trình biên dịch bỏ qua, dùng để giải thích ý nghĩa mã lệnh. Trang 41 minh họa hai dạng:
\begin{lstlisting}[language=C++]
/* Day la chu thich nhieu dong.
   Cac dong nay duoc bo qua. */
// Day la chu thich mot dong.
\end{lstlisting}
\noindent Lưu ý: chú thích nhiều dòng theo cú pháp \texttt{/* ... */} không được viết lồng nhau; dấu \texttt{//} dùng cho chú thích một dòng.
\subsection{Cấu trúc điều khiển: for, if, while, do--while}
\paragraph{Vòng lặp for (slide 42).} Lặp khối câu lệnh giữa ngoặc nhọn với bộ đếm thường gồm ba biểu thức khởi tạo, điều kiện, bước cập nhật. Biến đếm có thể tăng hoặc giảm.
\begin{lstlisting}[language=C++]
for (int i = 0; i <= 10; i++) {
  Serial.write(7);
}
\end{lstlisting}
\paragraph{Rẽ nhánh if (slide 43).} Kiểm tra biểu thức điều kiện; nếu đúng thì thực hiện nhánh \texttt{if}, nếu sai có thể thực hiện \texttt{else}. Bài giảng minh họa kiểm tra nhiệt độ so với ngưỡng \(70\), đồng thời liệt kê các toán tử \texttt{==}, \texttt{!=}, \texttt{<}, \texttt{<=}, \texttt{>}, \texttt{>=}. Cần phân biệt \texttt{=} (gán) và \texttt{==} (so sánh bằng).
\begin{lstlisting}[language=C++]
if (temperature > 70) {
  // Xu ly khi cao hon nguong
} else {
  // Xu ly khi khong cao hon nguong
}
\end{lstlisting}
\paragraph{Vòng lặp while (slide 44).} Lặp khi điều kiện vẫn đúng; kiểm tra điều kiện trước vòng lặp. \texttt{do\{...\}while(...);} thực hiện thân lặp ít nhất một lần rồi mới kiểm tra điều kiện.
\begin{lstlisting}[language=C++]
int var = 0;
while (var < 200) { var++; }

do { var++; } while (var < 400);
\end{lstlisting}
\slidesource{41--44}
\subsection{Hàm đọc và ghi tín hiệu số}
\texttt{pinMode(pin, mode)} thiết lập chân làm đầu vào hoặc đầu ra (trong slide: \texttt{INPUT}, \texttt{OUTPUT}). \texttt{digitalRead(pin)} đọc trạng thái logic \texttt{HIGH}/\texttt{LOW}. \texttt{digitalWrite(pin, value)} xuất logic \texttt{HIGH} hoặc \texttt{LOW} ra một chân đã được cấu hình phù hợp. Tham số \texttt{pin} là số hiệu chân digital cần tác động, \texttt{value} là mức logic yêu cầu.\slidesource{45}
\begin{lstlisting}[language=C++]
const int LED_PIN = 13;
void setup() { pinMode(LED_PIN, OUTPUT); }
void loop() {
  digitalWrite(LED_PIN, HIGH);
  delay(1000);
  digitalWrite(LED_PIN, LOW);
  delay(1000);
}
\end{lstlisting}
\subsection{Tín hiệu tương tự và PWM}
\texttt{analogRead(pin)} đọc đầu vào ADC (slide mô tả bộ chuyển đổi 10 bit). \texttt{analogWrite(pin, value)} xuất \emph{điều chế độ rộng xung} PWM lên chân hỗ trợ; giá trị duty trong ví dụ Arduino nằm trong khoảng 0--255, tương ứng thời gian xung ở mức cao trong mỗi chu kỳ, chứ không phải một đầu ra điện áp DAC lý tưởng. Hình trang 46 biểu diễn dạng sóng PWM có độ rộng khác nhau.\slidesource{46}
\slidefig{46}{Sơ đồ minh họa nguyên lý PWM và các giá trị analog của Arduino}
\subsection{Các hàm toán học thông dụng}
\begin{longtable}{p{.26\linewidth}p{.64\linewidth}}\toprule\textbf{Hàm} & \textbf{Tác dụng theo bài giảng}\\\midrule\endhead
\texttt{abs(x)} & Trả về trị tuyệt đối.\\
\texttt{constrain(x,a,b)} & Giới hạn \(x\) trong đoạn từ \(a\) đến \(b\).\\
\texttt{max(x,y)} & Chọn giá trị lớn nhất.\\
\texttt{min(x,y)} & Chọn giá trị nhỏ nhất.\\
\texttt{pow(base,exponent)} & Lũy thừa \(\mathrm{base}^{\mathrm{exponent}}\).\\
\texttt{sq(x)} & Bình phương \(x\).\\
\texttt{sqrt(x)} & Căn bậc hai \(\sqrt{x}\).\\\bottomrule\end{longtable}
\slidesource{47}
\subsection{Trì hoãn và chỉ thị tiền xử lý}
\texttt{delay(milliseconds)} tạm dừng chương trình trong khoảng thời gian tính theo mili-giây, \(1000\,\mathrm{ms}=1\,\mathrm{s}\). \texttt{\#define} đặt tên hằng hoặc macro trước khi chương trình được biên dịch, ví dụ \texttt{\#define ledPin 3}. \texttt{\#include} thêm tiêu đề thư viện, ví dụ \texttt{\#include <SPI.h>}; nhờ đó sketch truy cập các hàm, lớp và định nghĩa đã được thư viện cung cấp.\slidesource{48--49}
\begin{lstlisting}[language=C++]
#define ledPin 3
#include <SPI.h>
\end{lstlisting}
\subsection{Hai hàm chính setup() và loop()}
\texttt{setup()} được gọi một lần mỗi khi khởi động hoặc reset để thiết lập chân, Serial, thư viện và các trạng thái ban đầu. \texttt{loop()} được gọi lặp liên tục sau khi \texttt{setup()} hoàn tất, là nơi đặt chương trình chính xử lý tín hiệu và điều khiển thiết bị. Slide 50 dùng \texttt{pinMode(8, INPUT)} minh họa thao tác khởi tạo; slide 51 minh họa dùng \texttt{digitalRead(8)} và ghi chữ qua Serial theo trạng thái nút nhấn.\slidesource{50--51}
\begin{lstlisting}[language=C++]
void setup() {
  pinMode(8, INPUT);
  Serial.begin(9600);
}
void loop() {
  if (digitalRead(8) == HIGH) {
    Serial.write('H');
  } else {
    Serial.write('L');
  }
}
\end{lstlisting}
\subsection{Thêm thư viện và giao tiếp Serial}
Trang 52 minh họa thao tác \emph{Sketch > Include Library} trong Arduino IDE. Các thư viện có thể đi kèm IDE hoặc được thêm bởi người dùng. Cổng COM là tên truyền thống của cổng giao tiếp nối tiếp trên máy tính tương thích IBM PC, ngày nay còn bao gồm cổng nối tiếp ảo thông qua Bluetooth hoặc cầu USB-to-UART. Arduino thường cung cấp ít nhất một cổng nối tiếp, một số bo có nhiều cổng. Trên Arduino Uno, \textbf{pin 0 là RX và pin 1 là TX}; người lập trình thường dùng \texttt{Serial.begin(speed)} khởi tạo với tốc độ baud và \texttt{Serial.write(value)} ghi dữ liệu ra cổng; baud là số ký hiệu truyền mỗi giây (bài giảng mô tả gần đúng theo bit/giây).\slidesource{52--53}
\slidefig{52}{Thao tác thêm thư viện trong Arduino IDE}
\section{Nền tảng Wi-Fi ESP8266}
\subsection{Tổng quan ESP8266EX và khả năng phần mềm}
ESP8266EX là chip Wi-Fi chi phí thấp của Espressif Systems, tích hợp giao thức TCP/IP và có thể ghép vào các thiết kế dùng vi điều khiển. Phần lõi hỗ trợ Arduino của ESP8266 cung cấp thư viện Wi-Fi cho TCP, UDP, HTTP, mDNS, SSDP, máy chủ DNS; cho phép cập nhật firmware qua mạng (OTA), hệ thống tệp trên flash hoặc SD, điều khiển servo và trao đổi SPI/I\(^{2}\)C. Slide 55 liệt kê các biến thể mô-đun ESP-01, ESP-02, ESP-03, ESP-04, ESP-05, ESP-06, ESP-07, ESP-08, ESP-09, ESP-10, ESP-11, ESP-12, ESP-13.\slidesource{54--55}
\slidefig{55}{Tập hợp biến thể mô-đun ESP8266}
\subsection{ESP-01 và ESP-07}
\begin{longtable}{p{.34\linewidth}p{.57\linewidth}}\toprule\textbf{Mô-đun} & \textbf{Đặc điểm được cung cấp trong slide}\\\midrule\endhead
ESP-01 & Kích thước $24{,}75\times14{,}5$ mm; nguồn 3,3 V; anten mạch in tích hợp; LED nguồn và TXD; chế độ AP, STA, AP+STA; có thể điều khiển bằng lệnh AT; khoảng cách chân 2,54 mm.\\
ESP-07 & Nguồn 3,3 V; anten gốm tích hợp và đầu nối anten ngoài U.FL; Wi-Fi 802.11b/g/n băng 2,4 GHz với WPA/WPA2; TCP/UDP; hỗ trợ lệnh AT; khoảng cách chân 2 mm.\\\bottomrule\end{longtable}
\slidesource{56--57}
\slidefig{56}{Sơ đồ chân, module ESP-01 và đặc điểm ngoại hình}
\slidefig{57}{Sơ đồ chân ESP-07 trong bài giảng}
\subsection{SDK của Espressif và môi trường thay thế}
Bài giảng nêu năm 2014 Espressif giới thiệu SDK cho phép ESP8266 chạy như vi điều khiển độc lập thay vì chỉ là mô-đun Wi-Fi của MCU khác. Hai nhánh SDK được nhắc đến là \emph{ESP8266 NonOS SDK} và \emph{ESP8266 RTOS SDK}. Các môi trường hoặc SDK thay thế trên slide gồm ESP-Open-SDK, Unofficial Development Kit, Arduino, ESP8266 BASIC, ESP Easy và ESPHome.\slidesource{58}
\subsection{Cấu hình ESP8266 trên Arduino IDE}
Ba trang hướng dẫn tuần tự: \textbf{(1)} vào \emph{File > Preferences} và thêm URL gói bo \url{http://arduino.esp8266.com/stable/package_esp8266com_index.json}; \textbf{(2)} vào \emph{Tools > Board > Boards Manager}, tìm ``esp8266'' và chọn \emph{Install}; \textbf{(3)} vào \emph{Tools > Board} chọn \emph{Generic ESP8266 Module} nếu đúng loại mô-đun cần sử dụng. Các ảnh minh họa trong slide là giao diện IDE tại thời điểm biên soạn; vị trí menu có thể khác ở bản IDE mới.\slidesource{59--61}
\slidefig{59}{Thêm liên kết JSON của gói bo ESP8266}
\slidefig{60}{Cài gói hỗ trợ ESP8266 qua Boards Manager}
\slidefig{61}{Chọn bo Generic ESP8266 Module}
\section{Nền tảng ESP32}
ESP32 là họ \emph{system-on-chip} giá tương đối thấp do Espressif phát triển, là bước tiến so với họ ESP8266. Slide giới thiệu khả năng đa lõi, đồng xử lý tiết kiệm năng lượng và hỗ trợ kết nối Wi-Fi, Bluetooth. Những thiết bị minh họa gồm ESP32 Development Board, ESP32-CAM, ESP32-WROOM-32U/-32D, DOIT DevKit V1, ESP32 DevKit, ESP32S NodeMCU, ESP32 Thing, Wemos Lolin32, Wemos OLED và HUZZAH32.

Các \textbf{ưu điểm theo slide}: CPU dual-core khoảng 160--240 MHz; Wi-Fi/Bluetooth tiện cho điều khiển và giám sát; số chân GPIO nhiều hơn ESP8266; slide ghi tốc độ tối đa 150 Mb/s. \textbf{Nhược điểm theo slide}: giá cao hơn ESP8266. Thông số là mô tả cho các phiên bản ESP32 được trình bày trong bài giảng, không đại diện cho mọi biến thể ESP32. Đặc biệt, ảnh trang 64 minh họa pinout bo DOIT ESP32 DevKit V1 loại 30 chân, gồm các chân cấp nguồn, EN, chân UART, ADC, digital GPIO, SPI và I\(^{2}\)C tùy cấu hình.\slidesource{62--65}
\slidefig{63}{Một số bo ESP32 được giới thiệu trong slide}
\slidefig{64}{Sơ đồ chân bo DOIT ESP32 DevKit V1 (30 GPIOs theo tiêu đề slide)}
\section{Nền tảng Intel Edison}
\subsection{Cấu hình mô-đun Edison}
Intel Edison trong bài giảng có kích thước $35{,}5\times25\times3{,}9$ mm, linh kiện gắn ở hai mặt. SoC Intel Atom Tangier Z34XX tiến trình 22 nm tích hợp hai lõi Silvermont 500 MHz và một lõi Intel Quark ở 100 MHz (đơn vị trong slide không ghi đầy đủ tại dòng này); RAM 1 GB và eMMC 4 GB; có Wi-Fi, Bluetooth 4, điều khiển USB. Kết nối ra ngoài qua đầu nối Hirose DF40 mật độ cao 70 chân (USB, SD, UART, GPIO). Các môi trường lập trình nêu trong slide: Arduino IDE, Eclipse (C, C++, Python), Intel XDK (NodeJS/HTML5).\slidesource{66}
\slidefig{66}{Mô-đun Intel Edison và cấu hình phần cứng}
\subsection{Bo breakout và bo dạng Arduino cho Edison}
\textbf{Breakout board}: xuất các mức I/O gốc 1,8 V của Edison ra các điểm hàn dạng lưới 0,1 inch; có USB OTG Micro Type-AB, công tắc nguồn OTG, bộ sạc pin, cầu USB--UART qua Micro-USB, đầu cấp nguồn DC 7--15 V. Điện áp I/O 1,8 V là thông tin quan trọng khi giao tiếp với ngoại vi.

\textbf{Arduino board for Edison Module}: 20 I/O số (6 PWM), 6 đầu vào analog, 1 UART, 1 I\(^{2}\)C, đầu ICSP 6 chân cho SPI, cổng Micro-USB device hoặc cổng USB host Type-A chọn qua công tắc cơ, cổng Micro-USB UART, khe SD và đầu nguồn DC 7--15 V.\slidesource{67--68}
\slidefig{67}{Intel Edison Breakout Board}
\slidefig{68}{Bo tương thích bố cục Arduino cho Edison}
\section{Điện toán biên (Edge Computing)}
Slide 69 trình bày sự khác nhau về vị trí xử lý giữa đám mây, trung tâm dữ liệu doanh nghiệp và biên gần thiết bị; trang 70 nhấn mạnh \emph{xử lý thông tin nhạy cảm với độ trễ}, \emph{quyền riêng tư và an ninh}, \emph{băng thông và chất lượng kết nối}. Trang 71 mô tả dữ liệu do thiết bị IoT tạo ra được xử lý một phần ở máy hoặc gateway gần nguồn, trong khi đám mây và trung tâm dữ liệu tiếp tục thực hiện các công việc ở quy mô lớn. Khác với gửi tất cả dữ liệu thô lên cloud, cách bố trí này có thể giảm lượng dữ liệu truyền và cho phép đưa ra quyết định cục bộ.

Các \textbf{thiết bị biên} minh họa trang 72 gồm Raspberry Pi, NVIDIA Jetson Nano và NVIDIA Jetson Xavier. Liên kết tham khảo trên slide: \url{https://www.networkworld.com/article/3224893/what-is-edge-computing-and-how-it-s-changing-the-network.html}.\slidesource{69--72}
\slidefig{69}{Quan hệ giữa cloud, edge gateway và trung tâm dữ liệu}
\slidefig{71}{Sơ đồ hoạt động của điện toán biên}
\slidefig{72}{Các máy tính biên: Raspberry Pi và NVIDIA Jetson}
\section{Máy tính đơn bo Raspberry Pi}
\subsection{Khái quát}
Raspberry Pi là dòng \emph{single-board computers} (SBC), tức máy tính tích hợp phần lớn thành phần trên một bo mạch nhỏ gọn. Theo số liệu lịch sử \textbf{được trích trong slide}, đến tháng 2/2015 hơn 5 triệu máy Raspberry Pi đã được bán, trở thành máy tính Anh bán chạy nhất thời điểm đó. Dòng bo do Raspberry Pi Foundation phát triển nhằm khuyến khích học khoa học máy tính cơ bản trong trường học và tại những quốc gia đang phát triển. Slide nêu các bo có Wi-Fi/Bluetooth tích hợp; khả năng này tùy theo phiên bản thực tế. Hình trang 73 minh họa các dự án robot và bộ phụ kiện Raspberry Pi.\slidesource{73--74}
\slidefig{73}{Hệ sinh thái thiết bị và ứng dụng Raspberry Pi}
\subsection{Bộ kit nhập môn}
Trang 75 liệt kê các thành phần: bo Raspberry Pi, thẻ microSD dung lượng \textbf{8 GB trở lên} theo slide, bộ cấp nguồn, bàn phím/chuột USB, màn hình có ngõ HDMI, cáp HDMI và hộp vỏ bảo vệ. Thẻ microSD làm bộ nhớ khởi động và lưu hệ điều hành đối với các cấu hình được minh họa trong bài giảng.\slidesource{75}
\slidefig{75}{Bộ kit Raspberry Pi gồm bo, microSD, nguồn và ngoại vi}
\subsection{Cấu tạo và chân Raspberry Pi 4}
Trang 76 chỉ các phần \textbf{USB-C power}, \textbf{micro-HDMI} (hai đầu ra), \textbf{Ethernet Gigabit}, cổng \textbf{USB 3.0/2.0}, header GPIO 40 chân và lựa chọn RAM 2/4/8 GB trong hình. Trang 77 chứa bảng pinout 40 chân: nguồn 3,3 V, 5 V, GND, các chân GPIO và các chân chức năng SPI/UART/I\(^{2}\)C. Khi đấu nối phải đối chiếu số thứ tự chân vật lý (physical pin) và tên GPIO/BCM, tránh cấp 5 V vào tín hiệu GPIO không chịu 5 V.\slidesource{76--77}
\slidefig{76}{Các cổng kết nối và khu vực chính của Raspberry Pi 4}
\slidefig{77}{Sơ đồ chân header 40 pin Raspberry Pi 4}
\section{So sánh Arduino Uno và Raspberry Pi 4}
\begin{longtable}{p{.38\linewidth}p{.26\linewidth}p{.26\linewidth}}\toprule\textbf{Tiêu chí trong bài giảng} & \textbf{Arduino Uno} & \textbf{Raspberry Pi 4}\\\midrule\endhead
Loại nền tảng & Vi điều khiển (microcontroller) & Máy vi tính (microcomputer)\\
Hệ điều hành & Không chạy hệ điều hành đa nhiệm thông dụng & Có thể chạy hệ điều hành\\
Đa nhiệm & Slide ghi ``No multitasking'' & Có đa nhiệm\\
Giá thành tương đối & Rẻ & Đắt hơn\\
Tài nguyên & Xung 16 MHz, SRAM 2 KB & Xung khoảng 1,5 GHz, RAM 2/4/8 GB\\
Ứng dụng thời gian thực & Được slide đánh giá là thích hợp & Slide đánh giá không lý tưởng cho thời gian thực cứng\\
Kết nối mạng tích hợp & Slide ghi không có kết nối mạng & Có khả năng kết nối mạng\\
ADC tích hợp trên bo & Có & Slide ghi không có bộ chuyển đổi A/D\\\bottomrule\end{longtable}
Dòng so sánh phản ánh \textbf{đúng hai bo được nêu trong slide 78}, không phải khẳng định cho toàn bộ họ Arduino và Raspberry Pi. Một số chức năng có thể mở rộng bằng shield, HAT hoặc IC ngoài; với Raspberry Pi, yêu cầu đo analog thường cần thêm ADC bên ngoài.\slidesource{78}
\slidefig{78}{Bảng so sánh Arduino Uno và Raspberry Pi 4 nguyên bản}
\section{Bảng tra cứu tương ứng với toàn bộ 79 slide}
\begin{longtable}{p{.15\linewidth}p{.76\linewidth}}\toprule\textbf{Trang PDF} & \textbf{Chủ đề và hình minh họa được lưu trong phụ lục gốc}\\\midrule\endhead
1--3 & Bìa giảng viên, mục lục chương và kiến trúc Cloud--Edge--IoT Devices.\\
4--6 & Khái niệm cảm biến, hàm chuyển đổi, ảnh mô-đun cảm biến đa dạng.\\
7--12 & Chủ động/thụ động; phân loại theo hiện tượng và analog/digital; dạng sóng; cảm biến tương tự và cảm biến số.\\
13--15 & Cơ cấu chấp hành, hình các thiết bị đầu ra, so sánh cảm biến/cơ cấu chấp hành.\\
16--17 & Phân rã phần cứng nút IoT và sơ đồ khối IoT Sensor Node.\\
18--21 & Khái niệm MCU, chip ATmega328P, cấu hình và bộ nạp.\\
22--27 & Lịch sử/ưu điểm Arduino, bo Uno, các vùng chức năng và pinout ATmega328P.\\
28--31 & Arduino Leonardo, Mega, Nano và ảnh các biến thể Arduino.\\
32--35 & Bo thay thế Arduino và shield mở rộng.\\
36--40 & Starter Kit, breadboard, Arduino Sketch, cài đặt và giao diện Arduino IDE.\\
41--44 & Comment và cấu trúc \texttt{for}, \texttt{if}, \texttt{while}, \texttt{do--while}.\\
45--48 & Digital I/O, Analog I/O và PWM, hàm toán học và delay.\\
49--53 & \texttt{\#define}, \texttt{\#include}, \texttt{setup()}, \texttt{loop()}, thư viện và Serial.\\
54--58 & ESP8266, các biến thể mô-đun, đặc điểm ESP-01, ESP-07, SDK.\\
59--61 & Quy trình cài đặt board ESP8266 trong Arduino IDE với ảnh các bước.\\
62--65 & ESP32, các bo ESP32, sơ đồ chân DOIT và ưu/nhược điểm.\\
66--68 & Intel Edison, bo breakout và bo giao tiếp phong cách Arduino.\\
69--72 & Điện toán biên, mô hình hoạt động và các bo edge devices.\\
73--77 & Raspberry Pi, bộ kit, sơ đồ các cổng và pinout Pi 4.\\
78--79 & So sánh Uno và Pi 4, lời kết.\\\bottomrule\end{longtable}
\section{Kết luận Chương 2}
Phần cứng IoT được tổ chức xung quanh \emph{cảm biến} (thu nhận dữ liệu), \emph{cơ cấu chấp hành} (tạo tác động), \emph{bộ điều khiển/tính toán} (quản lý, xử lý, điều khiển), \emph{giao diện I/O}, \emph{khối nguồn} và \emph{khối truyền thông}. Với các tác vụ điều khiển nhúng, Arduino/ATmega và ESP8266/ESP32 là những ví dụ tiêu biểu của MCU. Với các tác vụ cần hệ điều hành, đa nhiệm, xử lý ở biên và giao tiếp phần mềm đa dạng, Raspberry Pi hoặc một bo công nghiệp có thể phù hợp hơn. Intel Edison được bài giảng dùng làm minh họa một kiến trúc mô-đun máy tính nhúng kết hợp nhiều giao diện ngoại vi.

\textbf{Điểm cần ghi nhớ:} lựa chọn không chỉ theo xung CPU, mà theo số chân, điện áp I/O, ADC, dung lượng RAM/Flash, giao tiếp không dây, nguồn, phần mềm công cụ và độ trễ xử lý mong muốn. Các ví dụ, bảng chân và ảnh cấu trúc ở phụ lục giúp đối chiếu thông tin thực tế của chính từng bo, không phải thông số chuẩn cho toàn bộ họ chip.
\ifdefined\iotstandalone
\clearpage
\appendix
\chapter{Phụ lục: 79 trang slide Chương 2 nguyên bản}
Phụ lục đính kèm \textbf{đầy đủ 79 trang} tài liệu do người học cung cấp, nhằm bảo toàn các ảnh chụp phần cứng, sơ đồ chân, bảng so sánh, ví dụ mã và ảnh giao diện mà phần diễn giải bằng chữ không thể thay thế hoàn toàn. Bố cục slide gốc không bị chỉnh sửa. Khi ghép các chương, có thể bỏ qua phụ lục bằng cách không sao chép đoạn \texttt{\textbackslash includepdf} bên dưới.
\IfFileExists{assets/chapter2_original_slides.pdf}{%
  \includepdf[pages=1-79,fitpaper=false,offset=0 0,pagecommand={},width=.95\paperwidth]{assets/chapter2_original_slides.pdf}%
}{\begin{noteiot}Không có tệp \texttt{assets/chapter2\_original\_slides.pdf} trong thư mục dự án. Module \texttt{chapter2.tex} vẫn chứa các phần diễn giải và bảng tra cứu; để hiện ảnh slide nguyên bản hãy đặt PDF nguồn tại vị trí này rồi biên dịch lại.\end{noteiot}}
\fi % Phụ lục gốc chỉ xuất hiện khi biên dịch riêng Chương 2.

% Nguồn: Chapter 2_IoT Hardware(2).pdf (79 trang)
% Không chứa nội dung Chương 1 hay Chương 3.
% Sử dụng LuaLaTeX; main.tex khai báo fontspec, babel, pdfpages,
% graphicx, tikz, tcolorbox, listings và các macro cho hình.
% =============================================================
\chapter{Phần cứng IoT (IoT Hardware)}\label{chap:iot-hardware}
\begin{concept}
\textbf{Phạm vi Chương 2:} (1) Cảm biến và phân loại cảm biến; (2) cơ cấu chấp hành; (3) nút cảm biến và sơ đồ phần cứng; (4) vi điều khiển ATmega/Arduino cùng kiến thức lập trình; (5) ESP8266 và ESP32; (6) Intel Edison; (7) điện toán biên và Raspberry Pi; (8) so sánh nền tảng. Nội dung được diễn giải sang tiếng Việt và đánh dấu số trang slide gốc nhằm dễ kiểm chứng. Phụ lục cuối chương giữ nguyên \textbf{toàn bộ 79 trang slide}, bao gồm ảnh, sơ đồ chân, ảnh màn hình IDE và bảng biểu.
\end{concept}
\section{Bối cảnh: kiến trúc phần cứng của hệ thống IoT}
Hình ở trang 3 trình bày chuỗi \emph{Things / IoT Devices} (các thiết bị vật lý), \emph{Edge} (điện toán biên), \emph{Network Connectivity} (liên kết mạng) và \emph{Cloud} (đám mây). Về phía phần cứng, thiết bị cảm nhận hoặc tác động lên môi trường; một bo mạch vi điều khiển hay máy tính nhúng thực hiện xử lý; phần giao tiếp mạng trao đổi dữ liệu với biên và đám mây. Các minh họa trong slide đưa ra bo mạch và gateway thuộc nhiều loại khác nhau; không nhất thiết mọi hệ thống phải có tất cả các tầng riêng biệt.
\slidefig{3}{Kiến trúc thiết bị IoT, biên và đám mây}
\section{Cảm biến (Sensors)}
\subsection{Khái niệm và vai trò}
Cảm biến là phần tử cơ bản, thiết yếu của hệ thống IoT, tạo cầu nối giữa thế giới thực và hệ điện tử. Nó đo một đại lượng vật lý, hóa học hoặc sinh học rồi biến đổi thành một tín hiệu mà con người hay thiết bị điện tử có thể đọc. Tín hiệu -- thường là tín hiệu điện -- sau đó được phân tích, xử lý, lưu trữ và truyền đi bằng phương tiện có dây hoặc không dây.

Nếu gọi \(m\) là đại lượng đo và \(s\) là tín hiệu điện đầu ra, quan hệ chuyển đổi được viết:
\begin{equation}\boxed{s=f(m).}\end{equation}
Vì \(f\) mô tả đặc tính chuyển đổi của cảm biến, một phép đo cần quan tâm cả đại lượng đầu vào lẫn cách ánh xạ sang tín hiệu điện; slide không đưa ra hàm hiệu chuẩn cụ thể cho từng cảm biến.\slidesource{4}
\subsection{Các loại cảm biến minh họa}
Hai trang ảnh giới thiệu những ví dụ sau; các tên được chuẩn hóa tiếng Việt và giữ thuật ngữ gốc để thuận tiện tìm linh kiện.
\begin{longtable}{p{.31\linewidth}p{.60\linewidth}}
\toprule\textbf{Tên trong slide} & \textbf{Đại lượng / chức năng tương ứng}\\\midrule\endhead
Temperature sensor & Cảm biến nhiệt độ.\\
Soil moisture sensor & Cảm biến độ ẩm đất.\\
Humidity sensor & Cảm biến độ ẩm môi trường.\\
Infrared sensor; IR optical sensor & Cảm biến hồng ngoại; cảm biến quang hồng ngoại.\\
Gas sensor; alcohol sensor & Cảm biến khí; cảm biến hơi cồn.\\
CO\(_2\) sensor & Cảm biến nồng độ khí carbon dioxide.\\
Digital push button & Nút nhấn số, cung cấp trạng thái đóng/ngắt để làm tín hiệu vào.\\
Ultrasonic sensor & Cảm biến siêu âm.\\
LDR sensor & Quang trở, thay đổi điện trở theo ánh sáng.\\
Gyroscope sensor & Cảm biến tốc độ góc / con quay hồi chuyển.\\
Rain sensor & Cảm biến mưa.\\
Photodiode & Điốt quang.\\
IR proximity sensor; proximity sensor & Cảm biến tiệm cận hồng ngoại; cảm biến tiệm cận nói chung.\\
PIR sensor & Cảm biến hồng ngoại thụ động phát hiện biến thiên bức xạ hồng ngoại.\\\bottomrule
\end{longtable}
\slidesource{5--6}
\slidefig{5}{Ví dụ cảm biến nhiệt độ, độ ẩm, khí, hồng ngoại và nút nhấn}
\slidefig{6}{Bộ sưu tập cảm biến cồn, siêu âm, quang, mưa, tiệm cận, PIR và con quay hồi chuyển}
\subsection{Phân loại theo cách cảm nhận: chủ động và thụ động}
\textbf{Cảm biến tích cực / chủ động (active sensor)} theo cách dùng của bài giảng là loại cần nguồn năng lượng hoặc tín hiệu kích thích từ bên ngoài để thực hiện phép đo. Radar phát sóng vô tuyến, sonar/siêu âm phát sóng âm đến vật thể và nhận phản xạ để xác định khoảng cách. Bài giảng còn liệt kê camera video như ví dụ cảm biến tích cực.\slidesource{7}

\textbf{Cảm biến thụ động (passive sensor)} theo cách trình bày của bài giảng không cần cấp năng lượng kích thích cho chính cơ chế chuyển đổi. \emph{Cặp nhiệt điện} sinh điện áp nhiệt điện khi có chênh lệch nhiệt độ; \emph{cảm biến áp suất} có thể sử dụng màng hoặc vật liệu áp điện tạo ra dịch chuyển/điện áp do tác dụng của áp suất.\slidesource{8}
\begin{noteiot}
Khái niệm chủ động/thụ động ở đây dựa theo cách phân loại trong slide. Khi thiết kế mạch thực tế phải phân biệt \emph{năng lượng kích thích phần tử đo} và \emph{nguồn nuôi toàn bộ mô-đun}; mô-đun có vi mạch xử lý có thể cần cấp nguồn dù phần tử nhạy có cơ chế thụ động.
\end{noteiot}
\subsection{Phân loại theo đại lượng và cơ chế biến đổi}
Ngoài chủ động/thụ động, bài giảng nêu cách phân nhóm theo đối tượng và phương tiện phát hiện (điện, sinh học, hóa học, phóng xạ, v.v.) hoặc hiện tượng chuyển đổi (quang điện, nhiệt điện, sinh học, hóa học, phóng xạ). Các ví dụ theo đại lượng gồm vị trí, nhiệt độ, áp suất, tốc độ, độ ẩm và màu sắc. Đây là các trục phân loại khác nhau, có thể đồng thời mô tả một cảm biến.\slidesource{7--9}
\subsection{Cảm biến tương tự và cảm biến số}
\textbf{Cảm biến tương tự (analog sensor)} tạo tín hiệu có thể biến thiên liên tục trong miền giá trị hoạt động; ví dụ trong slide gồm gia tốc kế, cảm biến ánh sáng, âm thanh, áp suất, nhiệt độ. \textbf{Cảm biến số (digital sensor)} cung cấp giá trị rời rạc hoặc dữ liệu đã số hóa. Bài giảng mô tả cảm biến số được kết hợp với \emph{bộ chuyển đổi tương tự--số} (ADC), còn cảm biến tương tự không tích hợp ADC ở ngõ ra; ví dụ cảm biến gia tốc số, nhiệt độ số DS1620 và HYT~271, cảm biến độ ẩm số. Slide lưu ý nhóm cảm biến số có thể bị giới hạn dải tính toán/đo, đồng thời đưa ví dụ về cách diễn giải đầu ra nhiệt độ.

Phân biệt thêm \emph{tín hiệu số hai mức} (logic HIGH/LOW, ví dụ nút nhấn) với \emph{dữ liệu số nhiều bit} (ví dụ nhiệt độ đọc qua giao tiếp số). Trên slide 10, hình tín hiệu số minh họa các xung nhị phân mức 0/1, còn hình tương tự biểu diễn sóng biến thiên liên tục.\slidesource{9--12}
\slidefig{10}{So sánh tín hiệu tương tự liên tục và tín hiệu nhị phân}
\section{Cơ cấu chấp hành (Actuators)}
\subsection{Định nghĩa và ví dụ}
Cơ cấu chấp hành nhận tín hiệu điều khiển -- chẳng hạn tín hiệu điện từ một bộ điều khiển -- và biến đổi thành tác động đối với vật lý bên ngoài: chuyển động cơ học, biến thiên nhiệt độ, ánh sáng hoặc các dạng năng lượng khác. Bài giảng đưa ví dụ \emph{động cơ bước}: động cơ quay từng góc xác định tương ứng chuỗi xung điều khiển; động cơ bước thích hợp khi yêu cầu điều khiển chính xác vị trí, như cánh tay rô-bốt. Các hình còn cho thấy động cơ DC, màn hình LCD, LED, cơ cấu tuyến tính, ma trận LED và còi/buzzer, qua đó thể hiện cơ cấu chấp hành không chỉ giới hạn ở động cơ quay.\slidesource{13--14}
\slidefig{14}{Ví dụ các cơ cấu chấp hành điện--cơ, hiển thị và âm thanh}
\subsection{So sánh cảm biến và cơ cấu chấp hành}
\begin{center}\begin{tabularx}{\linewidth}{>{\bfseries}p{.22\linewidth}XX}\toprule & Cảm biến & Cơ cấu chấp hành\\\midrule
Chức năng & Chuyển đại lượng vật lý thành tín hiệu đo/điện & Chuyển lệnh hoặc tín hiệu điện thành tác động vật lý\\
Vị trí trong điều khiển & Đầu vào của bộ điều khiển & Đầu ra của bộ điều khiển\\
Ví dụ & Công tắc, micro, cảm biến quay & LED, loa/buzzer, bộ điều khiển động cơ, laser\\\bottomrule\end{tabularx}\end{center}
Một vòng lặp điều khiển điển hình: \emph{môi trường} \(\rightarrow\) \emph{cảm biến} \(\rightarrow\) \emph{bộ điều khiển} \(\rightarrow\) \emph{cơ cấu chấp hành} \(\rightarrow\) \emph{môi trường}.\slidesource{15}
\section{Nút cảm biến IoT (IoT Sensor Node)}
Trong trang 16, bài giảng minh họa cách ghép nhóm \emph{Sensors + Boards + Connectivity} thành \emph{IoT Sensor Nodes}; có ví dụ cảm biến, các bo xử lý, mô-đun truyền thông và bo lắp ráp cuối cùng. Trang 17 đưa sơ đồ khối chi tiết: cảm biến thông minh, cảm biến số, cảm biến video/môi trường; chuyển đổi ADC; MCU; các bộ nhớ flash/RAM; cấp nguồn, DC/DC, quản lý pin; mô-đun truyền thông không dây (Wi-Fi, Bluetooth, IEEE 802.15.4) hoặc có dây. Sơ đồ chỉ ra MCU đóng vai trò phối hợp lấy mẫu, giao tiếp, xử lý và truyền dữ liệu.
\begin{center}\begin{tikzpicture}[node distance=8mm and 7mm, box/.style={draw=navy,rounded corners,fill=blue!4,align=center,minimum height=11mm,minimum width=21mm,font=\footnotesize}, >={Latex}]
\node[box] (s) {Cảm biến};\node[box,right=of s] (adc) {Điều hòa\ tín hiệu / ADC};\node[box,right=of adc] (m) {MCU / CPU};\node[box,right=of m] (n) {Kết nối\ mạng};\node[box,below=of m] (mem) {RAM / Flash};\node[box,below=of adc] (p) {Nguồn, pin,\ DC/DC};
\draw[->] (s)--(adc);\draw[->](adc)--(m);\draw[<->](m)--(n);\draw[<->](m)--(mem);\draw[->](p)--(adc);\draw[->](p.east) -| (m.south);
\end{tikzpicture}\end{center}
\slidefig{17}{Sơ đồ khối nguyên bản của một nút cảm biến IoT}
\section{Vi điều khiển và ATmega328P}
\subsection{Khái niệm microcontroller}
\emph{Microcontroller} (MCU) là một vi mạch tích hợp (IC), có thể hiểu như máy tính thu nhỏ gồm bộ xử lý, bộ nhớ và khối vào/ra. MCU có thể lập trình và được thiết kế cho các ứng dụng nhúng. Trang 18 minh họa các vi mạch Intel P8051, Hitachi HD68000, Motorola MC68HC000LC8. Cần phân biệt \emph{chip MCU} với \emph{bo phát triển} chứa MCU cùng mạch nguồn, giao tiếp và chân cắm.\slidesource{18}
\subsection{Đặc điểm của ATmega328P}
\begin{longtable}{p{.46\linewidth}p{.45\linewidth}}\toprule\textbf{Thông số được nêu trong slide} & \textbf{Giá trị}\\\midrule\endhead
Họ vi điều khiển & AVR 8 bit, hiệu năng cao, công suất thấp\\
Flash chương trình & 32 KB\\
SRAM & 2 KB\\
EEPROM & 1 KB\\
Đường I/O có thể lập trình & 23\\
Điện áp hoạt động & 1,8--5,5 V\\
Dải nhiệt độ & $-40$ đến $85\,^{\circ}\mathrm C$\\
Dòng chế độ Active & 0,2 mA (điều kiện cụ thể không được slide nêu)\\
Dòng chế độ Power-down & 0,1 $\mu$A (điều kiện cụ thể không được slide nêu)\\
Dòng chế độ Power-save & 0,75 $\mu$A (điều kiện cụ thể không được slide nêu)\\\bottomrule\end{longtable}
Hình trang 19 trình bày sơ đồ chân ATmega328P dạng DIP-28, với các nhóm chân: cấp nguồn VCC/AVCC/GND, RESET, dao động XTAL, chân PORTB/PORTC/PORTD, ADC, UART, SPI, I\(^{2}\)C và các chức năng ghép kênh. Trang 21 cho ví dụ mạch lập trình/chuyển chương trình vào chip, đầu nối và bảng pin của giao thức ISP/SPI.\slidesource{19--21}
\slidefig{19}{Sơ đồ 28 chân ATmega328P và các chức năng ghép kênh}
\section{Nền tảng Arduino}
\subsection{Khái quát và lịch sử}
Arduino là nền tảng điện tử \emph{mã nguồn mở} bao gồm cả phần cứng và phần mềm, hướng tới tính dễ dùng. Theo bài giảng, dự án hình thành năm 2003 tại Interaction Design Institute Ivrea (Ý) cho sinh viên; tên gọi bắt nguồn từ quán bar nơi nhóm sáng lập gặp nhau, quán được đặt theo Arduin of Ivrea. Phần cứng ban đầu do Smart Projects tại Ý sản xuất; một số bo mang thương hiệu Arduino được thiết kế bởi các đơn vị khác như SparkFun và Adafruit.\slidesource{22}

Ưu điểm được liệt kê gồm giá thành thấp, đa nền tảng, môi trường lập trình đơn giản, phần mềm mở và phần cứng mở. Các hình giới thiệu họ bo Arduino có nhiều kiểu dáng/kích thước, không phải một bộ xử lý đơn nhất.\slidesource{23--24,31}
\subsection{Arduino Uno và sơ đồ chức năng}
Arduino Uno sử dụng \emph{ATmega328P}, được Arduino.cc phát triển; có 14 chân I/O số và 6 đầu vào analog. Bo có thể lập trình trong Arduino IDE qua cáp USB Type-B; bootloader được nạp sẵn giúp tải chương trình mà không cần thiết bị nạp phần cứng ngoài.

Mười khu vực được chỉ ra trên ảnh bo Uno: (1) các chân digital; (2) LED trên pin 13; (3) LED báo nguồn; (4) ATmega; (5) chân analog; (6) GND và 5V; (7) đầu nối nguồn; (8) LED TX/RX; (9) cổng USB; (10) nút RESET. Sơ đồ pinout ở trang 27 minh họa quan hệ \emph{chân vật lý ATmega328P -- ký hiệu Port -- chức năng Arduino -- đầu vào ADC và các chức năng ngoại vi}.\slidesource{25--27}
\slidefig{26}{Sơ đồ mười khu vực trên bo Arduino Uno}
\slidefig{27}{Bản đồ chân ATmega328P trong hệ Arduino}
\subsection{Các mẫu bo: Leonardo, Mega, Nano}
\begin{longtable}{>{\raggedright\arraybackslash}p{.28\linewidth}p{.20\linewidth}p{.22\linewidth}p{.20\linewidth}}\toprule\textbf{Thông số} & \textbf{Leonardo} & \textbf{Mega} & \textbf{Nano}\\\midrule\endhead
Vi điều khiển & ATmega32U4 & ATmega2560 & ATmega328\\
Điện áp hoạt động & 5 V & 5 V & 5 V\\
Điện áp vào khuyến nghị & 7--12 V & 7--12 V & 7--12 V\\
Giới hạn điện áp vào & 6--20 V & 6--20 V & Không ghi trong slide\\
I/O số & 23 & 54 & 22 (6 PWM)\\
Kênh analog & 12 & 16 & 8\\
Flash & 32 KB, trong đó 4 KB cho bootloader & 256 KB & 32 KB\\
SRAM & 2,5 KB & 8 KB & 2 KB\\
EEPROM & 1 KB & 4 KB & 1 KB\\
Dòng I/O nêu trong slide & 40 mA & 40 mA & Không ghi\\
Dòng chân 3,3 V nêu trong slide & 50 mA & 50 mA & Không ghi\\
Kích thước & 68,6 $\times$ 53,3 mm & Không ghi & 18 $\times$ 45 mm\\
Khối lượng & 20 g & Không ghi & 7 g\\
Thông số bổ sung & -- & -- & 19 mA (tiêu thụ theo slide)\\\bottomrule\end{longtable}
Các mức dòng trong bảng được tái hiện theo slide để phục vụ đối chiếu, không thay thế cho giới hạn thiết kế điện và khuyến cáo của datasheet.\slidesource{28--30}
\subsection{Các lựa chọn ngoài Arduino và shield mở rộng}
Trang 32--34 trình bày \emph{LaunchPad MSP430, Particle Photon, NodeMCU ESP8266, SparkFun Thing Plus, Adafruit Feather HUZZAH, BeagleBoard/PocketBeagle, Seeeduino Nano, SparkFun RedBoard Artemis, STM32F3 Discovery}. Đây là các nền tảng hoặc bo có thể dùng cho nguyên mẫu IoT, không nhất thiết tương thích về phần cứng hay công cụ lập trình.

\emph{Arduino Shield} là bo cắm chồng lên PCB Arduino để bổ sung chức năng, được thiết kế theo hệ chân và triết lý dễ lắp, dễ tháo, dễ sản xuất chi phí thấp. Hình minh họa cho thấy các shield hiển thị, kết nối, nguyên mẫu và mở rộng giao tiếp. Bộ \emph{Arduino Starter Kit} trong trang 36 tập hợp bo mạch, breadboard, linh kiện rời, dây cắm và phụ kiện điện tử.\slidesource{32--36}
\slidefig{34}{Các bo phát triển thay thế được bài giảng liệt kê}
\slidefig{35}{Các dạng Arduino Shield mở rộng}
\subsection{Breadboard và nguyên mẫu}
\emph{Breadboard} là bảng cắm mạch thử nghiệm không cần hàn, thích hợp kiểm chứng nhanh mạch điện và thiết kế trước khi chế tạo PCB. Trang 37 minh họa Arduino nối một LED và điện trở trên breadboard; cần kiểm tra hướng kết nối, hàng/cột chân nối thông và điện trở giới hạn dòng.\slidesource{37}
\slidefig{37}{Mô hình thử nghiệm Arduino với breadboard}
\section{Môi trường và ngôn ngữ lập trình Arduino}
\subsection{Arduino Sketch và Arduino Desktop IDE}
\emph{Sketch} là tên Arduino dùng cho một chương trình, tức đơn vị mã nguồn được biên dịch, nạp và chạy trên bo. Ngôn ngữ Arduino dựa trên các hàm C/C++. Trang 38 mô tả một sketch mẫu và cách tổ chức chương trình thành \texttt{setup()} chạy khởi tạo và \texttt{loop()} thực thi lặp.

Arduino Desktop IDE hỗ trợ Windows, macOS, Linux. Trang 39 dẫn các hướng dẫn cài đặt của Arduino cho từng hệ điều hành, kèm lưu ý cách cài bằng \texttt{sudo apt-get install arduino} trên Linux không được bài giảng khuyến nghị. Trang 40 giới thiệu giao diện IDE, nơi viết sketch, xác minh/biên dịch, chọn bo/cổng giao tiếp, nạp và theo dõi cửa sổ Serial Monitor.\slidesource{38--40}
\slidefig{40}{Giao diện Arduino IDE trong slide gốc}
\subsection{Chú thích chương trình (comments)}
Chú thích được trình biên dịch bỏ qua, dùng để giải thích ý nghĩa mã lệnh. Trang 41 minh họa hai dạng:
\begin{lstlisting}[language=C++]
/* Day la chu thich nhieu dong.
   Cac dong nay duoc bo qua. */
// Day la chu thich mot dong.
\end{lstlisting}
\noindent Lưu ý: chú thích nhiều dòng theo cú pháp \texttt{/* ... */} không được viết lồng nhau; dấu \texttt{//} dùng cho chú thích một dòng.
\subsection{Cấu trúc điều khiển: for, if, while, do--while}
\paragraph{Vòng lặp for (slide 42).} Lặp khối câu lệnh giữa ngoặc nhọn với bộ đếm thường gồm ba biểu thức khởi tạo, điều kiện, bước cập nhật. Biến đếm có thể tăng hoặc giảm.
\begin{lstlisting}[language=C++]
for (int i = 0; i <= 10; i++) {
  Serial.write(7);
}
\end{lstlisting}
\paragraph{Rẽ nhánh if (slide 43).} Kiểm tra biểu thức điều kiện; nếu đúng thì thực hiện nhánh \texttt{if}, nếu sai có thể thực hiện \texttt{else}. Bài giảng minh họa kiểm tra nhiệt độ so với ngưỡng \(70\), đồng thời liệt kê các toán tử \texttt{==}, \texttt{!=}, \texttt{<}, \texttt{<=}, \texttt{>}, \texttt{>=}. Cần phân biệt \texttt{=} (gán) và \texttt{==} (so sánh bằng).
\begin{lstlisting}[language=C++]
if (temperature > 70) {
  // Xu ly khi cao hon nguong
} else {
  // Xu ly khi khong cao hon nguong
}
\end{lstlisting}
\paragraph{Vòng lặp while (slide 44).} Lặp khi điều kiện vẫn đúng; kiểm tra điều kiện trước vòng lặp. \texttt{do\{...\}while(...);} thực hiện thân lặp ít nhất một lần rồi mới kiểm tra điều kiện.
\begin{lstlisting}[language=C++]
int var = 0;
while (var < 200) { var++; }

do { var++; } while (var < 400);
\end{lstlisting}
\slidesource{41--44}
\subsection{Hàm đọc và ghi tín hiệu số}
\texttt{pinMode(pin, mode)} thiết lập chân làm đầu vào hoặc đầu ra (trong slide: \texttt{INPUT}, \texttt{OUTPUT}). \texttt{digitalRead(pin)} đọc trạng thái logic \texttt{HIGH}/\texttt{LOW}. \texttt{digitalWrite(pin, value)} xuất logic \texttt{HIGH} hoặc \texttt{LOW} ra một chân đã được cấu hình phù hợp. Tham số \texttt{pin} là số hiệu chân digital cần tác động, \texttt{value} là mức logic yêu cầu.\slidesource{45}
\begin{lstlisting}[language=C++]
const int LED_PIN = 13;
void setup() { pinMode(LED_PIN, OUTPUT); }
void loop() {
  digitalWrite(LED_PIN, HIGH);
  delay(1000);
  digitalWrite(LED_PIN, LOW);
  delay(1000);
}
\end{lstlisting}
\subsection{Tín hiệu tương tự và PWM}
\texttt{analogRead(pin)} đọc đầu vào ADC (slide mô tả bộ chuyển đổi 10 bit). \texttt{analogWrite(pin, value)} xuất \emph{điều chế độ rộng xung} PWM lên chân hỗ trợ; giá trị duty trong ví dụ Arduino nằm trong khoảng 0--255, tương ứng thời gian xung ở mức cao trong mỗi chu kỳ, chứ không phải một đầu ra điện áp DAC lý tưởng. Hình trang 46 biểu diễn dạng sóng PWM có độ rộng khác nhau.\slidesource{46}
\slidefig{46}{Sơ đồ minh họa nguyên lý PWM và các giá trị analog của Arduino}
\subsection{Các hàm toán học thông dụng}
\begin{longtable}{p{.26\linewidth}p{.64\linewidth}}\toprule\textbf{Hàm} & \textbf{Tác dụng theo bài giảng}\\\midrule\endhead
\texttt{abs(x)} & Trả về trị tuyệt đối.\\
\texttt{constrain(x,a,b)} & Giới hạn \(x\) trong đoạn từ \(a\) đến \(b\).\\
\texttt{max(x,y)} & Chọn giá trị lớn nhất.\\
\texttt{min(x,y)} & Chọn giá trị nhỏ nhất.\\
\texttt{pow(base,exponent)} & Lũy thừa \(\mathrm{base}^{\mathrm{exponent}}\).\\
\texttt{sq(x)} & Bình phương \(x\).\\
\texttt{sqrt(x)} & Căn bậc hai \(\sqrt{x}\).\\\bottomrule\end{longtable}
\slidesource{47}
\subsection{Trì hoãn và chỉ thị tiền xử lý}
\texttt{delay(milliseconds)} tạm dừng chương trình trong khoảng thời gian tính theo mili-giây, \(1000\,\mathrm{ms}=1\,\mathrm{s}\). \texttt{\#define} đặt tên hằng hoặc macro trước khi chương trình được biên dịch, ví dụ \texttt{\#define ledPin 3}. \texttt{\#include} thêm tiêu đề thư viện, ví dụ \texttt{\#include <SPI.h>}; nhờ đó sketch truy cập các hàm, lớp và định nghĩa đã được thư viện cung cấp.\slidesource{48--49}
\begin{lstlisting}[language=C++]
#define ledPin 3
#include <SPI.h>
\end{lstlisting}
\subsection{Hai hàm chính setup() và loop()}
\texttt{setup()} được gọi một lần mỗi khi khởi động hoặc reset để thiết lập chân, Serial, thư viện và các trạng thái ban đầu. \texttt{loop()} được gọi lặp liên tục sau khi \texttt{setup()} hoàn tất, là nơi đặt chương trình chính xử lý tín hiệu và điều khiển thiết bị. Slide 50 dùng \texttt{pinMode(8, INPUT)} minh họa thao tác khởi tạo; slide 51 minh họa dùng \texttt{digitalRead(8)} và ghi chữ qua Serial theo trạng thái nút nhấn.\slidesource{50--51}
\begin{lstlisting}[language=C++]
void setup() {
  pinMode(8, INPUT);
  Serial.begin(9600);
}
void loop() {
  if (digitalRead(8) == HIGH) {
    Serial.write('H');
  } else {
    Serial.write('L');
  }
}
\end{lstlisting}
\subsection{Thêm thư viện và giao tiếp Serial}
Trang 52 minh họa thao tác \emph{Sketch > Include Library} trong Arduino IDE. Các thư viện có thể đi kèm IDE hoặc được thêm bởi người dùng. Cổng COM là tên truyền thống của cổng giao tiếp nối tiếp trên máy tính tương thích IBM PC, ngày nay còn bao gồm cổng nối tiếp ảo thông qua Bluetooth hoặc cầu USB-to-UART. Arduino thường cung cấp ít nhất một cổng nối tiếp, một số bo có nhiều cổng. Trên Arduino Uno, \textbf{pin 0 là RX và pin 1 là TX}; người lập trình thường dùng \texttt{Serial.begin(speed)} khởi tạo với tốc độ baud và \texttt{Serial.write(value)} ghi dữ liệu ra cổng; baud là số ký hiệu truyền mỗi giây (bài giảng mô tả gần đúng theo bit/giây).\slidesource{52--53}
\slidefig{52}{Thao tác thêm thư viện trong Arduino IDE}
\section{Nền tảng Wi-Fi ESP8266}
\subsection{Tổng quan ESP8266EX và khả năng phần mềm}
ESP8266EX là chip Wi-Fi chi phí thấp của Espressif Systems, tích hợp giao thức TCP/IP và có thể ghép vào các thiết kế dùng vi điều khiển. Phần lõi hỗ trợ Arduino của ESP8266 cung cấp thư viện Wi-Fi cho TCP, UDP, HTTP, mDNS, SSDP, máy chủ DNS; cho phép cập nhật firmware qua mạng (OTA), hệ thống tệp trên flash hoặc SD, điều khiển servo và trao đổi SPI/I\(^{2}\)C. Slide 55 liệt kê các biến thể mô-đun ESP-01, ESP-02, ESP-03, ESP-04, ESP-05, ESP-06, ESP-07, ESP-08, ESP-09, ESP-10, ESP-11, ESP-12, ESP-13.\slidesource{54--55}
\slidefig{55}{Tập hợp biến thể mô-đun ESP8266}
\subsection{ESP-01 và ESP-07}
\begin{longtable}{p{.34\linewidth}p{.57\linewidth}}\toprule\textbf{Mô-đun} & \textbf{Đặc điểm được cung cấp trong slide}\\\midrule\endhead
ESP-01 & Kích thước $24{,}75\times14{,}5$ mm; nguồn 3,3 V; anten mạch in tích hợp; LED nguồn và TXD; chế độ AP, STA, AP+STA; có thể điều khiển bằng lệnh AT; khoảng cách chân 2,54 mm.\\
ESP-07 & Nguồn 3,3 V; anten gốm tích hợp và đầu nối anten ngoài U.FL; Wi-Fi 802.11b/g/n băng 2,4 GHz với WPA/WPA2; TCP/UDP; hỗ trợ lệnh AT; khoảng cách chân 2 mm.\\\bottomrule\end{longtable}
\slidesource{56--57}
\slidefig{56}{Sơ đồ chân, module ESP-01 và đặc điểm ngoại hình}
\slidefig{57}{Sơ đồ chân ESP-07 trong bài giảng}
\subsection{SDK của Espressif và môi trường thay thế}
Bài giảng nêu năm 2014 Espressif giới thiệu SDK cho phép ESP8266 chạy như vi điều khiển độc lập thay vì chỉ là mô-đun Wi-Fi của MCU khác. Hai nhánh SDK được nhắc đến là \emph{ESP8266 NonOS SDK} và \emph{ESP8266 RTOS SDK}. Các môi trường hoặc SDK thay thế trên slide gồm ESP-Open-SDK, Unofficial Development Kit, Arduino, ESP8266 BASIC, ESP Easy và ESPHome.\slidesource{58}
\subsection{Cấu hình ESP8266 trên Arduino IDE}
Ba trang hướng dẫn tuần tự: \textbf{(1)} vào \emph{File > Preferences} và thêm URL gói bo \url{http://arduino.esp8266.com/stable/package_esp8266com_index.json}; \textbf{(2)} vào \emph{Tools > Board > Boards Manager}, tìm ``esp8266'' và chọn \emph{Install}; \textbf{(3)} vào \emph{Tools > Board} chọn \emph{Generic ESP8266 Module} nếu đúng loại mô-đun cần sử dụng. Các ảnh minh họa trong slide là giao diện IDE tại thời điểm biên soạn; vị trí menu có thể khác ở bản IDE mới.\slidesource{59--61}
\slidefig{59}{Thêm liên kết JSON của gói bo ESP8266}
\slidefig{60}{Cài gói hỗ trợ ESP8266 qua Boards Manager}
\slidefig{61}{Chọn bo Generic ESP8266 Module}
\section{Nền tảng ESP32}
ESP32 là họ \emph{system-on-chip} giá tương đối thấp do Espressif phát triển, là bước tiến so với họ ESP8266. Slide giới thiệu khả năng đa lõi, đồng xử lý tiết kiệm năng lượng và hỗ trợ kết nối Wi-Fi, Bluetooth. Những thiết bị minh họa gồm ESP32 Development Board, ESP32-CAM, ESP32-WROOM-32U/-32D, DOIT DevKit V1, ESP32 DevKit, ESP32S NodeMCU, ESP32 Thing, Wemos Lolin32, Wemos OLED và HUZZAH32.

Các \textbf{ưu điểm theo slide}: CPU dual-core khoảng 160--240 MHz; Wi-Fi/Bluetooth tiện cho điều khiển và giám sát; số chân GPIO nhiều hơn ESP8266; slide ghi tốc độ tối đa 150 Mb/s. \textbf{Nhược điểm theo slide}: giá cao hơn ESP8266. Thông số là mô tả cho các phiên bản ESP32 được trình bày trong bài giảng, không đại diện cho mọi biến thể ESP32. Đặc biệt, ảnh trang 64 minh họa pinout bo DOIT ESP32 DevKit V1 loại 30 chân, gồm các chân cấp nguồn, EN, chân UART, ADC, digital GPIO, SPI và I\(^{2}\)C tùy cấu hình.\slidesource{62--65}
\slidefig{63}{Một số bo ESP32 được giới thiệu trong slide}
\slidefig{64}{Sơ đồ chân bo DOIT ESP32 DevKit V1 (30 GPIOs theo tiêu đề slide)}
\section{Nền tảng Intel Edison}
\subsection{Cấu hình mô-đun Edison}
Intel Edison trong bài giảng có kích thước $35{,}5\times25\times3{,}9$ mm, linh kiện gắn ở hai mặt. SoC Intel Atom Tangier Z34XX tiến trình 22 nm tích hợp hai lõi Silvermont 500 MHz và một lõi Intel Quark ở 100 MHz (đơn vị trong slide không ghi đầy đủ tại dòng này); RAM 1 GB và eMMC 4 GB; có Wi-Fi, Bluetooth 4, điều khiển USB. Kết nối ra ngoài qua đầu nối Hirose DF40 mật độ cao 70 chân (USB, SD, UART, GPIO). Các môi trường lập trình nêu trong slide: Arduino IDE, Eclipse (C, C++, Python), Intel XDK (NodeJS/HTML5).\slidesource{66}
\slidefig{66}{Mô-đun Intel Edison và cấu hình phần cứng}
\subsection{Bo breakout và bo dạng Arduino cho Edison}
\textbf{Breakout board}: xuất các mức I/O gốc 1,8 V của Edison ra các điểm hàn dạng lưới 0,1 inch; có USB OTG Micro Type-AB, công tắc nguồn OTG, bộ sạc pin, cầu USB--UART qua Micro-USB, đầu cấp nguồn DC 7--15 V. Điện áp I/O 1,8 V là thông tin quan trọng khi giao tiếp với ngoại vi.

\textbf{Arduino board for Edison Module}: 20 I/O số (6 PWM), 6 đầu vào analog, 1 UART, 1 I\(^{2}\)C, đầu ICSP 6 chân cho SPI, cổng Micro-USB device hoặc cổng USB host Type-A chọn qua công tắc cơ, cổng Micro-USB UART, khe SD và đầu nguồn DC 7--15 V.\slidesource{67--68}
\slidefig{67}{Intel Edison Breakout Board}
\slidefig{68}{Bo tương thích bố cục Arduino cho Edison}
\section{Điện toán biên (Edge Computing)}
Slide 69 trình bày sự khác nhau về vị trí xử lý giữa đám mây, trung tâm dữ liệu doanh nghiệp và biên gần thiết bị; trang 70 nhấn mạnh \emph{xử lý thông tin nhạy cảm với độ trễ}, \emph{quyền riêng tư và an ninh}, \emph{băng thông và chất lượng kết nối}. Trang 71 mô tả dữ liệu do thiết bị IoT tạo ra được xử lý một phần ở máy hoặc gateway gần nguồn, trong khi đám mây và trung tâm dữ liệu tiếp tục thực hiện các công việc ở quy mô lớn. Khác với gửi tất cả dữ liệu thô lên cloud, cách bố trí này có thể giảm lượng dữ liệu truyền và cho phép đưa ra quyết định cục bộ.

Các \textbf{thiết bị biên} minh họa trang 72 gồm Raspberry Pi, NVIDIA Jetson Nano và NVIDIA Jetson Xavier. Liên kết tham khảo trên slide: \url{https://www.networkworld.com/article/3224893/what-is-edge-computing-and-how-it-s-changing-the-network.html}.\slidesource{69--72}
\slidefig{69}{Quan hệ giữa cloud, edge gateway và trung tâm dữ liệu}
\slidefig{71}{Sơ đồ hoạt động của điện toán biên}
\slidefig{72}{Các máy tính biên: Raspberry Pi và NVIDIA Jetson}
\section{Máy tính đơn bo Raspberry Pi}
\subsection{Khái quát}
Raspberry Pi là dòng \emph{single-board computers} (SBC), tức máy tính tích hợp phần lớn thành phần trên một bo mạch nhỏ gọn. Theo số liệu lịch sử \textbf{được trích trong slide}, đến tháng 2/2015 hơn 5 triệu máy Raspberry Pi đã được bán, trở thành máy tính Anh bán chạy nhất thời điểm đó. Dòng bo do Raspberry Pi Foundation phát triển nhằm khuyến khích học khoa học máy tính cơ bản trong trường học và tại những quốc gia đang phát triển. Slide nêu các bo có Wi-Fi/Bluetooth tích hợp; khả năng này tùy theo phiên bản thực tế. Hình trang 73 minh họa các dự án robot và bộ phụ kiện Raspberry Pi.\slidesource{73--74}
\slidefig{73}{Hệ sinh thái thiết bị và ứng dụng Raspberry Pi}
\subsection{Bộ kit nhập môn}
Trang 75 liệt kê các thành phần: bo Raspberry Pi, thẻ microSD dung lượng \textbf{8 GB trở lên} theo slide, bộ cấp nguồn, bàn phím/chuột USB, màn hình có ngõ HDMI, cáp HDMI và hộp vỏ bảo vệ. Thẻ microSD làm bộ nhớ khởi động và lưu hệ điều hành đối với các cấu hình được minh họa trong bài giảng.\slidesource{75}
\slidefig{75}{Bộ kit Raspberry Pi gồm bo, microSD, nguồn và ngoại vi}
\subsection{Cấu tạo và chân Raspberry Pi 4}
Trang 76 chỉ các phần \textbf{USB-C power}, \textbf{micro-HDMI} (hai đầu ra), \textbf{Ethernet Gigabit}, cổng \textbf{USB 3.0/2.0}, header GPIO 40 chân và lựa chọn RAM 2/4/8 GB trong hình. Trang 77 chứa bảng pinout 40 chân: nguồn 3,3 V, 5 V, GND, các chân GPIO và các chân chức năng SPI/UART/I\(^{2}\)C. Khi đấu nối phải đối chiếu số thứ tự chân vật lý (physical pin) và tên GPIO/BCM, tránh cấp 5 V vào tín hiệu GPIO không chịu 5 V.\slidesource{76--77}
\slidefig{76}{Các cổng kết nối và khu vực chính của Raspberry Pi 4}
\slidefig{77}{Sơ đồ chân header 40 pin Raspberry Pi 4}
\section{So sánh Arduino Uno và Raspberry Pi 4}
\begin{longtable}{p{.38\linewidth}p{.26\linewidth}p{.26\linewidth}}\toprule\textbf{Tiêu chí trong bài giảng} & \textbf{Arduino Uno} & \textbf{Raspberry Pi 4}\\\midrule\endhead
Loại nền tảng & Vi điều khiển (microcontroller) & Máy vi tính (microcomputer)\\
Hệ điều hành & Không chạy hệ điều hành đa nhiệm thông dụng & Có thể chạy hệ điều hành\\
Đa nhiệm & Slide ghi ``No multitasking'' & Có đa nhiệm\\
Giá thành tương đối & Rẻ & Đắt hơn\\
Tài nguyên & Xung 16 MHz, SRAM 2 KB & Xung khoảng 1,5 GHz, RAM 2/4/8 GB\\
Ứng dụng thời gian thực & Được slide đánh giá là thích hợp & Slide đánh giá không lý tưởng cho thời gian thực cứng\\
Kết nối mạng tích hợp & Slide ghi không có kết nối mạng & Có khả năng kết nối mạng\\
ADC tích hợp trên bo & Có & Slide ghi không có bộ chuyển đổi A/D\\\bottomrule\end{longtable}
Dòng so sánh phản ánh \textbf{đúng hai bo được nêu trong slide 78}, không phải khẳng định cho toàn bộ họ Arduino và Raspberry Pi. Một số chức năng có thể mở rộng bằng shield, HAT hoặc IC ngoài; với Raspberry Pi, yêu cầu đo analog thường cần thêm ADC bên ngoài.\slidesource{78}
\slidefig{78}{Bảng so sánh Arduino Uno và Raspberry Pi 4 nguyên bản}
\section{Bảng tra cứu tương ứng với toàn bộ 79 slide}
\begin{longtable}{p{.15\linewidth}p{.76\linewidth}}\toprule\textbf{Trang PDF} & \textbf{Chủ đề và hình minh họa được lưu trong phụ lục gốc}\\\midrule\endhead
1--3 & Bìa giảng viên, mục lục chương và kiến trúc Cloud--Edge--IoT Devices.\\
4--6 & Khái niệm cảm biến, hàm chuyển đổi, ảnh mô-đun cảm biến đa dạng.\\
7--12 & Chủ động/thụ động; phân loại theo hiện tượng và analog/digital; dạng sóng; cảm biến tương tự và cảm biến số.\\
13--15 & Cơ cấu chấp hành, hình các thiết bị đầu ra, so sánh cảm biến/cơ cấu chấp hành.\\
16--17 & Phân rã phần cứng nút IoT và sơ đồ khối IoT Sensor Node.\\
18--21 & Khái niệm MCU, chip ATmega328P, cấu hình và bộ nạp.\\
22--27 & Lịch sử/ưu điểm Arduino, bo Uno, các vùng chức năng và pinout ATmega328P.\\
28--31 & Arduino Leonardo, Mega, Nano và ảnh các biến thể Arduino.\\
32--35 & Bo thay thế Arduino và shield mở rộng.\\
36--40 & Starter Kit, breadboard, Arduino Sketch, cài đặt và giao diện Arduino IDE.\\
41--44 & Comment và cấu trúc \texttt{for}, \texttt{if}, \texttt{while}, \texttt{do--while}.\\
45--48 & Digital I/O, Analog I/O và PWM, hàm toán học và delay.\\
49--53 & \texttt{\#define}, \texttt{\#include}, \texttt{setup()}, \texttt{loop()}, thư viện và Serial.\\
54--58 & ESP8266, các biến thể mô-đun, đặc điểm ESP-01, ESP-07, SDK.\\
59--61 & Quy trình cài đặt board ESP8266 trong Arduino IDE với ảnh các bước.\\
62--65 & ESP32, các bo ESP32, sơ đồ chân DOIT và ưu/nhược điểm.\\
66--68 & Intel Edison, bo breakout và bo giao tiếp phong cách Arduino.\\
69--72 & Điện toán biên, mô hình hoạt động và các bo edge devices.\\
73--77 & Raspberry Pi, bộ kit, sơ đồ các cổng và pinout Pi 4.\\
78--79 & So sánh Uno và Pi 4, lời kết.\\\bottomrule\end{longtable}
\section{Kết luận Chương 2}
Phần cứng IoT được tổ chức xung quanh \emph{cảm biến} (thu nhận dữ liệu), \emph{cơ cấu chấp hành} (tạo tác động), \emph{bộ điều khiển/tính toán} (quản lý, xử lý, điều khiển), \emph{giao diện I/O}, \emph{khối nguồn} và \emph{khối truyền thông}. Với các tác vụ điều khiển nhúng, Arduino/ATmega và ESP8266/ESP32 là những ví dụ tiêu biểu của MCU. Với các tác vụ cần hệ điều hành, đa nhiệm, xử lý ở biên và giao tiếp phần mềm đa dạng, Raspberry Pi hoặc một bo công nghiệp có thể phù hợp hơn. Intel Edison được bài giảng dùng làm minh họa một kiến trúc mô-đun máy tính nhúng kết hợp nhiều giao diện ngoại vi.

\textbf{Điểm cần ghi nhớ:} lựa chọn không chỉ theo xung CPU, mà theo số chân, điện áp I/O, ADC, dung lượng RAM/Flash, giao tiếp không dây, nguồn, phần mềm công cụ và độ trễ xử lý mong muốn. Các ví dụ, bảng chân và ảnh cấu trúc ở phụ lục giúp đối chiếu thông tin thực tế của chính từng bo, không phải thông số chuẩn cho toàn bộ họ chip.
\ifdefined\iotstandalone
\clearpage
\appendix
\chapter{Phụ lục: 79 trang slide Chương 2 nguyên bản}
Phụ lục đính kèm \textbf{đầy đủ 79 trang} tài liệu do người học cung cấp, nhằm bảo toàn các ảnh chụp phần cứng, sơ đồ chân, bảng so sánh, ví dụ mã và ảnh giao diện mà phần diễn giải bằng chữ không thể thay thế hoàn toàn. Bố cục slide gốc không bị chỉnh sửa. Khi ghép các chương, có thể bỏ qua phụ lục bằng cách không sao chép đoạn \texttt{\textbackslash includepdf} bên dưới.
\IfFileExists{assets/chapter2_original_slides.pdf}{%
  \includepdf[pages=1-79,fitpaper=false,offset=0 0,pagecommand={},width=.95\paperwidth]{assets/chapter2_original_slides.pdf}%
}{\begin{noteiot}Không có tệp \texttt{assets/chapter2\_original\_slides.pdf} trong thư mục dự án. Module \texttt{chapter2.tex} vẫn chứa các phần diễn giải và bảng tra cứu; để hiện ảnh slide nguyên bản hãy đặt PDF nguồn tại vị trí này rồi biên dịch lại.\end{noteiot}}
\fi % Phụ lục gốc chỉ xuất hiện khi biên dịch riêng Chương 2.

% % CHỈ CHƯƠNG 3. Không chứa \documentclass hoặc \begin{document}.
% Sử dụng trong main.tex: % CHỈ CHƯƠNG 3. Không chứa \documentclass hoặc \begin{document}.
% Sử dụng trong main.tex: % CHỈ CHƯƠNG 3. Không chứa \documentclass hoặc \begin{document}.
% Sử dụng trong main.tex: \input{chapter3.tex}
\chapter{Công nghệ mạng trong Internet vạn vật}
\label{chap:iot-networking}
\begin{ghinho}[title={Mục tiêu và phạm vi}]
Chương này tổng hợp bốn phần của 81 trang slide: (1) giao tiếp trong IoT; (2) công nghệ và giao thức truyền thông, gồm định danh, địa chỉ hóa, công nghệ truyền dẫn, bảo mật và định tuyến; (3) phân loại mạng IoT; (4) giao thức dữ liệu HTTP, CoAP, MQTT, WebSocket, AMQP và XMPP. Các bảng thông số được ghi theo chính slide, không được hiểu là mức cam kết cho mọi phiên bản, thiết bị hoặc quốc gia. 
\end{ghinho}

\section{Giao tiếp trong hệ thống IoT}
\subsection{Khái niệm và vai trò của truyền thông}
Internet vạn vật (Internet of Things, IoT) kết nối các thiết bị thông minh qua mạng, cho phép chúng trao đổi dữ liệu và thực hiện công việc với ít hoặc không cần sự can thiệp trực tiếp của con người. Những thiết bị này thường tích hợp linh kiện điện tử, cảm biến, phần mềm, kết nối mạng và cơ cấu chấp hành. Truyền thông là điều kiện để dữ liệu đo ở một thiết bị trở thành thông tin có ích cho các thiết bị, dịch vụ và người sử dụng ở nơi khác. Không có liên kết đáng tin cậy giữa các thiết bị và nền tảng xử lý, một hệ thống cảm biến rời rạc chưa phát huy đầy đủ giá trị IoT. \sourcepage{3}

\subsection{Bốn hướng tương tác trong IoT}
\paragraph{Người đến máy (Human-to-Machine, H2M).} Con người gửi đầu vào cho hệ thống dưới dạng giọng nói, văn bản, hình ảnh hoặc thao tác trực tiếp; thiết bị thu nhận, phân tích và đáp ứng bằng văn bản, hình ảnh hoặc hành động. H2M phối hợp phần cứng và phần mềm để hỗ trợ các công việc hằng ngày. Ví dụ: nhận dạng giọng nói; nhận dạng khuôn mặt; chấm công sinh trắc học. \sourcepage{4}

\paragraph{Máy đến máy (Machine-to-Machine, M2M).} Hai hoặc nhiều máy trao đổi dữ liệu hoặc lệnh tự động, không đòi hỏi con người tham gia vào mỗi phiên liên lạc. Đường truyền có thể là hữu tuyến hoặc vô tuyến; một kết nối M2M điểm--điểm có thể sử dụng Ethernet hoặc mạng di động. IoT tiếp thu cách vận hành tự động của M2M và mở rộng sang các mạng thiết bị quy mô lớn có dịch vụ/nền tảng đám mây. Ví dụ: máy bán hàng tự động thông báo hết hàng cho nhà phân phối; máy giặt thông minh gửi thông báo hoàn tất; công tơ thông minh theo dõi điện năng và cảnh báo chủ sử dụng. \sourcepage{5--6}

\paragraph{Máy đến người (Machine-to-Human, M2H).} Thiết bị chủ động phát tín hiệu, văn bản, hình ảnh hoặc lời nói hướng tới con người, kể cả khi người đó không trực tiếp phát yêu cầu. Ví dụ: chuông báo cháy, đèn giao thông, thiết bị giám sát sức khỏe. \sourcepage{7}

\paragraph{Người đến người (Human-to-Human, H2H).} Con người liên lạc bằng lời nói, chữ viết, hình vẽ, nét mặt hoặc ngôn ngữ cơ thể. Trong một hệ thống IoT thực tế, H2H vẫn cần thiết để phối hợp xử lý sự cố, quản lý tình huống và đưa ra quyết định khi tự động hóa M2M chưa đáp ứng đủ. \sourcepage{8}

\begin{figure}[htbp]\centering
\begin{tikzpicture}[>=Latex,box/.style={draw=navy,rounded corners,fill=blue!5,minimum width=36mm,minimum height=13mm,align=center},node distance=20mm]
\node[box](human){Con người};\node[box,right=36mm of human](machine){Thiết bị IoT};
\draw[->,thick,teal,bend left=20](human) to node[above]{H2M: lệnh/đầu vào}(machine);
\draw[->,thick,orangeiot,bend left=20](machine) to node[below]{M2H: cảnh báo/kết quả}(human);
\node[box,below=19mm of machine](node2){Thiết bị IoT khác};
\draw[<->,thick,navy](machine)--node[right]{M2M}(node2);
\node[box,below=19mm of human](human2){Con người khác};
\draw[<->,thick,navy](human)--node[left]{H2H}(human2);
\end{tikzpicture}
\caption{Minh họa bốn loại hình tương tác H2M, M2M, M2H, H2H (vẽ lại theo các slide 4--8).}\end{figure}

\section{Công nghệ và giao thức truyền thông IoT}
\subsection{Tổng quan và yêu cầu lựa chọn}
Một hệ thống IoT có thể dùng Bluetooth, Wi-Fi, Wi-Fi Direct, giao thức vô tuyến công suất thấp, mạng di động LTE/LTE-A hoặc các công nghệ khác. Không có giao thức tốt nhất cho mọi ứng dụng: phải cân nhắc tầm phủ sóng, mức tiêu thụ năng lượng, băng thông, độ trễ, chi phí, tính di động, quy mô nút và mức bảo mật. Slide giới thiệu các họ công nghệ: PLC, X10, NFC, UWB, Wi-Fi, BLE, Zigbee, Z-Wave, Weightless, Sigfox, LoRaWAN, Li-Fi, mạng tế bào/5G, LTE-M và NB-IoT. \sourcepage{2,9,18}

\subsection{Định danh, đặt tên và địa chỉ hóa}
\paragraph{Định danh (identification).} Khi một nút gia nhập mạng, nó cần được nhận biết bằng một định danh. Cơ chế định danh giúp quản lý thiết bị, xác định đối tượng giao tiếp và hạn chế xử lý lưu lượng không cần thiết trong mạng rất lớn. Hệ thống còn cần tên và địa chỉ để tìm đúng thiết bị, đúng tài nguyên. Lưu ý về mặt kỹ thuật: việc cấp định danh tự nó không chặn được tấn công hay lưu lượng tràn ngập; phải kết hợp xác thực, phân quyền và cơ chế điều khiển lưu lượng. \sourcepage{10}

\paragraph{Đặt tên tài nguyên và đối tượng.} Ba cơ chế được nêu trong slide gồm:
\begin{itemize}
\item \textbf{URI (Uniform Resource Identifier):} chuỗi nhận dạng tài nguyên trên mạng. Ví dụ trong slide: \texttt{http://www.example.com/temperature/}.
\item \textbf{EPC (Electronic Product Code):} mã định danh đối tượng vật lý, thường mã hóa trên thẻ RFID thụ động. Ví dụ: \texttt{urn:epc:id:sgtin:}\allowbreak\texttt{0614141.012345.62852}.
\item \textbf{uCode (Ubiquitous Code):} mã định danh của hệ thống định danh đối tượng; ví dụ slide: \texttt{00001C0000000000}\allowbreak\texttt{0001000285E7A6E3}. Trang giới thiệu: \url{https://www.uidcenter.org}.
\end{itemize}
\sourcepage{11}

\paragraph{IPv4.} Giao thức Internet phiên bản 4 dùng địa chỉ 32 bit, biểu diễn thường gặp bằng bốn nhóm số thập phân phân tách bằng dấu chấm; tổng không gian lý thuyết là $2^{32}$ (xấp xỉ 4,3 tỷ) địa chỉ. Địa chỉ có thể chia mạng con, định tuyến và được cấu hình thủ công hoặc qua DHCP. Ví dụ slide minh họa dạng \texttt{192.168.5.18}. \sourcepage{12,14}

\paragraph{IPv6.} Giao thức Internet phiên bản 6 dùng địa chỉ 128 bit, tương ứng $2^{128}\approx 3{,}4\times 10^{38}$ địa chỉ; ví dụ ký hiệu rút gọn sử dụng các nhóm số hệ 16 và dấu hai chấm. Mục tiêu thiết kế bao gồm giải quyết hạn chế không gian địa chỉ của IPv4 và hỗ trợ số lượng rất lớn thiết bị. IPv6 có cơ chế tự cấu hình địa chỉ, hỗ trợ multicast và khả năng phối hợp với các cơ chế bảo vệ như IPsec. \sourcepage{13--17}

\begin{table}[htbp]\centering\small
\caption{So sánh IPv4 và IPv6 theo minh họa các slide 12--17.}
\begin{tabularx}{\textwidth}{p{3.15cm}XX}\toprule
Tiêu chí & IPv4 & IPv6\\\midrule
Độ dài địa chỉ & 32 bit & 128 bit\\
Số địa chỉ lý thuyết & $2^{32}\approx4{,}3\cdot10^9$ & $2^{128}\approx3{,}4\cdot10^{38}$\\
Dạng địa chỉ & 4 nhóm số thập phân (ví dụ 192.168.5.18) & Các nhóm số hệ 16, phân cách bằng dấu hai chấm\\
Cấu hình & Thủ công hoặc DHCP & Hỗ trợ tự cấu hình; có thể dùng DHCPv6\\
Ứng dụng trong IoT & Có thể bị hạn chế bởi số địa chỉ công cộng & Tạo không gian địa chỉ rất lớn cho nhiều nút\\\bottomrule
\end{tabularx}\end{table}

\paragraph{Các lợi ích IPv6 mà slide liệt kê.} (1) Dễ nhận diện và trao đổi dữ liệu giữa nhiều thiết bị hơn nhờ không gian địa chỉ lớn; (2) tự cấu hình đơn giản hơn ở quy mô lớn, có thể giảm công cấu hình và chi phí triển khai; (3) hỗ trợ truyền tải dữ liệu và hiệu quả sử dụng băng thông; (4) hỗ trợ các thuộc tính bảo mật về bí mật, xác thực và toàn vẹn khi cấu hình các cơ chế phù hợp; (5) multicast giúp tránh một số thao tác phát quảng bá thường xuyên, giảm số gói phải xử lý và có thể tiết kiệm pin; (6) cơ chế định tuyến và cấu trúc địa chỉ linh hoạt theo cách tổ chức mạng/ứng dụng. Đây là các lợi ích tiềm năng theo slide, không phải bảo đảm rằng bật IPv6 tự động tăng tốc hay bảo mật toàn bộ hệ thống. \sourcepage{16--17}

\begin{ghinho}[title={Phân biệt ba khái niệm}]
\textbf{Tên} cho biết tài nguyên hoặc đối tượng nào cần tìm; \textbf{định danh} phân biệt đối tượng trong hệ thống; \textbf{địa chỉ mạng} giúp chuyển gói đến vị trí mạng đang sử dụng. Tên, mã RFID và địa chỉ IP vì thế có thể cùng tồn tại trên một thiết bị IoT.
\end{ghinho}

\subsection{Công nghệ truyền dẫn qua dây điện: PLC và X10}
\paragraph{PLC (Power-Line Communication).} PLC sử dụng ngay đường dây điện hiện có để mang cả điện năng xoay chiều và tín hiệu dữ liệu. Một nút có thể nhận nguồn điện và trao đổi thông tin mà không cần tuyến dây dữ liệu riêng. Phù hợp cho thiết bị cố định, nhà thông minh, lưới điện thông minh, phương tiện hoặc môi trường có nhiễu vô tuyến cao. Hạn chế thực tế phụ thuộc chất lượng dây, tải và nhiễu từ thiết bị điện. \sourcepage{19}

\paragraph{X10.} Chuẩn điều khiển thiết bị trên đường điện ra đời khoảng thập niên 1970. Tín hiệu điều khiển được đưa lên dây điện bằng các xung sóng mang RF ngắn. Tốc độ và phạm vi thấp so với một số hệ PLC khác nhưng chi phí thấp; ứng dụng chính là bật/tắt, điều chỉnh đèn và điều khiển ổ cắm trong nhà thông minh qua bộ điều khiển hoặc ứng dụng. \sourcepage{20}

\begin{longtable}{p{.38\textwidth}p{.55\textwidth}}
\caption{Các thông số X10 theo bảng ở slide PDF trang 21.}\\\toprule
Thông số & Giá trị trong slide\\\midrule\endfirsthead
\toprule Thông số & Giá trị trong slide\\\midrule\endhead
Môi trường truyền & Đường dây điện (powerline only)\\
Loại thiết bị & Bộ điều khiển; thiết bị bị điều khiển; thiết bị có báo trạng thái\\
Điều chế lớp vật lý & OOK, sóng mang 120 kHz\\
Tốc độ dữ liệu & 120 bps tức thời; 60 bps duy trì; 24 bps tải hữu ích\\
Địa chỉ hóa & 16 mã nhà A--P và 16 mã thiết bị 1--16, tối đa 256 địa chỉ\\
Lệnh & 16 lệnh, như ON, OFF, DIM, BRIGHT, ALL UNITS OFF, ALL LIGHTS ON\\
Độ dài thông điệp & Slide ghi 100 bit (8,5 byte); hai trị số không tương đương nên cần kiểm tra nếu dùng thiết kế bản tin\\
Mã hóa & Manchester\\
Cấu trúc gói & Start Code 4 bit; House Code 8 bit; Key Code 10 bit (theo slide)\\
ACK và bảo mật & Không hỗ trợ theo bảng slide\\\bottomrule
\end{longtable}

\subsection{Giao tiếp khoảng cách rất ngắn: NFC và UWB}
\paragraph{NFC (Near-Field Communication).} Tập hợp giao thức vô tuyến cho các thiết bị tiếp xúc/đặt gần nhau, hướng tới xác thực, đọc/ghi thẻ, chia sẻ dữ liệu và kiểm soát ra vào; ví dụ thanh toán thông minh. Khác với cách dùng RFID chỉ để định danh trong nhiều hệ thống, NFC có các chế độ giao tiếp cho cả thiết bị và thẻ. \sourcepage{22}

\begin{longtable}{p{.47\textwidth}p{.45\textwidth}}
\caption{Thông số và chế độ NFC theo slide PDF trang 23.}\\\toprule
Đặc tính & Giá trị được nêu\\\midrule\endfirsthead
\toprule Đặc tính & Giá trị được nêu\\\midrule\endhead
Tần số mang & 13,56 MHz\\
Khoảng cách & Dưới 10 cm (bảng đối chiếu trang 42 sử dụng 4 cm)\\
Tốc độ & 106, 212 hoặc 424 kbit/s\\
Đối tượng & Thẻ (tag) và đầu đọc (reader)\\
Loại thẻ & Type 1 đến Type 5\\
Cấu hình mạng & Peer-to-peer\\
Chế độ thiết bị & Card emulation; reader/writer; peer-to-peer (active/passive)\\
Cách giao tiếp & Active--passive hoặc active--active\\
Thiết lập kết nối & Bảng slide ghi 'vài giây' và 'xấp xỉ dưới 0,1 giây'; cần xác định phép đo cụ thể\\
Mã hóa đường truyền & NRZ-L, Manchester, Modified Miller\\
Điều chế & ASK, BPSK (theo bảng)\\
Giải quyết va chạm & Có cơ chế chống va chạm\\
Tiêu chuẩn & ISO 14443A/B, ISO 18092, JIS X6319-4\\\bottomrule
\end{longtable}

\paragraph{UWB (Ultra-Wideband).} Vô tuyến sử dụng phổ tần rất rộng và xung phát năng lượng thấp để truyền dữ liệu và/hoặc định vị. Theo slide, khả năng hạn chế can nhiễu và theo dõi vị trí từng nút hỗ trợ mở rộng hệ thống; ứng dụng nổi bật là dịch vụ định vị và theo dõi tài sản. Không đồng nhất UWB với định vị GPS hay với mạng LPWAN: công năng và môi trường triển khai khác nhau. \sourcepage{24}

\subsection{Wi-Fi và Wi-Fi HaLow}
Wi-Fi phổ biến trong kết nối mạng cục bộ và có thể thiết lập kết nối ad hoc như Wi-Fi Direct mà không cần access point. Điểm bất lợi với cảm biến chạy pin là công suất tiêu thụ, dù nó không quan trọng bằng đối với nhiều thiết bị nhà thông minh có nguồn ổn định. \term{Wi-Fi HaLow}{IEEE 802.11ah} là hướng phát triển Wi-Fi tiết kiệm năng lượng dùng dải tần dưới 1 GHz, hướng tới thiết bị IoT. Bảng minh họa của slide nêu kênh OFDM băng hẹp, chế độ tiết kiệm năng lượng, IP gốc, bảo mật Wi-Fi, phạm vi khoảng 1 km, khả năng xuyên tường và hoạt động với pin nhỏ trong thời gian dài theo điều kiện ứng dụng. \sourcepage{25--26}

\subsection{Bluetooth Low Energy (BLE)}
BLE là kết nối vô tuyến tầm ngắn cho mạng cá nhân (PAN), chú trọng thiết bị cảm biến và wearable có nhu cầu trao đổi gói dữ liệu nhỏ, không liên tục. Từ Bluetooth 4.0, các cơ chế năng lượng thấp thích hợp để một nút hỗ trợ kết nối các nút cục bộ và chuyển dữ liệu về backend. Ứng dụng: tòa nhà thông minh, giao thông, thiết bị đeo. \sourcepage{27}

\begin{table}[htbp]\centering\small
\caption{Bảng Classic Bluetooth và BLE ở slide PDF trang 28 (số liệu theo phiên bản trong slide).}
\begin{tabularx}{\textwidth}{p{3.4cm}XX}\toprule
Thông số & Classic Bluetooth & BLE\\\midrule
Tầm hoạt động & 100 m & Hơn 100 m\\
Tốc độ dữ liệu & 1--3 Mb/s & 1 Mb/s\\
Thông lượng ứng dụng & 0,7--2,1 Mb/s & 0,27 Mb/s\\
Tần số & 2,4 GHz & 2,4 GHz\\
Bảo mật & 56/128 bit (bảng nguồn) & AES 128 bit với CCM\\
Độ trễ & 100 ms & 6 ms\\
Thời gian thiết lập & 100 ms & 3 ms\\
Âm thanh & Có & Không (bảng theo đời cũ)\\
Mô hình mạng & Sao & Sao\\
Công suất tiêu thụ (theo slide) & 1 W & 0,01--0,5 W\\
Dòng đỉnh (theo slide) & Dưới 30 mA & Dưới 15 mA\\\bottomrule
\end{tabularx}\end{table}

\begin{nguonslide}[title={Lưu ý với bảng Bluetooth}]
Bảng trang 28 phản ánh cách so sánh trong thời điểm soạn slide, không đại diện cho mọi phiên bản Bluetooth mới hơn. Đặc biệt mục \emph{Voice capable: No} áp dụng theo bảng lịch sử và không mô tả đầy đủ LE Audio. Nên dùng datasheet thiết bị cụ thể để thiết kế công suất.
\end{nguonslide}

\subsection{Zigbee}
Zigbee là công nghệ phù hợp mạng cảm biến/điều khiển công suất thấp với yêu cầu an toàn, độ bền kết nối và nhiều nút; được sử dụng trong ứng dụng công nghiệp và tự động hóa. Zigbee 3.0 được slide mô tả là sự thống nhất các chuẩn Zigbee trước đó. \sourcepage{29}

Các tính chất giao thức được liệt kê: hỗ trợ điểm--điểm, điểm--đa điểm và mesh; chu kỳ hoạt động (duty cycle) thấp giúp kéo dài tuổi thọ pin; độ trễ thấp; trải phổ chuỗi trực tiếp (DSSS); khả năng mở rộng tới khoảng 65.000 nút/mạng theo slide; mã hóa AES 128 bit; tránh xung đột, phát lại và xác nhận. Các giới hạn triển khai thực tế phụ thuộc cấu hình và thiết bị. \sourcepage{30}

\begin{figure}[htbp]\centering
\begin{tikzpicture}[node distance=12mm, every node/.style={font=\small}, coord/.style={circle,draw=orangeiot,fill=orange!25,minimum size=8mm},router/.style={circle,draw=teal,fill=cyan!20,minimum size=8mm},end/.style={circle,draw=navy,fill=blue!12,minimum size=7mm}]
\node[coord](c){C}; \node[router,below left=14mm and 22mm of c](r1){R};\node[router,below right=14mm and 22mm of c](r2){R};\node[router,below=30mm of c](r3){R};
\node[end,below left=13mm and 7mm of r1](a){E}; \node[end,left=17mm of r1](b){E}; \node[end,right=17mm of r2](d){E};\node[end,below right=13mm and 7mm of r2](e){E};
\draw[thick] (c)--(r1)--(r3)--(r2)--(c);\draw[thick] (r1)--(r2);\draw (r1)--(a) (r1)--(b) (r2)--(d) (r2)--(e);
\end{tikzpicture}
\caption{Ví dụ mạng Zigbee mesh. C: coordinator; R: router; E: end device (vẽ lại từ sơ đồ các topo slide 31).}\end{figure}

\subsection{Z-Wave}
Z-Wave được ứng dụng chủ yếu trong nhà thông minh để gắn thiết bị điện, điều khiển từ Internet và hỗ trợ an ninh/tự động hóa. Theo slide: tần số 868,42 MHz tại châu Âu, 908,42 MHz tại Hoa Kỳ; các tốc độ 9,6; 40; 200 kbit/s; điều chế GFSK ở dải gần 900 MHz; tầm khoảng 30 m trong nhà và 100 m ngoài trời. Thông số phụ thuộc khu vực pháp lý và phiên bản thiết bị. \sourcepage{32}

\subsection{Weightless, Sigfox và LoRaWAN}
\paragraph{Weightless.} Họ tiêu chuẩn mở LPWAN kết nối trạm gốc với hàng nghìn nút. Có Weightless-W, Weightless-P và Weightless-N cho các điều kiện khác nhau. Weightless-W hỗ trợ hai chiều ở phổ trắng truyền hình (TV white spaces), slide lấy ví dụ lĩnh vực dầu khí; Weightless-P cũng hai chiều, thích hợp mạng riêng và các trường hợp cần cả đường lên/đường xuống. Slide chỉ nêu tên Weightless-N mà không khai triển tính chất của biến thể này. \sourcepage{33}

\paragraph{Sigfox.} Dịch vụ kết nối LPWAN theo thuê bao tại các quốc gia được triển khai, tối ưu gói dữ liệu nhỏ gửi không thường xuyên. Ví dụ tiêu biểu: thiết bị đo thông minh hoặc cảm biến phát cảnh báo. \sourcepage{34}

\paragraph{LoRaWAN.} Công nghệ mạng diện rộng công suất thấp do hệ sinh thái LoRa Alliance thúc đẩy, cho truyền hai chiều, hướng tới chi phí thấp, pin lâu, mạng lớn và dịch vụ di động. Ứng dụng: thành phố thông minh và công nghiệp. Cần phân biệt điều chế vô tuyến LoRa ở lớp vật lý với giao thức mạng LoRaWAN. \sourcepage{35}

\subsection{Li-Fi và truyền thông ánh sáng}
Li-Fi (Light Fidelity) dùng ánh sáng nhìn thấy (Visible Light Communication, VLC) để truyền dữ liệu tốc độ cao, có thể tận dụng đèn LED hiệu suất cao. Điểm hạn chế trong mô tả slide là yêu cầu điều kiện truyền quang phù hợp, dễ bị cản bởi vật chắn hoặc mất đường nhìn; chưa được triển khai rộng rãi trong IoT tại thời điểm biên soạn. Tiềm năng: lớp học và phòng trình chiếu thông minh. \sourcepage{36}

\subsection{Mạng di động và 5G}
Các thế hệ mạng tế bào (3G, 4G, LTE-Advanced, 5G) có vùng phủ rộng và khả năng hỗ trợ thiết bị di động. LTE-Advanced được slide lấy làm ví dụ về tiến bộ tốc độ và chế độ tiết kiệm năng lượng: có biến thể/chế độ cho các nút có thể đạt thời gian pin rất dài tùy chu kỳ gửi dữ liệu; bản LTE-A khác có tốc độ tối đa ở mức rất cao. Hai trị số 10 Mb/s và 1 Gb/s trong slide mô tả các trường hợp khác nhau, không phải một cấu hình chung. \sourcepage{37}

Slide 5G (có trích dẫn dự báo trước năm 2020) nêu tăng tốc độ, độ phủ, số người dùng và mức cơ động so với 4G. Bảng so sánh 4G/5G ở slide PDF trang 39 có các ý: dải tần 4G khoảng 600 MHz--5,925 GHz; 5G gồm FR1 (khoảng 450 MHz--6 GHz) và FR2 (khoảng 24,25--52,6 GHz, tùy quy định); băng thông 4G tới 20 MHz, 5G FR1 tới 100 MHz/FR2 tới 400 MHz; khoảng cách sóng mang 4G cố định 15 kHz, 5G mở rộng; tốc độ đỉnh downlink theo slide là 1 Gb/s và 20 Gb/s; độ trễ slide lần lượt 10 ms và dưới 1 ms; quy mô kết nối 2.000 so với 1 triệu thiết bị/km$^2$; tốc độ di động 350 so với 500 km/h. Đây là mục tiêu/giá trị minh họa trong bảng, không phải hiệu năng bảo đảm mọi vùng phủ. \sourcepage{38--39}

\subsection{LTE-M và NB-IoT}
\paragraph{LTE-M.} LTE-M, tên đầy đủ LTE Machine-Type Communication, gồm eMTC (enhanced MTC), là công nghệ kết nối di động LPWAN do 3GPP thiết kế cho thiết bị M2M/IoT. eMTC LTE Cat M1 đã được đặc tả ở 3GPP Release 13 (tháng 6 năm 2016 theo slide). Hệ 3GPP còn có NB-IoT và EC-GSM-IoT. \sourcepage{40}

\paragraph{NB-IoT.} Narrowband IoT là tiêu chuẩn LPWAN cho các thiết bị truyền dữ liệu ít, chú trọng mức tiêu thụ điện thấp, dung lượng mạng, hiệu quả phổ và vùng phủ sâu (ví dụ trong công trình). Slide cho rằng thời gian dùng pin hơn 10 năm khả thi ở một số kịch bản; điều này phụ thuộc dữ liệu gửi, vùng phủ, pin và cấu hình hoạt động. \sourcepage{41}

\begin{table}[htbp]\centering\small
\caption{So sánh khái quát hai công nghệ LTE-M và NB-IoT từ slide PDF trang 50.}
\begin{tabularx}{\textwidth}{p{3.4cm}XX}\toprule
Đặc tính & LTE-M & NB-IoT\\\midrule
Tên chuẩn & LTE Cat M1 (theo slide) & LTE Cat NB1\\
Băng thông & Khoảng 1,08 MHz, trong kênh LTE 1,4 MHz & Tới 180 kHz\\
Thông lượng & 36 kb/s (giá trị trong bảng slide) & 30/60 kb/s (giá trị trong bảng slide)\\
Triển khai & Trong các mạng LTE hiện hữu & In-band, guard-band hoặc độc lập trong một số phổ GSM\\
Khả năng giọng nói & Có hỗ trợ thoại/VoLTE theo cấu hình & Không được thiết kế cho thoại\\
Di động/chuyển giao & Có & Hạn chế\\
Thời lượng pin theo bảng & Có thể tới 10 năm & Có thể tới 10 năm\\\bottomrule
\end{tabularx}\end{table}

\subsection{Bảng tổng hợp phổ, tầm phủ và tốc độ các công nghệ}
\begingroup\small\sloppy\begin{longtable}{p{2.6cm}p{3.5cm}p{2.25cm}p{3.6cm}}
\caption{Bảng đối chiếu slide PDF trang 42; số liệu tham khảo theo slide, không đại diện mọi phiên bản.}\\
\toprule Công nghệ & Tần số theo slide & Tầm phủ & Tốc độ theo slide\\\midrule\endfirsthead
\toprule Công nghệ & Tần số theo slide & Tầm phủ & Tốc độ\\\midrule\endhead
Bluetooth & 2,4 GHz & 100 m & 125 kbit/s--Mb/s\\
Wi-Fi & 2,4 hoặc 5 GHz & 50 m & 150--600 Mb/s\\
NFC & 13,56 MHz & 4 cm & 100--420 kbit/s\\
LoRaWAN & 867--869 MHz (Âu); 902--928 MHz (Bắc Mỹ) & 15 km & 0,3--50 kbit/s\\
Cellular & 900/1800/\allowbreak 1900/2100 MHz & Bảng nguồn ghi 30 m (nút--trạm), không phải tầm phủ mạng tế bào nói chung & 21 Mb/s (3G+); 600 Mb/s (4G)\\
Z-Wave & 865--926 MHz (ISM) & 100 m & 100 kbit/s\\
Zigbee & 2,4 GHz (ISM) & 100 m & 20--250 kbit/s\\
Sigfox & 900 MHz & 3--50 km & 10--1000 bit/s\\\bottomrule
\end{longtable}\endgroup

\subsection{Bảo mật ở lớp truyền thông}
Các giao thức được slide nêu bao gồm:
\begin{itemize}
\item \textbf{IPsec (Internet Protocol Security):} cơ chế bảo mật ở lớp mạng IP, cung cấp các phương thức xác thực, toàn vẹn và mã hóa khi được triển khai/cấu hình đúng. Được sử dụng với IPv6 nhưng không nên hiểu IPsec mặc nhiên được bật trên mọi mạng IPv6.
\item \textbf{TLS (Transport Layer Security):} bảo vệ kênh kết nối dựa trên TCP, ví dụ HTTPS hoặc MQTT qua TLS.
\item \textbf{DTLS (Datagram Transport Layer Security):} giải pháp bảo mật cho truyền thông kiểu datagram như UDP, ví dụ CoAP/DTLS.
\item \textbf{IEEE 1888.3:} thành viên họ tiêu chuẩn dành cho mạng điều khiển cộng đồng xanh phổ cập (Ubiquitous Green Community Control Network). Mục đích theo slide là hạn chế truy cập trái phép, ngăn rò rỉ dữ liệu và bảo đảm tính bí mật/toàn vẹn dữ liệu truyền.
\end{itemize}
Các thuộc tính cốt lõi là bảo mật nội dung, xác thực đối tác và phát hiện sửa đổi; hệ thống IoT vẫn phải quản lý khóa, phân quyền nút và cập nhật phần mềm. \sourcepage{43}

\subsection{Định tuyến và giao thức RPL}
Định tuyến (routing) là quá trình chọn đường chuyển lưu lượng trong một mạng hoặc qua nhiều mạng. Với rất nhiều thiết bị có tài nguyên hạn chế, chạy pin và liên kết chất lượng không ổn định, lựa chọn đường tốt có ảnh hưởng lớn tới khả năng mở rộng, thời gian gửi dữ liệu và thời lượng pin.

\term{RPL}{Routing Protocol for Low-Power and Lossy Networks} là chuẩn định tuyến IPv6 dành cho mạng công suất thấp và liên kết dễ mất gói. Slide nêu có thể sử dụng trên IEEE 802.15.4, hỗ trợ các kiểu truyền một--một, nhiều--một và một--nhiều. Hình trang 45 biểu diễn hai cấu hình \emph{DODAG} (Destination-Oriented Directed Acyclic Graph): non-storing và storing mode; các nút tổ chức tương quan với nút gốc, còn phương thức duy trì đường đi phụ thuộc chế độ hoạt động. \sourcepage{44--45}

\begin{figure}[htbp]\centering
\begin{tikzpicture}[>=Latex, every node/.style={font=\small}, leaf/.style={circle,draw=orangeiot,fill=orange!18,minimum size=8mm}, root/.style={circle,draw=navy,fill=blue!20,minimum size=10mm}]
\node[root](r) at (0,3){Gốc};
\node[leaf](a) at (-2,1.7){A};\node[leaf](b) at (0,1.7){B};\node[leaf](c) at (2,1.7){C};
\node[leaf](d) at (-2.6,0){D};\node[leaf](e) at (-.9,0){E};\node[leaf](f) at (1,0){F};\node[leaf](g) at (2.7,0){G};
\foreach \x/\y in {a/r,b/r,c/r,d/a,e/a,e/b,f/b,f/c,g/c}{\draw[->,teal,thick](\x)--(\y);}
\end{tikzpicture}
\caption{Một đồ thị định tuyến hướng về nút gốc kiểu DODAG của RPL (minh họa vẽ lại từ slide trang 45).}\end{figure}

\section{Các loại mạng IoT}
Việc kết nối tin cậy giữa cảm biến, thiết bị và nền tảng IoT quyết định hoạt động của ứng dụng. Slide chia các mạng đang dùng thành bốn nhóm: mạng tế bào, mạng cục bộ/cá nhân (LAN/PAN), mạng diện rộng công suất thấp (LPWAN), và mạng mắt lưới (mesh). Đây là các góc nhìn phân loại theo phạm vi/vận hành và topo; một hệ thống thực tế có thể kết hợp chúng. \sourcepage{47}

\subsection{Mạng tế bào (Cellular Networks)}
Mạng tế bào sử dụng cơ sở hạ tầng nhà khai thác di động giống điện thoại thông minh để kết nối thiết bị IoT. Ưu điểm là vùng phủ địa lý rộng và hỗ trợ tính di động; hạn chế là điện năng và phí dịch vụ, cũng như khả năng phủ sóng kém ở hầm sâu, tủ kỹ thuật kín, mỏ hoặc công trình nhiều lớp vật cản. LTE-M và NB-IoT là hai lựa chọn IoT di động được đề cập; slide cho rằng NB-IoT có thể tiêu thụ ít năng lượng và giá thành thấp hơn LTE-M, nhưng phạm vi hỗ trợ dịch vụ phụ thuộc mạng triển khai ở thời điểm khảo sát. \sourcepage{48--50}

\subsection{Mạng cục bộ và cá nhân (LAN/PAN)}
PAN (Personal Area Network) dùng cho kết nối các thiết bị gần cá nhân; LAN (Local Area Network) kết nối các nút trong nhà, khu vực hoặc cơ sở. Phạm vi ngắn, chi phí sử dụng thường thấp; độ tin cậy phụ thuộc nhiễu, vật cản, mật độ nút và cấu hình. Hai ví dụ chính trong slide là Wi-Fi và Bluetooth. Có thể dùng gateway chuyển dữ liệu của mạng LAN/PAN ra mạng diện rộng. \sourcepage{51--52}

\subsection{Mạng diện rộng công suất thấp (LPWAN)}
LPWAN cho tầm xa hơn các kết nối Wi-Fi/Bluetooth thông dụng nhưng thường dùng ít năng lượng hơn nhiều công nghệ di động băng rộng. Đổi lại, tốc độ thấp và phù hợp các bản tin cảm biến nhỏ, gửi cách quãng. LoRaWAN, Sigfox, Weightless và các biến thể LPWAN tế bào như LTE-M/NB-IoT là các ví dụ đã giới thiệu trong phần trên. \sourcepage{53}

\subsection{Mạng mắt lưới (Mesh Networks)}
Trong mạng mesh, các nút có thể phối hợp chuyển tiếp dữ liệu cho nhau để đến gateway. Tầm mỗi bước nhảy có thể ngắn nhưng tổng vùng phủ được kéo dài qua nhiều nút trung gian. Điều này hữu ích cho cảm biến phân tán trong một tòa nhà; đổi lại cần đủ nút, bộ chuyển tiếp và cấu hình đường đi để tránh các vùng không phủ. Sơ đồ Zigbee và RPL ở các phần trước là những minh họa thích hợp cho kiến trúc đa chặng. \sourcepage{54}

\begin{figure}[htbp]\centering
\begin{tikzpicture}[>=Latex,scale=1,transform shape]
\draw[->,thick] (0,0)--(12,0) node[right]{Tầm phủ};
\draw[->,thick] (0,0)--(0,5) node[above,align=center]{Tốc độ dữ liệu};
\fill[blue!12] (.6,2.8) rectangle (3.3,4.6);\node[align=center] at (1.95,3.8){Wi-Fi\\Bluetooth};
\fill[cyan!16] (7,.65) rectangle (10.15,2.45);\node[align=center] at (8.55,1.55){LPWAN\\LoRa, Sigfox, NB-IoT};
\fill[orange!15] (4.0,2.5) rectangle (10.7,4.65);\node[align=center] at (7.45,3.65){Mạng tế bào\\2G/3G/4G/5G};
\end{tikzpicture}
\caption{So sánh định tính tầm phủ--tốc độ theo bố cục slide trang 55; vị trí vùng chỉ mang tính mô tả, không theo thang số.}\end{figure}

\begin{table}[htbp]\centering\small
\caption{Tóm tắt bốn nhóm mạng IoT.}
\begin{tabularx}{\textwidth}{p{2.6cm}XXX}\toprule
Loại mạng & Điểm mạnh & Hạn chế & Ví dụ\\\midrule
Cellular & Phủ rộng, di động & Năng lượng/chi phí; vùng phủ trong nhà sâu & 4G, LTE-M, NB-IoT, 5G\\
LAN/PAN & Chi phí triển khai thấp, tốc độ tùy công nghệ & Tầm ngắn, dễ bị cản/nhiễu & Wi-Fi, BLE\\
LPWAN & Tầm xa, phù hợp cảm biến chạy pin & Thông lượng thấp, hạn chế payload/chu kỳ & LoRaWAN, Sigfox, Weightless\\
Mesh & Mở rộng đa chặng, có tuyến thay thế & Cần nút chuyển tiếp, phức tạp định tuyến & Zigbee mesh, RPL\\\bottomrule
\end{tabularx}\end{table}

\subsection{Sơ đồ phân nhóm các công nghệ mạng}
Slide trang 46 sắp xếp mạng theo \emph{cellular}, \emph{wired}, \emph{satellite}, \emph{short-range wireless}, và \emph{LP-WAN / LP wireless}. Các mục được nêu gồm: 2G/2.5G, 3G, 4G/LTE, 5G; POTS, Fiber, Ethernet, ATM/Frame Relay, MPLS, Powerline; C-Band và L-Band; Wi-Fi/802.11, Bluetooth/BLE, Zigbee/802.15.4, 900 MHz proprietary, 6LoWPAN, Wi-Fi HaLow; WiMAX, Sigfox, LoRa, NB-IoT, LTE-M, Fixed Microwave và WiSUN. Hình này bổ sung các loại truyền dẫn chứ không thay đổi cách phân nhóm bốn loại của slide 47.

\section{Giao thức dữ liệu của IoT}
\subsection{Vai trò của giao thức ứng dụng và mô hình trao đổi}
Slide bắt đầu bằng việc phân biệt truyền thông dữ liệu giữa các thiết bị với khả năng truy cập Internet, nhấn mạnh các nút công suất thấp có thể giao tiếp trực tiếp bằng liên kết có dây hoặc vô tuyến mà không phải lúc nào cũng cần truy cập Internet. Trong ngăn xếp giao thức, Wi-Fi/BLE/LoRa thuộc nhóm công nghệ kết nối bên dưới, còn HTTP, CoAP, MQTT, WebSocket, AMQP và XMPP cung cấp cách thức trao đổi thông điệp ở tầng ứng dụng. \sourcepage{56,60}

Một nút IoT có thể thu thập dữ liệu, phân tích dữ liệu và điều khiển nút khác. Tùy hệ thống, trước khi tương tác với dịch vụ, nút cần được nhận dạng và đăng ký. Slide gọi quá trình đó là \emph{node subscription}. Trong ngữ cảnh truyền dữ liệu, cần phân biệt \emph{đăng ký một nút với dịch vụ} và \emph{đăng ký theo dõi một chủ đề dữ liệu (subscribe topic)}. \sourcepage{57}

\paragraph{Publish/subscribe (xuất bản/đăng ký).} Nút xuất bản (publisher) gửi bản tin lên một chủ đề/kênh; các nút đăng ký (subscriber) nhận dữ liệu tương ứng khi bản tin được xuất bản. Thành phần trung gian có thể là broker; cách tiếp cận tách người gửi và người nhận và thuận lợi cho nhiều nút trao đổi thường xuyên. \sourcepage{58--59}

\paragraph{Request/response (yêu cầu/phản hồi).} Client chủ động yêu cầu dữ liệu hoặc tác vụ và server trả kết quả. Mô hình quen thuộc với ứng dụng client/server và các dịch vụ REST. Slide nhận xét publish/subscribe thường thuận lợi hơn với triển khai IoT lớn trao đổi dữ liệu nhiều, trong khi request/response dễ tích hợp ứng dụng client/server. \sourcepage{59}

\begin{figure}[htbp]\centering
\begin{tikzpicture}[>=Latex,font=\small,block/.style={draw=navy,rounded corners,minimum width=29mm,minimum height=10mm,fill=blue!6}]
\node[block](cli) at (0,2.7){Client};\node[block](serv) at (5,2.7){Server};
\draw[->,thick,teal] (cli)--node[above]{Request}(serv);
\draw[<-,thick,orangeiot] (cli) to[bend right=22] node[below]{Response}(serv);
\node[block](pub) at (0,0){Publisher};\node[block](broker) at (5,0){Broker};\node[block](sub) at (10,0){Subscriber};
\draw[->,thick,teal](pub)--node[above]{Publish}(broker);
\draw[->,thick,teal](broker)--node[above]{Push chủ đề}(sub);
\end{tikzpicture}
\caption{Hai mô hình trao đổi dữ liệu ở slide trang 57--59.}\end{figure}

\subsection{HTTP: Hypertext Transfer Protocol}
HTTP là giao thức ứng dụng phổ biến của Web, thường theo mô hình \emph{pull}: client gửi yêu cầu lấy tài nguyên, server trả phản hồi. HTTP không lưu trạng thái yêu cầu trước đó ở bản thân giao thức (stateless), hỗ trợ thương lượng kiểu nội dung và biểu diễn tài nguyên để các hệ thống được phát triển độc lập vẫn trao đổi dữ liệu. \sourcepage{61--62}

Những phương thức HTTP phổ biến: \texttt{GET} (đọc tài nguyên), \texttt{POST} (tạo/xử lý dữ liệu), \texttt{PUT} (thay thế/cập nhật tài nguyên), \texttt{PATCH} (cập nhật một phần), \texttt{DELETE} (xóa). CRUD: Create--Read--Update--Delete là các thao tác ứng dụng tương ứng ở mức khái niệm, không phải ánh xạ duy nhất cho mọi API. \sourcepage{62}

Theo slide, HTTP có hạn chế đối với một số hệ thống IoT: quá nhiều cặp yêu cầu--phản hồi khi nhiều cảm biến hoạt động khiến phát sinh độ trễ và năng lượng; thông thường client khởi xướng yêu cầu thay vì server tự đẩy bản tin như mô hình publish/subscribe; HTTP/TCP có overhead và chi phí kết nối, không tối ưu cho một số cảm biến dùng pin. Những hạn chế này được mô tả cho cách triển khai HTTP phổ biến; các biến thể HTTP mới, keep-alive, server push hoặc kỹ thuật khác có thể thay đổi chi phí thực tế. \sourcepage{63--64}

\begin{figure}[htbp]\centering
\begin{tikzpicture}[>=Latex,font=\small]
\node at (0,4.8){Client};\node at (6,4.8){Server};
\draw[dashed](0,4.4)--(0,0);\draw[dashed](6,4.4)--(6,0);
\draw[->,thick,navy](0,4)--node[above]{TCP: SYN}(6,3.5);
\draw[<-,thick,navy](0,2.9)--node[above]{SYN/ACK}(6,3.4);
\draw[->,thick,navy](0,2.8)--node[above]{ACK}(6,2.4);
\draw[->,thick,teal](0,2)--node[above]{HTTP GET /temperature}(6,1.6);
\draw[<-,thick,orangeiot](0,.8)--node[above]{HTTP 200 OK + dữ liệu}(6,1.3);
\end{tikzpicture}
\caption{Chuỗi bắt tay TCP và yêu cầu--phản hồi HTTP, dựng theo ý tưởng hình slide trang 63.}\end{figure}

\subsection{TCP và UDP}
\textbf{TCP} là giao thức hướng kết nối, sử dụng bắt tay thiết lập phiên, đánh số thứ tự, xác nhận/điều khiển truyền lại để cung cấp luồng byte tin cậy. \textbf{UDP} truyền datagram không thiết lập phiên TCP và không tự đảm bảo xác nhận, thứ tự hoặc truyền lại. Do đó TCP thường có chi phí thiết lập/trạng thái cao hơn, còn UDP thích hợp với giao thức tự quản lý độ tin cậy hoặc yêu cầu thông điệp ngắn như CoAP. Slide trang 65 minh họa TCP bằng SYN--SYN/ACK--ACK và UDP bằng các datagram yêu cầu--phản hồi không có bắt tay. Việc nói 'UDP không tin cậy' chỉ có nghĩa UDP không tự bảo đảm chuyển phát, không có nghĩa ứng dụng UDP không thể cài đặt xác nhận. \sourcepage{65}

\subsection{CoAP: Constrained Application Protocol}
CoAP là giao thức truyền tài nguyên kiểu Web thiết kế cho nút và mạng bị hạn chế về băng thông, độ sẵn sàng và tài nguyên xử lý. Nó sử dụng mô hình client/server tương tự REST/HTTP: server công bố tài nguyên qua URI, client dùng GET, POST, PUT, DELETE. Bản tin CoAP được thiết kế gọn và thường dựa trên UDP; việc xác nhận chuyển phát phải do cơ chế CoAP xử lý khi cần. \sourcepage{66}

\begin{figure}[htbp]\centering
\begin{tikzpicture}[>=Latex,font=\small,box/.style={draw=navy,fill=blue!5,rounded corners,minimum width=23mm,minimum height=10mm}]
\node[box](web) at (0,0){HTTP client};\node[box](proxy) at (4.2,0){HTTP/CoAP proxy};\node[box](sensor) at (8.8,0){Cảm biến CoAP};
\draw[<->,thick,navy](web)--node[above]{HTTP / TCP}(proxy);
\draw[<->,thick,teal](proxy)--node[above]{CoAP / UDP}(sensor);
\end{tikzpicture}
\caption{Liên kết client Web với mạng cảm biến giới hạn qua proxy CoAP (theo slide 67).}\end{figure}

\paragraph{Đặc trưng CoAP.} Slide nêu năm nhóm ưu điểm: phục vụ giao tiếp M2M bị ràng buộc; trao đổi bản tin không đồng bộ; overhead thấp và dễ phân tích cú pháp; hỗ trợ URI và kiểu nội dung; hỗ trợ proxy và cache. CoAP có thể làm cầu nối với Web bằng gateway/proxy. \sourcepage{67--68}

\paragraph{Bốn loại bản tin.} CoAP định nghĩa:
\begin{enumerate}[leftmargin=2em]
\item \textbf{Confirmable (CON):} bản tin cần ACK; nếu không nhận được sẽ có cơ chế gửi lại.
\item \textbf{Non-confirmable (NON):} bản tin không đòi ACK ở lớp thông điệp CoAP (\emph{fire-and-forget}).
\item \textbf{Acknowledgement (ACK):} xác nhận đã nhận bản tin CON.
\item \textbf{Reset (RST):} từ chối/đặt lại bản tin hoặc xử lý quan sát không còn hợp lệ.
\end{enumerate}

CoAP dùng Message ID 16 bit cho ghép bản tin và phát hiện bản sao, cùng cơ chế truyền lại có back-off đối với CON. Bảo mật có thể sử dụng DTLS. \sourcepage{69}

\begin{figure}[htbp]\centering
\begin{tikzpicture}[>=Latex,font=\small]
\node at (0,4.8){Client};\node at (7,4.8){Server};
\draw[dashed](0,4.4)--(0,.1);\draw[dashed](7,4.4)--(7,.1);
\draw[->,thick,teal](0,3.9)--node[above]{CON, MID=0x4A51, GET}(7,3.35);
\draw[<-,thick,orangeiot](0,2.3)--node[above]{ACK, MID=0x4A51, 2.05 Content}(7,2.9);
\draw[->,thick,teal](0,1.7)--node[above]{NON, MID=0x4A52}(7,1.2);
\node[align=left] at (3.5,.45){ACK cho CON; không ACK cho NON.};
\end{tikzpicture}
\caption{Ví dụ luồng CoAP confirmable và non-confirmable, vẽ lại theo slide trang 70--71.}\end{figure}

\paragraph{Ví dụ tra cứu tài nguyên.} Các hình slide 70--71 cho thấy bản tin \texttt{CON GET /temperature} với \texttt{Token} được phản hồi bằng \texttt{ACK 2.05 Content} chứa nhiệt độ \texttt{20.1 C}, hoặc phản hồi lỗi \texttt{4.04 Not Found} khi tài nguyên không tồn tại. Hình trang 70 cũng thể hiện CON--ACK, CON--RST và NON một chiều, để làm rõ khác nhau giữa xác nhận và từ chối. \sourcepage{70--71}

\subsection{MQTT: Message Queuing Telemetry Transport}
MQTT là giao thức gửi thông điệp nhẹ, hữu ích để phân phối dữ liệu telemetry tới các client có tài nguyên và kết nối mạng hạn chế. Nó dùng mô hình publish/subscribe và rất phổ biến trong giao tiếp M2M/IoT. Thành phần trung tâm là \textbf{broker}: publisher gửi dữ liệu theo chủ đề (topic); broker xác định người đăng ký để chuyển tiếp; subscriber nhận bản tin không cần duy trì kết nối trực tiếp với từng publisher. \sourcepage{72--73}

\begin{figure}[htbp]\centering
\begin{tikzpicture}[>=Latex,font=\small,box/.style={draw=navy,rounded corners,fill=blue!5,minimum height=12mm,minimum width=28mm}]
\node[box](p) at (0,0){Publisher};\node[box,fill=cyan!15](b) at (4.5,0){MQTT Broker};
\node[box](s1) at (9,1.4){Subscriber 1};\node[box](s2) at (9,-1.4){Subscriber 2};
\draw[->,thick,teal](p)--node[above]{PUBLISH topic}(b);
\draw[->,thick,orangeiot](b)--node[above,sloped]{PUBLISH}(s1);
\draw[->,thick,orangeiot](b)--node[below,sloped]{PUBLISH}(s2);
\end{tikzpicture}
\caption{Kiến trúc MQTT publisher--broker--subscriber (dựng lại từ hình slide trang 73).}\end{figure}

\paragraph{Bốn giai đoạn phiên MQTT.} Theo slide: (1) \textbf{Connection}: client thiết lập kết nối TCP/IP đến broker (cổng chuẩn hoặc cổng riêng do server định nghĩa); (2) \textbf{Authentication}: xác thực nếu broker yêu cầu; (3) \textbf{Communication}: publish và subscribe bản tin; (4) \textbf{Termination}: đóng phiên/kết nối. Với Client ID và thiết lập phiên thích hợp, broker có thể tiếp tục một phiên trước đó. \sourcepage{74}

\paragraph{Ba mức Quality of Service (QoS).}
\begin{table}[htbp]\centering\small
\caption{Ba mức QoS trong MQTT (slide trang 75).}
\begin{tabularx}{\textwidth}{p{1.4cm}p{3cm}X}\toprule
QoS & Ý nghĩa & Đặc điểm\\\midrule
0 & At most once (tối đa một lần) & Không bảo đảm thông điệp đến đích; không có cơ chế xác nhận chuyển phát QoS 0\\
1 & At least once (ít nhất một lần) & Bảo đảm theo cơ chế xác nhận/gửi lại; có thể trùng lặp bản tin\\
2 & Exactly once (đúng một lần) & Có quy trình bắt tay để bảo đảm xử lý chuyển phát đúng một lần theo ngữ nghĩa MQTT\\\bottomrule
\end{tabularx}\end{table}

\paragraph{Ví dụ ứng dụng ở thời điểm giảng dạy.} Slide nêu Facebook sử dụng MQTT trong hệ thống nhắn tin để tiết kiệm pin và truyền tin nhanh trên mạng không ổn định; đồng thời nêu các nền tảng AWS, Google Cloud, IBM Bluemix, Microsoft Azure và ThingOn IoT Platform có hỗ trợ MQTT. Những khẳng định mang yếu tố thời điểm này được giữ theo tư liệu giảng dạy, không coi là xác nhận về sản phẩm hiện tại. \sourcepage{76}

\subsection{WebSocket}
WebSocket liên quan HTTP trong bước thiết lập kết nối: trình duyệt/client thực hiện cơ chế bắt tay nâng cấp kết nối và duy trì một kênh truyền trên TCP để cả hai phía trao đổi bản tin cho đến khi đóng. So với cơ chế nhiều lượt request/response, WebSocket cho phép truyền hai chiều trên kết nối lâu dài, thuận lợi cho ứng dụng thời gian thực; slide nhận xét nó vẫn có overhead đáng kể đối với các nút IoT rất hạn chế tài nguyên. Hình trang 78 minh họa cột HTTP với những cặp request--response lặp lại so với cột WebSocket với bước handshake rồi truyền hai chiều trực tiếp. \sourcepage{77--78}

\begin{figure}[htbp]\centering
\begin{tikzpicture}[>=Latex,font=\small]
\node at (0,3.5){Client};\node at (6,3.5){Server};\draw[dashed](0,3.2)--(0,0);\draw[dashed](6,3.2)--(6,0);
\draw[->,thick,navy](0,2.9)--node[above]{HTTP Upgrade: websocket}(6,2.65);
\draw[<-,thick,navy](0,2.22)--node[above]{101 Switching Protocols}(6,2.5);
\draw[->,thick,teal](0,1.6)--node[above]{Thông điệp client $\to$ server}(6,1.4);
\draw[<-,thick,orangeiot](0,.8)--node[above]{Thông điệp server $\to$ client}(6,1.0);
\end{tikzpicture}
\caption{Thiết lập và truyền hai chiều với WebSocket, vẽ lại theo slide 77--78.}\end{figure}

\subsection{AMQP: Advanced Message Queuing Protocol}
AMQP là tiêu chuẩn mở thiên về giao tiếp hướng thông điệp, hỗ trợ publish/subscribe. Tương tự cách mô tả MQTT trong slide, các ứng dụng có thể sử dụng các ngữ nghĩa chuyển phát tối đa một lần, ít nhất một lần hoặc đúng một lần tùy giao thức/triển khai. Kênh TCP có thể được bảo vệ với TLS/SSL, đồng thời có cơ chế xác thực SASL. Hai thành phần cốt lõi được slide nêu:
\begin{itemize}
\item \textbf{Exchange:} nhận bản tin và chọn (định tuyến) queue phù hợp dựa trên các quy tắc đã đặt.
\item \textbf{Message queue:} lưu bản tin chờ được chuyển tới phía tiêu thụ/nhận.
\end{itemize}
Kiến trúc này phù hợp hệ thống có xử lý hàng đợi, định tuyến thông điệp linh hoạt và tích hợp backend. \sourcepage{79}

\begin{figure}[htbp]\centering
\begin{tikzpicture}[>=Latex,font=\small,box/.style={draw=navy,rounded corners,minimum width=26mm,minimum height=10mm,fill=blue!5}]
\node[box](p) at (0,0){Producer};\node[box](ex) at (4,0){Exchange};\node[box](q) at (8,0){Queue};\node[box](c) at (12,0){Consumer};
\draw[->,thick,teal](p)--(ex);\draw[->,thick,teal](ex)--node[above]{Routing}(q);\draw[->,thick,teal](q)--(c);
\end{tikzpicture}
\caption{Luồng thông điệp cơ bản AMQP theo slide trang 79.}\end{figure}

\subsection{XMPP: Extensible Messaging and Presence Protocol}
XMPP là giao thức tiêu chuẩn hóa qua IETF để trao đổi thông tin trên Internet, hoạt động gần thời gian thực, độ trễ thấp theo bối cảnh ứng dụng và có tính độc lập nền tảng/phân tán. Theo slide, XMPP dựa trên TCP, hỗ trợ cả publish/subscribe và request/response; client kết nối tới server và trao đổi bản tin dựa trên \term{XML}{Extensible Markup Language}. Nhờ cấu trúc thông điệp mở rộng, nó có thể phục vụ ứng dụng liên lạc IoT, nhưng XML thường tạo thêm overhead so với các định dạng nhị phân nhỏ gọn. \sourcepage{80}

\subsection{Bảng so sánh các giao thức dữ liệu}
\begingroup\footnotesize\sloppy\begin{longtable}{p{1.8cm}p{3.45cm}p{1.8cm}p{4.15cm}}
\caption{Tổng hợp từ các slide trang 56--80; một số nhận xét là phân tích hệ thống hóa.}\\
\toprule Giao thức & Mô hình & Vận chuyển & Ghi chú trong IoT\\\midrule\endfirsthead
\toprule Giao thức & Mô hình & Vận chuyển & Ghi chú trong IoT\\\midrule\endhead
HTTP & Request/response, client/server & TCP & Phổ biến REST, stateless; overhead kết nối/yêu cầu có thể cao\\
CoAP & REST, request/response; có hỗ trợ quan sát & UDP & Nhẹ, CON/NON/ACK/RST, URI, proxy/cache, DTLS\\
MQTT & Publish/subscribe qua broker & TCP & Chủ đề, phiên, QoS 0/1/2\\
WebSocket & Kết nối song công dài hạn & TCP & Nâng cấp từ HTTP; trao đổi hai chiều liên tục\\
AMQP & Trao đổi thông điệp, hàng đợi & TCP & Exchange, queue, routing, TLS/SASL\\
XMPP & Cả publish/subscribe và request/response & TCP & Bản tin XML, mở rộng, gần thời gian thực\\\bottomrule
\end{longtable}\endgroup

\section{Tổng kết kiến thức và gợi ý ứng dụng}
\begin{ghinho}[title={Những khái niệm cần phân biệt}]
\textbf{Định danh/địa chỉ:} URI, EPC, uCode, IPv4, IPv6.\par
\textbf{Kết nối vật lý và mạng:} PLC/X10, NFC, UWB, Wi-Fi/HaLow, BLE, Zigbee, Z-Wave, Weightless, Sigfox, LoRaWAN, Li-Fi, 3G/4G/5G, LTE-M, NB-IoT.\par
\textbf{Bảo vệ, định tuyến:} IPsec, TLS/DTLS, IEEE 1888.3, RPL.\par
\textbf{Kiến trúc mạng:} cellular, LAN/PAN, LPWAN, mesh.\par
\textbf{Cách trao đổi dữ liệu:} request/response và publish/subscribe; HTTP, CoAP, MQTT, WebSocket, AMQP, XMPP.
\end{ghinho}

Một ví dụ tổng hợp: cảm biến chạy pin trong nhà có thể kết nối BLE hoặc Zigbee tới gateway; gateway kết nối Wi-Fi hoặc Ethernet lên máy chủ; dữ liệu được xuất bản qua MQTT để nhiều dịch vụ nhận cùng thời điểm. Một cảm biến cần REST tối giản có thể dùng CoAP/UDP; hệ thống Web có thể dùng HTTP để truy vấn và WebSocket để cập nhật liên tục. Nút ở vị trí xa không có Wi-Fi có thể cân nhắc LoRaWAN hoặc LTE-M/NB-IoT. Các lựa chọn này là \emph{minh họa thiết kế được suy ra từ bài học}, không phải kịch bản thí nghiệm bắt buộc trong slide.

\section{Đối chiếu toàn bộ 81 trang slide}
Bảng dưới đây xác định trang tài liệu gốc liên quan tới từng mục, để người đọc rà soát các hình và bảng của bài giảng. Trang 1 là bìa; trang 2 là mục lục; trang 81 là lời cảm ơn. Các trang mang hình không trích xuất được thành chữ được mô tả và minh họa lại trong phần tương ứng.
\begin{longtable}{p{3.0cm}p{10cm}}
\toprule Trang PDF & Nội dung có trong tài liệu/chương\\\midrule\endhead
1--2 & Bìa, học phần, giảng viên, mục lục\\
3--8 & Định nghĩa giao tiếp, H2M, M2M, M2H, H2H và các ví dụ\\
9--10 & Tổng quan giao thức IoT; định danh nút mạng\\
11--15 & URI/EPC/uCode; IPv4 và IPv6; hai hình so sánh địa chỉ và lợi ích IPv6\\
16--18 & Các lợi ích của IPv6; biểu tượng các chuẩn kết nối IoT\\
19--21 & PLC và X10; bảng thông số X10\\
22--23 & NFC; bảng tần số, tầm xa, chế độ và mã hóa\\
24--26 & UWB; Wi-Fi, Wi-Fi Direct và bảng tính năng Wi-Fi HaLow\\
27--28 & BLE và bảng so sánh Bluetooth Classic/BLE\\
29--31 & Zigbee, tính chất giao thức, hình sao--cây--mesh\\
32--35 & Z-Wave, Weightless, Sigfox, LoRaWAN\\
36--39 & Li-Fi, mạng tế bào, 5G, bảng 4G so với 5G\\
40--42 & LTE-M, NB-IoT, bảng so sánh đa công nghệ vô tuyến\\
43--46 & IPsec/TLS/DTLS/IEEE 1888.3, RPL và hai sơ đồ phân loại/định tuyến\\
47--55 & Bốn loại mạng IoT, ví dụ cellular/LAN-PAN/LPWAN/mesh, biểu đồ tầm xa--thông lượng\\
56--60 & Giao thức dữ liệu, đăng ký nút, publish/subscribe, request/response, liệt kê giao thức\\
61--65 & HTTP, phương thức HTTP, hạn chế IoT, sơ đồ HTTP và TCP so với UDP\\
66--71 & CoAP, proxy, REST, CON/NON/ACK/RST, luồng gói và ví dụ GET nhiệt độ\\
72--76 & MQTT, kiến trúc broker, bốn giai đoạn phiên, QoS 0--2, ví dụ áp dụng\\
77--78 & WebSocket, tương tác hai chiều và so sánh hình với HTTP\\
79--80 & AMQP (exchanges/queues) và XMPP/XML\\
81 & Kết thúc và lời cảm ơn\\\bottomrule
\end{longtable}

\begin{nguonslide}[title={Nguồn và phương pháp biên soạn}]
Tài liệu chính: TS. Nguyễn Đắc Cử, \textit{Cơ sở và thiết kế hệ thống IoT -- Chapter 3: Networking technologies in IoT}, Khoa Điện -- Điện tử, Trường Đại học Phenikaa, bản PDF 81 trang. Cách trình bày ở đây là \textbf{bản tổng hợp tiếng Việt theo chủ đề} chứ không phải bản dịch đối chiếu từng trang. Các sơ đồ đã được vẽ lại, dữ liệu theo bảng slide được ghi chú riêng; những điểm không nhất quán trong slide (như đơn vị gói X10, thiết lập NFC và phạm vi cellular trong bảng so sánh) được chỉ rõ chứ không tự ý thay thế. Khi tệp PDF gốc có mặt tại \texttt{source/chapter3-slides.pdf}, có thể đối chiếu trực tiếp với bảng trang ở trên.
\end{nguonslide}
\chapter{Công nghệ mạng trong Internet vạn vật}
\label{chap:iot-networking}
\begin{ghinho}[title={Mục tiêu và phạm vi}]
Chương này tổng hợp bốn phần của 81 trang slide: (1) giao tiếp trong IoT; (2) công nghệ và giao thức truyền thông, gồm định danh, địa chỉ hóa, công nghệ truyền dẫn, bảo mật và định tuyến; (3) phân loại mạng IoT; (4) giao thức dữ liệu HTTP, CoAP, MQTT, WebSocket, AMQP và XMPP. Các bảng thông số được ghi theo chính slide, không được hiểu là mức cam kết cho mọi phiên bản, thiết bị hoặc quốc gia. 
\end{ghinho}

\section{Giao tiếp trong hệ thống IoT}
\subsection{Khái niệm và vai trò của truyền thông}
Internet vạn vật (Internet of Things, IoT) kết nối các thiết bị thông minh qua mạng, cho phép chúng trao đổi dữ liệu và thực hiện công việc với ít hoặc không cần sự can thiệp trực tiếp của con người. Những thiết bị này thường tích hợp linh kiện điện tử, cảm biến, phần mềm, kết nối mạng và cơ cấu chấp hành. Truyền thông là điều kiện để dữ liệu đo ở một thiết bị trở thành thông tin có ích cho các thiết bị, dịch vụ và người sử dụng ở nơi khác. Không có liên kết đáng tin cậy giữa các thiết bị và nền tảng xử lý, một hệ thống cảm biến rời rạc chưa phát huy đầy đủ giá trị IoT. \sourcepage{3}

\subsection{Bốn hướng tương tác trong IoT}
\paragraph{Người đến máy (Human-to-Machine, H2M).} Con người gửi đầu vào cho hệ thống dưới dạng giọng nói, văn bản, hình ảnh hoặc thao tác trực tiếp; thiết bị thu nhận, phân tích và đáp ứng bằng văn bản, hình ảnh hoặc hành động. H2M phối hợp phần cứng và phần mềm để hỗ trợ các công việc hằng ngày. Ví dụ: nhận dạng giọng nói; nhận dạng khuôn mặt; chấm công sinh trắc học. \sourcepage{4}

\paragraph{Máy đến máy (Machine-to-Machine, M2M).} Hai hoặc nhiều máy trao đổi dữ liệu hoặc lệnh tự động, không đòi hỏi con người tham gia vào mỗi phiên liên lạc. Đường truyền có thể là hữu tuyến hoặc vô tuyến; một kết nối M2M điểm--điểm có thể sử dụng Ethernet hoặc mạng di động. IoT tiếp thu cách vận hành tự động của M2M và mở rộng sang các mạng thiết bị quy mô lớn có dịch vụ/nền tảng đám mây. Ví dụ: máy bán hàng tự động thông báo hết hàng cho nhà phân phối; máy giặt thông minh gửi thông báo hoàn tất; công tơ thông minh theo dõi điện năng và cảnh báo chủ sử dụng. \sourcepage{5--6}

\paragraph{Máy đến người (Machine-to-Human, M2H).} Thiết bị chủ động phát tín hiệu, văn bản, hình ảnh hoặc lời nói hướng tới con người, kể cả khi người đó không trực tiếp phát yêu cầu. Ví dụ: chuông báo cháy, đèn giao thông, thiết bị giám sát sức khỏe. \sourcepage{7}

\paragraph{Người đến người (Human-to-Human, H2H).} Con người liên lạc bằng lời nói, chữ viết, hình vẽ, nét mặt hoặc ngôn ngữ cơ thể. Trong một hệ thống IoT thực tế, H2H vẫn cần thiết để phối hợp xử lý sự cố, quản lý tình huống và đưa ra quyết định khi tự động hóa M2M chưa đáp ứng đủ. \sourcepage{8}

\begin{figure}[htbp]\centering
\begin{tikzpicture}[>=Latex,box/.style={draw=navy,rounded corners,fill=blue!5,minimum width=36mm,minimum height=13mm,align=center},node distance=20mm]
\node[box](human){Con người};\node[box,right=36mm of human](machine){Thiết bị IoT};
\draw[->,thick,teal,bend left=20](human) to node[above]{H2M: lệnh/đầu vào}(machine);
\draw[->,thick,orangeiot,bend left=20](machine) to node[below]{M2H: cảnh báo/kết quả}(human);
\node[box,below=19mm of machine](node2){Thiết bị IoT khác};
\draw[<->,thick,navy](machine)--node[right]{M2M}(node2);
\node[box,below=19mm of human](human2){Con người khác};
\draw[<->,thick,navy](human)--node[left]{H2H}(human2);
\end{tikzpicture}
\caption{Minh họa bốn loại hình tương tác H2M, M2M, M2H, H2H (vẽ lại theo các slide 4--8).}\end{figure}

\section{Công nghệ và giao thức truyền thông IoT}
\subsection{Tổng quan và yêu cầu lựa chọn}
Một hệ thống IoT có thể dùng Bluetooth, Wi-Fi, Wi-Fi Direct, giao thức vô tuyến công suất thấp, mạng di động LTE/LTE-A hoặc các công nghệ khác. Không có giao thức tốt nhất cho mọi ứng dụng: phải cân nhắc tầm phủ sóng, mức tiêu thụ năng lượng, băng thông, độ trễ, chi phí, tính di động, quy mô nút và mức bảo mật. Slide giới thiệu các họ công nghệ: PLC, X10, NFC, UWB, Wi-Fi, BLE, Zigbee, Z-Wave, Weightless, Sigfox, LoRaWAN, Li-Fi, mạng tế bào/5G, LTE-M và NB-IoT. \sourcepage{2,9,18}

\subsection{Định danh, đặt tên và địa chỉ hóa}
\paragraph{Định danh (identification).} Khi một nút gia nhập mạng, nó cần được nhận biết bằng một định danh. Cơ chế định danh giúp quản lý thiết bị, xác định đối tượng giao tiếp và hạn chế xử lý lưu lượng không cần thiết trong mạng rất lớn. Hệ thống còn cần tên và địa chỉ để tìm đúng thiết bị, đúng tài nguyên. Lưu ý về mặt kỹ thuật: việc cấp định danh tự nó không chặn được tấn công hay lưu lượng tràn ngập; phải kết hợp xác thực, phân quyền và cơ chế điều khiển lưu lượng. \sourcepage{10}

\paragraph{Đặt tên tài nguyên và đối tượng.} Ba cơ chế được nêu trong slide gồm:
\begin{itemize}
\item \textbf{URI (Uniform Resource Identifier):} chuỗi nhận dạng tài nguyên trên mạng. Ví dụ trong slide: \texttt{http://www.example.com/temperature/}.
\item \textbf{EPC (Electronic Product Code):} mã định danh đối tượng vật lý, thường mã hóa trên thẻ RFID thụ động. Ví dụ: \texttt{urn:epc:id:sgtin:}\allowbreak\texttt{0614141.012345.62852}.
\item \textbf{uCode (Ubiquitous Code):} mã định danh của hệ thống định danh đối tượng; ví dụ slide: \texttt{00001C0000000000}\allowbreak\texttt{0001000285E7A6E3}. Trang giới thiệu: \url{https://www.uidcenter.org}.
\end{itemize}
\sourcepage{11}

\paragraph{IPv4.} Giao thức Internet phiên bản 4 dùng địa chỉ 32 bit, biểu diễn thường gặp bằng bốn nhóm số thập phân phân tách bằng dấu chấm; tổng không gian lý thuyết là $2^{32}$ (xấp xỉ 4,3 tỷ) địa chỉ. Địa chỉ có thể chia mạng con, định tuyến và được cấu hình thủ công hoặc qua DHCP. Ví dụ slide minh họa dạng \texttt{192.168.5.18}. \sourcepage{12,14}

\paragraph{IPv6.} Giao thức Internet phiên bản 6 dùng địa chỉ 128 bit, tương ứng $2^{128}\approx 3{,}4\times 10^{38}$ địa chỉ; ví dụ ký hiệu rút gọn sử dụng các nhóm số hệ 16 và dấu hai chấm. Mục tiêu thiết kế bao gồm giải quyết hạn chế không gian địa chỉ của IPv4 và hỗ trợ số lượng rất lớn thiết bị. IPv6 có cơ chế tự cấu hình địa chỉ, hỗ trợ multicast và khả năng phối hợp với các cơ chế bảo vệ như IPsec. \sourcepage{13--17}

\begin{table}[htbp]\centering\small
\caption{So sánh IPv4 và IPv6 theo minh họa các slide 12--17.}
\begin{tabularx}{\textwidth}{p{3.15cm}XX}\toprule
Tiêu chí & IPv4 & IPv6\\\midrule
Độ dài địa chỉ & 32 bit & 128 bit\\
Số địa chỉ lý thuyết & $2^{32}\approx4{,}3\cdot10^9$ & $2^{128}\approx3{,}4\cdot10^{38}$\\
Dạng địa chỉ & 4 nhóm số thập phân (ví dụ 192.168.5.18) & Các nhóm số hệ 16, phân cách bằng dấu hai chấm\\
Cấu hình & Thủ công hoặc DHCP & Hỗ trợ tự cấu hình; có thể dùng DHCPv6\\
Ứng dụng trong IoT & Có thể bị hạn chế bởi số địa chỉ công cộng & Tạo không gian địa chỉ rất lớn cho nhiều nút\\\bottomrule
\end{tabularx}\end{table}

\paragraph{Các lợi ích IPv6 mà slide liệt kê.} (1) Dễ nhận diện và trao đổi dữ liệu giữa nhiều thiết bị hơn nhờ không gian địa chỉ lớn; (2) tự cấu hình đơn giản hơn ở quy mô lớn, có thể giảm công cấu hình và chi phí triển khai; (3) hỗ trợ truyền tải dữ liệu và hiệu quả sử dụng băng thông; (4) hỗ trợ các thuộc tính bảo mật về bí mật, xác thực và toàn vẹn khi cấu hình các cơ chế phù hợp; (5) multicast giúp tránh một số thao tác phát quảng bá thường xuyên, giảm số gói phải xử lý và có thể tiết kiệm pin; (6) cơ chế định tuyến và cấu trúc địa chỉ linh hoạt theo cách tổ chức mạng/ứng dụng. Đây là các lợi ích tiềm năng theo slide, không phải bảo đảm rằng bật IPv6 tự động tăng tốc hay bảo mật toàn bộ hệ thống. \sourcepage{16--17}

\begin{ghinho}[title={Phân biệt ba khái niệm}]
\textbf{Tên} cho biết tài nguyên hoặc đối tượng nào cần tìm; \textbf{định danh} phân biệt đối tượng trong hệ thống; \textbf{địa chỉ mạng} giúp chuyển gói đến vị trí mạng đang sử dụng. Tên, mã RFID và địa chỉ IP vì thế có thể cùng tồn tại trên một thiết bị IoT.
\end{ghinho}

\subsection{Công nghệ truyền dẫn qua dây điện: PLC và X10}
\paragraph{PLC (Power-Line Communication).} PLC sử dụng ngay đường dây điện hiện có để mang cả điện năng xoay chiều và tín hiệu dữ liệu. Một nút có thể nhận nguồn điện và trao đổi thông tin mà không cần tuyến dây dữ liệu riêng. Phù hợp cho thiết bị cố định, nhà thông minh, lưới điện thông minh, phương tiện hoặc môi trường có nhiễu vô tuyến cao. Hạn chế thực tế phụ thuộc chất lượng dây, tải và nhiễu từ thiết bị điện. \sourcepage{19}

\paragraph{X10.} Chuẩn điều khiển thiết bị trên đường điện ra đời khoảng thập niên 1970. Tín hiệu điều khiển được đưa lên dây điện bằng các xung sóng mang RF ngắn. Tốc độ và phạm vi thấp so với một số hệ PLC khác nhưng chi phí thấp; ứng dụng chính là bật/tắt, điều chỉnh đèn và điều khiển ổ cắm trong nhà thông minh qua bộ điều khiển hoặc ứng dụng. \sourcepage{20}

\begin{longtable}{p{.38\textwidth}p{.55\textwidth}}
\caption{Các thông số X10 theo bảng ở slide PDF trang 21.}\\\toprule
Thông số & Giá trị trong slide\\\midrule\endfirsthead
\toprule Thông số & Giá trị trong slide\\\midrule\endhead
Môi trường truyền & Đường dây điện (powerline only)\\
Loại thiết bị & Bộ điều khiển; thiết bị bị điều khiển; thiết bị có báo trạng thái\\
Điều chế lớp vật lý & OOK, sóng mang 120 kHz\\
Tốc độ dữ liệu & 120 bps tức thời; 60 bps duy trì; 24 bps tải hữu ích\\
Địa chỉ hóa & 16 mã nhà A--P và 16 mã thiết bị 1--16, tối đa 256 địa chỉ\\
Lệnh & 16 lệnh, như ON, OFF, DIM, BRIGHT, ALL UNITS OFF, ALL LIGHTS ON\\
Độ dài thông điệp & Slide ghi 100 bit (8,5 byte); hai trị số không tương đương nên cần kiểm tra nếu dùng thiết kế bản tin\\
Mã hóa & Manchester\\
Cấu trúc gói & Start Code 4 bit; House Code 8 bit; Key Code 10 bit (theo slide)\\
ACK và bảo mật & Không hỗ trợ theo bảng slide\\\bottomrule
\end{longtable}

\subsection{Giao tiếp khoảng cách rất ngắn: NFC và UWB}
\paragraph{NFC (Near-Field Communication).} Tập hợp giao thức vô tuyến cho các thiết bị tiếp xúc/đặt gần nhau, hướng tới xác thực, đọc/ghi thẻ, chia sẻ dữ liệu và kiểm soát ra vào; ví dụ thanh toán thông minh. Khác với cách dùng RFID chỉ để định danh trong nhiều hệ thống, NFC có các chế độ giao tiếp cho cả thiết bị và thẻ. \sourcepage{22}

\begin{longtable}{p{.47\textwidth}p{.45\textwidth}}
\caption{Thông số và chế độ NFC theo slide PDF trang 23.}\\\toprule
Đặc tính & Giá trị được nêu\\\midrule\endfirsthead
\toprule Đặc tính & Giá trị được nêu\\\midrule\endhead
Tần số mang & 13,56 MHz\\
Khoảng cách & Dưới 10 cm (bảng đối chiếu trang 42 sử dụng 4 cm)\\
Tốc độ & 106, 212 hoặc 424 kbit/s\\
Đối tượng & Thẻ (tag) và đầu đọc (reader)\\
Loại thẻ & Type 1 đến Type 5\\
Cấu hình mạng & Peer-to-peer\\
Chế độ thiết bị & Card emulation; reader/writer; peer-to-peer (active/passive)\\
Cách giao tiếp & Active--passive hoặc active--active\\
Thiết lập kết nối & Bảng slide ghi 'vài giây' và 'xấp xỉ dưới 0,1 giây'; cần xác định phép đo cụ thể\\
Mã hóa đường truyền & NRZ-L, Manchester, Modified Miller\\
Điều chế & ASK, BPSK (theo bảng)\\
Giải quyết va chạm & Có cơ chế chống va chạm\\
Tiêu chuẩn & ISO 14443A/B, ISO 18092, JIS X6319-4\\\bottomrule
\end{longtable}

\paragraph{UWB (Ultra-Wideband).} Vô tuyến sử dụng phổ tần rất rộng và xung phát năng lượng thấp để truyền dữ liệu và/hoặc định vị. Theo slide, khả năng hạn chế can nhiễu và theo dõi vị trí từng nút hỗ trợ mở rộng hệ thống; ứng dụng nổi bật là dịch vụ định vị và theo dõi tài sản. Không đồng nhất UWB với định vị GPS hay với mạng LPWAN: công năng và môi trường triển khai khác nhau. \sourcepage{24}

\subsection{Wi-Fi và Wi-Fi HaLow}
Wi-Fi phổ biến trong kết nối mạng cục bộ và có thể thiết lập kết nối ad hoc như Wi-Fi Direct mà không cần access point. Điểm bất lợi với cảm biến chạy pin là công suất tiêu thụ, dù nó không quan trọng bằng đối với nhiều thiết bị nhà thông minh có nguồn ổn định. \term{Wi-Fi HaLow}{IEEE 802.11ah} là hướng phát triển Wi-Fi tiết kiệm năng lượng dùng dải tần dưới 1 GHz, hướng tới thiết bị IoT. Bảng minh họa của slide nêu kênh OFDM băng hẹp, chế độ tiết kiệm năng lượng, IP gốc, bảo mật Wi-Fi, phạm vi khoảng 1 km, khả năng xuyên tường và hoạt động với pin nhỏ trong thời gian dài theo điều kiện ứng dụng. \sourcepage{25--26}

\subsection{Bluetooth Low Energy (BLE)}
BLE là kết nối vô tuyến tầm ngắn cho mạng cá nhân (PAN), chú trọng thiết bị cảm biến và wearable có nhu cầu trao đổi gói dữ liệu nhỏ, không liên tục. Từ Bluetooth 4.0, các cơ chế năng lượng thấp thích hợp để một nút hỗ trợ kết nối các nút cục bộ và chuyển dữ liệu về backend. Ứng dụng: tòa nhà thông minh, giao thông, thiết bị đeo. \sourcepage{27}

\begin{table}[htbp]\centering\small
\caption{Bảng Classic Bluetooth và BLE ở slide PDF trang 28 (số liệu theo phiên bản trong slide).}
\begin{tabularx}{\textwidth}{p{3.4cm}XX}\toprule
Thông số & Classic Bluetooth & BLE\\\midrule
Tầm hoạt động & 100 m & Hơn 100 m\\
Tốc độ dữ liệu & 1--3 Mb/s & 1 Mb/s\\
Thông lượng ứng dụng & 0,7--2,1 Mb/s & 0,27 Mb/s\\
Tần số & 2,4 GHz & 2,4 GHz\\
Bảo mật & 56/128 bit (bảng nguồn) & AES 128 bit với CCM\\
Độ trễ & 100 ms & 6 ms\\
Thời gian thiết lập & 100 ms & 3 ms\\
Âm thanh & Có & Không (bảng theo đời cũ)\\
Mô hình mạng & Sao & Sao\\
Công suất tiêu thụ (theo slide) & 1 W & 0,01--0,5 W\\
Dòng đỉnh (theo slide) & Dưới 30 mA & Dưới 15 mA\\\bottomrule
\end{tabularx}\end{table}

\begin{nguonslide}[title={Lưu ý với bảng Bluetooth}]
Bảng trang 28 phản ánh cách so sánh trong thời điểm soạn slide, không đại diện cho mọi phiên bản Bluetooth mới hơn. Đặc biệt mục \emph{Voice capable: No} áp dụng theo bảng lịch sử và không mô tả đầy đủ LE Audio. Nên dùng datasheet thiết bị cụ thể để thiết kế công suất.
\end{nguonslide}

\subsection{Zigbee}
Zigbee là công nghệ phù hợp mạng cảm biến/điều khiển công suất thấp với yêu cầu an toàn, độ bền kết nối và nhiều nút; được sử dụng trong ứng dụng công nghiệp và tự động hóa. Zigbee 3.0 được slide mô tả là sự thống nhất các chuẩn Zigbee trước đó. \sourcepage{29}

Các tính chất giao thức được liệt kê: hỗ trợ điểm--điểm, điểm--đa điểm và mesh; chu kỳ hoạt động (duty cycle) thấp giúp kéo dài tuổi thọ pin; độ trễ thấp; trải phổ chuỗi trực tiếp (DSSS); khả năng mở rộng tới khoảng 65.000 nút/mạng theo slide; mã hóa AES 128 bit; tránh xung đột, phát lại và xác nhận. Các giới hạn triển khai thực tế phụ thuộc cấu hình và thiết bị. \sourcepage{30}

\begin{figure}[htbp]\centering
\begin{tikzpicture}[node distance=12mm, every node/.style={font=\small}, coord/.style={circle,draw=orangeiot,fill=orange!25,minimum size=8mm},router/.style={circle,draw=teal,fill=cyan!20,minimum size=8mm},end/.style={circle,draw=navy,fill=blue!12,minimum size=7mm}]
\node[coord](c){C}; \node[router,below left=14mm and 22mm of c](r1){R};\node[router,below right=14mm and 22mm of c](r2){R};\node[router,below=30mm of c](r3){R};
\node[end,below left=13mm and 7mm of r1](a){E}; \node[end,left=17mm of r1](b){E}; \node[end,right=17mm of r2](d){E};\node[end,below right=13mm and 7mm of r2](e){E};
\draw[thick] (c)--(r1)--(r3)--(r2)--(c);\draw[thick] (r1)--(r2);\draw (r1)--(a) (r1)--(b) (r2)--(d) (r2)--(e);
\end{tikzpicture}
\caption{Ví dụ mạng Zigbee mesh. C: coordinator; R: router; E: end device (vẽ lại từ sơ đồ các topo slide 31).}\end{figure}

\subsection{Z-Wave}
Z-Wave được ứng dụng chủ yếu trong nhà thông minh để gắn thiết bị điện, điều khiển từ Internet và hỗ trợ an ninh/tự động hóa. Theo slide: tần số 868,42 MHz tại châu Âu, 908,42 MHz tại Hoa Kỳ; các tốc độ 9,6; 40; 200 kbit/s; điều chế GFSK ở dải gần 900 MHz; tầm khoảng 30 m trong nhà và 100 m ngoài trời. Thông số phụ thuộc khu vực pháp lý và phiên bản thiết bị. \sourcepage{32}

\subsection{Weightless, Sigfox và LoRaWAN}
\paragraph{Weightless.} Họ tiêu chuẩn mở LPWAN kết nối trạm gốc với hàng nghìn nút. Có Weightless-W, Weightless-P và Weightless-N cho các điều kiện khác nhau. Weightless-W hỗ trợ hai chiều ở phổ trắng truyền hình (TV white spaces), slide lấy ví dụ lĩnh vực dầu khí; Weightless-P cũng hai chiều, thích hợp mạng riêng và các trường hợp cần cả đường lên/đường xuống. Slide chỉ nêu tên Weightless-N mà không khai triển tính chất của biến thể này. \sourcepage{33}

\paragraph{Sigfox.} Dịch vụ kết nối LPWAN theo thuê bao tại các quốc gia được triển khai, tối ưu gói dữ liệu nhỏ gửi không thường xuyên. Ví dụ tiêu biểu: thiết bị đo thông minh hoặc cảm biến phát cảnh báo. \sourcepage{34}

\paragraph{LoRaWAN.} Công nghệ mạng diện rộng công suất thấp do hệ sinh thái LoRa Alliance thúc đẩy, cho truyền hai chiều, hướng tới chi phí thấp, pin lâu, mạng lớn và dịch vụ di động. Ứng dụng: thành phố thông minh và công nghiệp. Cần phân biệt điều chế vô tuyến LoRa ở lớp vật lý với giao thức mạng LoRaWAN. \sourcepage{35}

\subsection{Li-Fi và truyền thông ánh sáng}
Li-Fi (Light Fidelity) dùng ánh sáng nhìn thấy (Visible Light Communication, VLC) để truyền dữ liệu tốc độ cao, có thể tận dụng đèn LED hiệu suất cao. Điểm hạn chế trong mô tả slide là yêu cầu điều kiện truyền quang phù hợp, dễ bị cản bởi vật chắn hoặc mất đường nhìn; chưa được triển khai rộng rãi trong IoT tại thời điểm biên soạn. Tiềm năng: lớp học và phòng trình chiếu thông minh. \sourcepage{36}

\subsection{Mạng di động và 5G}
Các thế hệ mạng tế bào (3G, 4G, LTE-Advanced, 5G) có vùng phủ rộng và khả năng hỗ trợ thiết bị di động. LTE-Advanced được slide lấy làm ví dụ về tiến bộ tốc độ và chế độ tiết kiệm năng lượng: có biến thể/chế độ cho các nút có thể đạt thời gian pin rất dài tùy chu kỳ gửi dữ liệu; bản LTE-A khác có tốc độ tối đa ở mức rất cao. Hai trị số 10 Mb/s và 1 Gb/s trong slide mô tả các trường hợp khác nhau, không phải một cấu hình chung. \sourcepage{37}

Slide 5G (có trích dẫn dự báo trước năm 2020) nêu tăng tốc độ, độ phủ, số người dùng và mức cơ động so với 4G. Bảng so sánh 4G/5G ở slide PDF trang 39 có các ý: dải tần 4G khoảng 600 MHz--5,925 GHz; 5G gồm FR1 (khoảng 450 MHz--6 GHz) và FR2 (khoảng 24,25--52,6 GHz, tùy quy định); băng thông 4G tới 20 MHz, 5G FR1 tới 100 MHz/FR2 tới 400 MHz; khoảng cách sóng mang 4G cố định 15 kHz, 5G mở rộng; tốc độ đỉnh downlink theo slide là 1 Gb/s và 20 Gb/s; độ trễ slide lần lượt 10 ms và dưới 1 ms; quy mô kết nối 2.000 so với 1 triệu thiết bị/km$^2$; tốc độ di động 350 so với 500 km/h. Đây là mục tiêu/giá trị minh họa trong bảng, không phải hiệu năng bảo đảm mọi vùng phủ. \sourcepage{38--39}

\subsection{LTE-M và NB-IoT}
\paragraph{LTE-M.} LTE-M, tên đầy đủ LTE Machine-Type Communication, gồm eMTC (enhanced MTC), là công nghệ kết nối di động LPWAN do 3GPP thiết kế cho thiết bị M2M/IoT. eMTC LTE Cat M1 đã được đặc tả ở 3GPP Release 13 (tháng 6 năm 2016 theo slide). Hệ 3GPP còn có NB-IoT và EC-GSM-IoT. \sourcepage{40}

\paragraph{NB-IoT.} Narrowband IoT là tiêu chuẩn LPWAN cho các thiết bị truyền dữ liệu ít, chú trọng mức tiêu thụ điện thấp, dung lượng mạng, hiệu quả phổ và vùng phủ sâu (ví dụ trong công trình). Slide cho rằng thời gian dùng pin hơn 10 năm khả thi ở một số kịch bản; điều này phụ thuộc dữ liệu gửi, vùng phủ, pin và cấu hình hoạt động. \sourcepage{41}

\begin{table}[htbp]\centering\small
\caption{So sánh khái quát hai công nghệ LTE-M và NB-IoT từ slide PDF trang 50.}
\begin{tabularx}{\textwidth}{p{3.4cm}XX}\toprule
Đặc tính & LTE-M & NB-IoT\\\midrule
Tên chuẩn & LTE Cat M1 (theo slide) & LTE Cat NB1\\
Băng thông & Khoảng 1,08 MHz, trong kênh LTE 1,4 MHz & Tới 180 kHz\\
Thông lượng & 36 kb/s (giá trị trong bảng slide) & 30/60 kb/s (giá trị trong bảng slide)\\
Triển khai & Trong các mạng LTE hiện hữu & In-band, guard-band hoặc độc lập trong một số phổ GSM\\
Khả năng giọng nói & Có hỗ trợ thoại/VoLTE theo cấu hình & Không được thiết kế cho thoại\\
Di động/chuyển giao & Có & Hạn chế\\
Thời lượng pin theo bảng & Có thể tới 10 năm & Có thể tới 10 năm\\\bottomrule
\end{tabularx}\end{table}

\subsection{Bảng tổng hợp phổ, tầm phủ và tốc độ các công nghệ}
\begingroup\small\sloppy\begin{longtable}{p{2.6cm}p{3.5cm}p{2.25cm}p{3.6cm}}
\caption{Bảng đối chiếu slide PDF trang 42; số liệu tham khảo theo slide, không đại diện mọi phiên bản.}\\
\toprule Công nghệ & Tần số theo slide & Tầm phủ & Tốc độ theo slide\\\midrule\endfirsthead
\toprule Công nghệ & Tần số theo slide & Tầm phủ & Tốc độ\\\midrule\endhead
Bluetooth & 2,4 GHz & 100 m & 125 kbit/s--Mb/s\\
Wi-Fi & 2,4 hoặc 5 GHz & 50 m & 150--600 Mb/s\\
NFC & 13,56 MHz & 4 cm & 100--420 kbit/s\\
LoRaWAN & 867--869 MHz (Âu); 902--928 MHz (Bắc Mỹ) & 15 km & 0,3--50 kbit/s\\
Cellular & 900/1800/\allowbreak 1900/2100 MHz & Bảng nguồn ghi 30 m (nút--trạm), không phải tầm phủ mạng tế bào nói chung & 21 Mb/s (3G+); 600 Mb/s (4G)\\
Z-Wave & 865--926 MHz (ISM) & 100 m & 100 kbit/s\\
Zigbee & 2,4 GHz (ISM) & 100 m & 20--250 kbit/s\\
Sigfox & 900 MHz & 3--50 km & 10--1000 bit/s\\\bottomrule
\end{longtable}\endgroup

\subsection{Bảo mật ở lớp truyền thông}
Các giao thức được slide nêu bao gồm:
\begin{itemize}
\item \textbf{IPsec (Internet Protocol Security):} cơ chế bảo mật ở lớp mạng IP, cung cấp các phương thức xác thực, toàn vẹn và mã hóa khi được triển khai/cấu hình đúng. Được sử dụng với IPv6 nhưng không nên hiểu IPsec mặc nhiên được bật trên mọi mạng IPv6.
\item \textbf{TLS (Transport Layer Security):} bảo vệ kênh kết nối dựa trên TCP, ví dụ HTTPS hoặc MQTT qua TLS.
\item \textbf{DTLS (Datagram Transport Layer Security):} giải pháp bảo mật cho truyền thông kiểu datagram như UDP, ví dụ CoAP/DTLS.
\item \textbf{IEEE 1888.3:} thành viên họ tiêu chuẩn dành cho mạng điều khiển cộng đồng xanh phổ cập (Ubiquitous Green Community Control Network). Mục đích theo slide là hạn chế truy cập trái phép, ngăn rò rỉ dữ liệu và bảo đảm tính bí mật/toàn vẹn dữ liệu truyền.
\end{itemize}
Các thuộc tính cốt lõi là bảo mật nội dung, xác thực đối tác và phát hiện sửa đổi; hệ thống IoT vẫn phải quản lý khóa, phân quyền nút và cập nhật phần mềm. \sourcepage{43}

\subsection{Định tuyến và giao thức RPL}
Định tuyến (routing) là quá trình chọn đường chuyển lưu lượng trong một mạng hoặc qua nhiều mạng. Với rất nhiều thiết bị có tài nguyên hạn chế, chạy pin và liên kết chất lượng không ổn định, lựa chọn đường tốt có ảnh hưởng lớn tới khả năng mở rộng, thời gian gửi dữ liệu và thời lượng pin.

\term{RPL}{Routing Protocol for Low-Power and Lossy Networks} là chuẩn định tuyến IPv6 dành cho mạng công suất thấp và liên kết dễ mất gói. Slide nêu có thể sử dụng trên IEEE 802.15.4, hỗ trợ các kiểu truyền một--một, nhiều--một và một--nhiều. Hình trang 45 biểu diễn hai cấu hình \emph{DODAG} (Destination-Oriented Directed Acyclic Graph): non-storing và storing mode; các nút tổ chức tương quan với nút gốc, còn phương thức duy trì đường đi phụ thuộc chế độ hoạt động. \sourcepage{44--45}

\begin{figure}[htbp]\centering
\begin{tikzpicture}[>=Latex, every node/.style={font=\small}, leaf/.style={circle,draw=orangeiot,fill=orange!18,minimum size=8mm}, root/.style={circle,draw=navy,fill=blue!20,minimum size=10mm}]
\node[root](r) at (0,3){Gốc};
\node[leaf](a) at (-2,1.7){A};\node[leaf](b) at (0,1.7){B};\node[leaf](c) at (2,1.7){C};
\node[leaf](d) at (-2.6,0){D};\node[leaf](e) at (-.9,0){E};\node[leaf](f) at (1,0){F};\node[leaf](g) at (2.7,0){G};
\foreach \x/\y in {a/r,b/r,c/r,d/a,e/a,e/b,f/b,f/c,g/c}{\draw[->,teal,thick](\x)--(\y);}
\end{tikzpicture}
\caption{Một đồ thị định tuyến hướng về nút gốc kiểu DODAG của RPL (minh họa vẽ lại từ slide trang 45).}\end{figure}

\section{Các loại mạng IoT}
Việc kết nối tin cậy giữa cảm biến, thiết bị và nền tảng IoT quyết định hoạt động của ứng dụng. Slide chia các mạng đang dùng thành bốn nhóm: mạng tế bào, mạng cục bộ/cá nhân (LAN/PAN), mạng diện rộng công suất thấp (LPWAN), và mạng mắt lưới (mesh). Đây là các góc nhìn phân loại theo phạm vi/vận hành và topo; một hệ thống thực tế có thể kết hợp chúng. \sourcepage{47}

\subsection{Mạng tế bào (Cellular Networks)}
Mạng tế bào sử dụng cơ sở hạ tầng nhà khai thác di động giống điện thoại thông minh để kết nối thiết bị IoT. Ưu điểm là vùng phủ địa lý rộng và hỗ trợ tính di động; hạn chế là điện năng và phí dịch vụ, cũng như khả năng phủ sóng kém ở hầm sâu, tủ kỹ thuật kín, mỏ hoặc công trình nhiều lớp vật cản. LTE-M và NB-IoT là hai lựa chọn IoT di động được đề cập; slide cho rằng NB-IoT có thể tiêu thụ ít năng lượng và giá thành thấp hơn LTE-M, nhưng phạm vi hỗ trợ dịch vụ phụ thuộc mạng triển khai ở thời điểm khảo sát. \sourcepage{48--50}

\subsection{Mạng cục bộ và cá nhân (LAN/PAN)}
PAN (Personal Area Network) dùng cho kết nối các thiết bị gần cá nhân; LAN (Local Area Network) kết nối các nút trong nhà, khu vực hoặc cơ sở. Phạm vi ngắn, chi phí sử dụng thường thấp; độ tin cậy phụ thuộc nhiễu, vật cản, mật độ nút và cấu hình. Hai ví dụ chính trong slide là Wi-Fi và Bluetooth. Có thể dùng gateway chuyển dữ liệu của mạng LAN/PAN ra mạng diện rộng. \sourcepage{51--52}

\subsection{Mạng diện rộng công suất thấp (LPWAN)}
LPWAN cho tầm xa hơn các kết nối Wi-Fi/Bluetooth thông dụng nhưng thường dùng ít năng lượng hơn nhiều công nghệ di động băng rộng. Đổi lại, tốc độ thấp và phù hợp các bản tin cảm biến nhỏ, gửi cách quãng. LoRaWAN, Sigfox, Weightless và các biến thể LPWAN tế bào như LTE-M/NB-IoT là các ví dụ đã giới thiệu trong phần trên. \sourcepage{53}

\subsection{Mạng mắt lưới (Mesh Networks)}
Trong mạng mesh, các nút có thể phối hợp chuyển tiếp dữ liệu cho nhau để đến gateway. Tầm mỗi bước nhảy có thể ngắn nhưng tổng vùng phủ được kéo dài qua nhiều nút trung gian. Điều này hữu ích cho cảm biến phân tán trong một tòa nhà; đổi lại cần đủ nút, bộ chuyển tiếp và cấu hình đường đi để tránh các vùng không phủ. Sơ đồ Zigbee và RPL ở các phần trước là những minh họa thích hợp cho kiến trúc đa chặng. \sourcepage{54}

\begin{figure}[htbp]\centering
\begin{tikzpicture}[>=Latex,scale=1,transform shape]
\draw[->,thick] (0,0)--(12,0) node[right]{Tầm phủ};
\draw[->,thick] (0,0)--(0,5) node[above,align=center]{Tốc độ dữ liệu};
\fill[blue!12] (.6,2.8) rectangle (3.3,4.6);\node[align=center] at (1.95,3.8){Wi-Fi\\Bluetooth};
\fill[cyan!16] (7,.65) rectangle (10.15,2.45);\node[align=center] at (8.55,1.55){LPWAN\\LoRa, Sigfox, NB-IoT};
\fill[orange!15] (4.0,2.5) rectangle (10.7,4.65);\node[align=center] at (7.45,3.65){Mạng tế bào\\2G/3G/4G/5G};
\end{tikzpicture}
\caption{So sánh định tính tầm phủ--tốc độ theo bố cục slide trang 55; vị trí vùng chỉ mang tính mô tả, không theo thang số.}\end{figure}

\begin{table}[htbp]\centering\small
\caption{Tóm tắt bốn nhóm mạng IoT.}
\begin{tabularx}{\textwidth}{p{2.6cm}XXX}\toprule
Loại mạng & Điểm mạnh & Hạn chế & Ví dụ\\\midrule
Cellular & Phủ rộng, di động & Năng lượng/chi phí; vùng phủ trong nhà sâu & 4G, LTE-M, NB-IoT, 5G\\
LAN/PAN & Chi phí triển khai thấp, tốc độ tùy công nghệ & Tầm ngắn, dễ bị cản/nhiễu & Wi-Fi, BLE\\
LPWAN & Tầm xa, phù hợp cảm biến chạy pin & Thông lượng thấp, hạn chế payload/chu kỳ & LoRaWAN, Sigfox, Weightless\\
Mesh & Mở rộng đa chặng, có tuyến thay thế & Cần nút chuyển tiếp, phức tạp định tuyến & Zigbee mesh, RPL\\\bottomrule
\end{tabularx}\end{table}

\subsection{Sơ đồ phân nhóm các công nghệ mạng}
Slide trang 46 sắp xếp mạng theo \emph{cellular}, \emph{wired}, \emph{satellite}, \emph{short-range wireless}, và \emph{LP-WAN / LP wireless}. Các mục được nêu gồm: 2G/2.5G, 3G, 4G/LTE, 5G; POTS, Fiber, Ethernet, ATM/Frame Relay, MPLS, Powerline; C-Band và L-Band; Wi-Fi/802.11, Bluetooth/BLE, Zigbee/802.15.4, 900 MHz proprietary, 6LoWPAN, Wi-Fi HaLow; WiMAX, Sigfox, LoRa, NB-IoT, LTE-M, Fixed Microwave và WiSUN. Hình này bổ sung các loại truyền dẫn chứ không thay đổi cách phân nhóm bốn loại của slide 47.

\section{Giao thức dữ liệu của IoT}
\subsection{Vai trò của giao thức ứng dụng và mô hình trao đổi}
Slide bắt đầu bằng việc phân biệt truyền thông dữ liệu giữa các thiết bị với khả năng truy cập Internet, nhấn mạnh các nút công suất thấp có thể giao tiếp trực tiếp bằng liên kết có dây hoặc vô tuyến mà không phải lúc nào cũng cần truy cập Internet. Trong ngăn xếp giao thức, Wi-Fi/BLE/LoRa thuộc nhóm công nghệ kết nối bên dưới, còn HTTP, CoAP, MQTT, WebSocket, AMQP và XMPP cung cấp cách thức trao đổi thông điệp ở tầng ứng dụng. \sourcepage{56,60}

Một nút IoT có thể thu thập dữ liệu, phân tích dữ liệu và điều khiển nút khác. Tùy hệ thống, trước khi tương tác với dịch vụ, nút cần được nhận dạng và đăng ký. Slide gọi quá trình đó là \emph{node subscription}. Trong ngữ cảnh truyền dữ liệu, cần phân biệt \emph{đăng ký một nút với dịch vụ} và \emph{đăng ký theo dõi một chủ đề dữ liệu (subscribe topic)}. \sourcepage{57}

\paragraph{Publish/subscribe (xuất bản/đăng ký).} Nút xuất bản (publisher) gửi bản tin lên một chủ đề/kênh; các nút đăng ký (subscriber) nhận dữ liệu tương ứng khi bản tin được xuất bản. Thành phần trung gian có thể là broker; cách tiếp cận tách người gửi và người nhận và thuận lợi cho nhiều nút trao đổi thường xuyên. \sourcepage{58--59}

\paragraph{Request/response (yêu cầu/phản hồi).} Client chủ động yêu cầu dữ liệu hoặc tác vụ và server trả kết quả. Mô hình quen thuộc với ứng dụng client/server và các dịch vụ REST. Slide nhận xét publish/subscribe thường thuận lợi hơn với triển khai IoT lớn trao đổi dữ liệu nhiều, trong khi request/response dễ tích hợp ứng dụng client/server. \sourcepage{59}

\begin{figure}[htbp]\centering
\begin{tikzpicture}[>=Latex,font=\small,block/.style={draw=navy,rounded corners,minimum width=29mm,minimum height=10mm,fill=blue!6}]
\node[block](cli) at (0,2.7){Client};\node[block](serv) at (5,2.7){Server};
\draw[->,thick,teal] (cli)--node[above]{Request}(serv);
\draw[<-,thick,orangeiot] (cli) to[bend right=22] node[below]{Response}(serv);
\node[block](pub) at (0,0){Publisher};\node[block](broker) at (5,0){Broker};\node[block](sub) at (10,0){Subscriber};
\draw[->,thick,teal](pub)--node[above]{Publish}(broker);
\draw[->,thick,teal](broker)--node[above]{Push chủ đề}(sub);
\end{tikzpicture}
\caption{Hai mô hình trao đổi dữ liệu ở slide trang 57--59.}\end{figure}

\subsection{HTTP: Hypertext Transfer Protocol}
HTTP là giao thức ứng dụng phổ biến của Web, thường theo mô hình \emph{pull}: client gửi yêu cầu lấy tài nguyên, server trả phản hồi. HTTP không lưu trạng thái yêu cầu trước đó ở bản thân giao thức (stateless), hỗ trợ thương lượng kiểu nội dung và biểu diễn tài nguyên để các hệ thống được phát triển độc lập vẫn trao đổi dữ liệu. \sourcepage{61--62}

Những phương thức HTTP phổ biến: \texttt{GET} (đọc tài nguyên), \texttt{POST} (tạo/xử lý dữ liệu), \texttt{PUT} (thay thế/cập nhật tài nguyên), \texttt{PATCH} (cập nhật một phần), \texttt{DELETE} (xóa). CRUD: Create--Read--Update--Delete là các thao tác ứng dụng tương ứng ở mức khái niệm, không phải ánh xạ duy nhất cho mọi API. \sourcepage{62}

Theo slide, HTTP có hạn chế đối với một số hệ thống IoT: quá nhiều cặp yêu cầu--phản hồi khi nhiều cảm biến hoạt động khiến phát sinh độ trễ và năng lượng; thông thường client khởi xướng yêu cầu thay vì server tự đẩy bản tin như mô hình publish/subscribe; HTTP/TCP có overhead và chi phí kết nối, không tối ưu cho một số cảm biến dùng pin. Những hạn chế này được mô tả cho cách triển khai HTTP phổ biến; các biến thể HTTP mới, keep-alive, server push hoặc kỹ thuật khác có thể thay đổi chi phí thực tế. \sourcepage{63--64}

\begin{figure}[htbp]\centering
\begin{tikzpicture}[>=Latex,font=\small]
\node at (0,4.8){Client};\node at (6,4.8){Server};
\draw[dashed](0,4.4)--(0,0);\draw[dashed](6,4.4)--(6,0);
\draw[->,thick,navy](0,4)--node[above]{TCP: SYN}(6,3.5);
\draw[<-,thick,navy](0,2.9)--node[above]{SYN/ACK}(6,3.4);
\draw[->,thick,navy](0,2.8)--node[above]{ACK}(6,2.4);
\draw[->,thick,teal](0,2)--node[above]{HTTP GET /temperature}(6,1.6);
\draw[<-,thick,orangeiot](0,.8)--node[above]{HTTP 200 OK + dữ liệu}(6,1.3);
\end{tikzpicture}
\caption{Chuỗi bắt tay TCP và yêu cầu--phản hồi HTTP, dựng theo ý tưởng hình slide trang 63.}\end{figure}

\subsection{TCP và UDP}
\textbf{TCP} là giao thức hướng kết nối, sử dụng bắt tay thiết lập phiên, đánh số thứ tự, xác nhận/điều khiển truyền lại để cung cấp luồng byte tin cậy. \textbf{UDP} truyền datagram không thiết lập phiên TCP và không tự đảm bảo xác nhận, thứ tự hoặc truyền lại. Do đó TCP thường có chi phí thiết lập/trạng thái cao hơn, còn UDP thích hợp với giao thức tự quản lý độ tin cậy hoặc yêu cầu thông điệp ngắn như CoAP. Slide trang 65 minh họa TCP bằng SYN--SYN/ACK--ACK và UDP bằng các datagram yêu cầu--phản hồi không có bắt tay. Việc nói 'UDP không tin cậy' chỉ có nghĩa UDP không tự bảo đảm chuyển phát, không có nghĩa ứng dụng UDP không thể cài đặt xác nhận. \sourcepage{65}

\subsection{CoAP: Constrained Application Protocol}
CoAP là giao thức truyền tài nguyên kiểu Web thiết kế cho nút và mạng bị hạn chế về băng thông, độ sẵn sàng và tài nguyên xử lý. Nó sử dụng mô hình client/server tương tự REST/HTTP: server công bố tài nguyên qua URI, client dùng GET, POST, PUT, DELETE. Bản tin CoAP được thiết kế gọn và thường dựa trên UDP; việc xác nhận chuyển phát phải do cơ chế CoAP xử lý khi cần. \sourcepage{66}

\begin{figure}[htbp]\centering
\begin{tikzpicture}[>=Latex,font=\small,box/.style={draw=navy,fill=blue!5,rounded corners,minimum width=23mm,minimum height=10mm}]
\node[box](web) at (0,0){HTTP client};\node[box](proxy) at (4.2,0){HTTP/CoAP proxy};\node[box](sensor) at (8.8,0){Cảm biến CoAP};
\draw[<->,thick,navy](web)--node[above]{HTTP / TCP}(proxy);
\draw[<->,thick,teal](proxy)--node[above]{CoAP / UDP}(sensor);
\end{tikzpicture}
\caption{Liên kết client Web với mạng cảm biến giới hạn qua proxy CoAP (theo slide 67).}\end{figure}

\paragraph{Đặc trưng CoAP.} Slide nêu năm nhóm ưu điểm: phục vụ giao tiếp M2M bị ràng buộc; trao đổi bản tin không đồng bộ; overhead thấp và dễ phân tích cú pháp; hỗ trợ URI và kiểu nội dung; hỗ trợ proxy và cache. CoAP có thể làm cầu nối với Web bằng gateway/proxy. \sourcepage{67--68}

\paragraph{Bốn loại bản tin.} CoAP định nghĩa:
\begin{enumerate}[leftmargin=2em]
\item \textbf{Confirmable (CON):} bản tin cần ACK; nếu không nhận được sẽ có cơ chế gửi lại.
\item \textbf{Non-confirmable (NON):} bản tin không đòi ACK ở lớp thông điệp CoAP (\emph{fire-and-forget}).
\item \textbf{Acknowledgement (ACK):} xác nhận đã nhận bản tin CON.
\item \textbf{Reset (RST):} từ chối/đặt lại bản tin hoặc xử lý quan sát không còn hợp lệ.
\end{enumerate}

CoAP dùng Message ID 16 bit cho ghép bản tin và phát hiện bản sao, cùng cơ chế truyền lại có back-off đối với CON. Bảo mật có thể sử dụng DTLS. \sourcepage{69}

\begin{figure}[htbp]\centering
\begin{tikzpicture}[>=Latex,font=\small]
\node at (0,4.8){Client};\node at (7,4.8){Server};
\draw[dashed](0,4.4)--(0,.1);\draw[dashed](7,4.4)--(7,.1);
\draw[->,thick,teal](0,3.9)--node[above]{CON, MID=0x4A51, GET}(7,3.35);
\draw[<-,thick,orangeiot](0,2.3)--node[above]{ACK, MID=0x4A51, 2.05 Content}(7,2.9);
\draw[->,thick,teal](0,1.7)--node[above]{NON, MID=0x4A52}(7,1.2);
\node[align=left] at (3.5,.45){ACK cho CON; không ACK cho NON.};
\end{tikzpicture}
\caption{Ví dụ luồng CoAP confirmable và non-confirmable, vẽ lại theo slide trang 70--71.}\end{figure}

\paragraph{Ví dụ tra cứu tài nguyên.} Các hình slide 70--71 cho thấy bản tin \texttt{CON GET /temperature} với \texttt{Token} được phản hồi bằng \texttt{ACK 2.05 Content} chứa nhiệt độ \texttt{20.1 C}, hoặc phản hồi lỗi \texttt{4.04 Not Found} khi tài nguyên không tồn tại. Hình trang 70 cũng thể hiện CON--ACK, CON--RST và NON một chiều, để làm rõ khác nhau giữa xác nhận và từ chối. \sourcepage{70--71}

\subsection{MQTT: Message Queuing Telemetry Transport}
MQTT là giao thức gửi thông điệp nhẹ, hữu ích để phân phối dữ liệu telemetry tới các client có tài nguyên và kết nối mạng hạn chế. Nó dùng mô hình publish/subscribe và rất phổ biến trong giao tiếp M2M/IoT. Thành phần trung tâm là \textbf{broker}: publisher gửi dữ liệu theo chủ đề (topic); broker xác định người đăng ký để chuyển tiếp; subscriber nhận bản tin không cần duy trì kết nối trực tiếp với từng publisher. \sourcepage{72--73}

\begin{figure}[htbp]\centering
\begin{tikzpicture}[>=Latex,font=\small,box/.style={draw=navy,rounded corners,fill=blue!5,minimum height=12mm,minimum width=28mm}]
\node[box](p) at (0,0){Publisher};\node[box,fill=cyan!15](b) at (4.5,0){MQTT Broker};
\node[box](s1) at (9,1.4){Subscriber 1};\node[box](s2) at (9,-1.4){Subscriber 2};
\draw[->,thick,teal](p)--node[above]{PUBLISH topic}(b);
\draw[->,thick,orangeiot](b)--node[above,sloped]{PUBLISH}(s1);
\draw[->,thick,orangeiot](b)--node[below,sloped]{PUBLISH}(s2);
\end{tikzpicture}
\caption{Kiến trúc MQTT publisher--broker--subscriber (dựng lại từ hình slide trang 73).}\end{figure}

\paragraph{Bốn giai đoạn phiên MQTT.} Theo slide: (1) \textbf{Connection}: client thiết lập kết nối TCP/IP đến broker (cổng chuẩn hoặc cổng riêng do server định nghĩa); (2) \textbf{Authentication}: xác thực nếu broker yêu cầu; (3) \textbf{Communication}: publish và subscribe bản tin; (4) \textbf{Termination}: đóng phiên/kết nối. Với Client ID và thiết lập phiên thích hợp, broker có thể tiếp tục một phiên trước đó. \sourcepage{74}

\paragraph{Ba mức Quality of Service (QoS).}
\begin{table}[htbp]\centering\small
\caption{Ba mức QoS trong MQTT (slide trang 75).}
\begin{tabularx}{\textwidth}{p{1.4cm}p{3cm}X}\toprule
QoS & Ý nghĩa & Đặc điểm\\\midrule
0 & At most once (tối đa một lần) & Không bảo đảm thông điệp đến đích; không có cơ chế xác nhận chuyển phát QoS 0\\
1 & At least once (ít nhất một lần) & Bảo đảm theo cơ chế xác nhận/gửi lại; có thể trùng lặp bản tin\\
2 & Exactly once (đúng một lần) & Có quy trình bắt tay để bảo đảm xử lý chuyển phát đúng một lần theo ngữ nghĩa MQTT\\\bottomrule
\end{tabularx}\end{table}

\paragraph{Ví dụ ứng dụng ở thời điểm giảng dạy.} Slide nêu Facebook sử dụng MQTT trong hệ thống nhắn tin để tiết kiệm pin và truyền tin nhanh trên mạng không ổn định; đồng thời nêu các nền tảng AWS, Google Cloud, IBM Bluemix, Microsoft Azure và ThingOn IoT Platform có hỗ trợ MQTT. Những khẳng định mang yếu tố thời điểm này được giữ theo tư liệu giảng dạy, không coi là xác nhận về sản phẩm hiện tại. \sourcepage{76}

\subsection{WebSocket}
WebSocket liên quan HTTP trong bước thiết lập kết nối: trình duyệt/client thực hiện cơ chế bắt tay nâng cấp kết nối và duy trì một kênh truyền trên TCP để cả hai phía trao đổi bản tin cho đến khi đóng. So với cơ chế nhiều lượt request/response, WebSocket cho phép truyền hai chiều trên kết nối lâu dài, thuận lợi cho ứng dụng thời gian thực; slide nhận xét nó vẫn có overhead đáng kể đối với các nút IoT rất hạn chế tài nguyên. Hình trang 78 minh họa cột HTTP với những cặp request--response lặp lại so với cột WebSocket với bước handshake rồi truyền hai chiều trực tiếp. \sourcepage{77--78}

\begin{figure}[htbp]\centering
\begin{tikzpicture}[>=Latex,font=\small]
\node at (0,3.5){Client};\node at (6,3.5){Server};\draw[dashed](0,3.2)--(0,0);\draw[dashed](6,3.2)--(6,0);
\draw[->,thick,navy](0,2.9)--node[above]{HTTP Upgrade: websocket}(6,2.65);
\draw[<-,thick,navy](0,2.22)--node[above]{101 Switching Protocols}(6,2.5);
\draw[->,thick,teal](0,1.6)--node[above]{Thông điệp client $\to$ server}(6,1.4);
\draw[<-,thick,orangeiot](0,.8)--node[above]{Thông điệp server $\to$ client}(6,1.0);
\end{tikzpicture}
\caption{Thiết lập và truyền hai chiều với WebSocket, vẽ lại theo slide 77--78.}\end{figure}

\subsection{AMQP: Advanced Message Queuing Protocol}
AMQP là tiêu chuẩn mở thiên về giao tiếp hướng thông điệp, hỗ trợ publish/subscribe. Tương tự cách mô tả MQTT trong slide, các ứng dụng có thể sử dụng các ngữ nghĩa chuyển phát tối đa một lần, ít nhất một lần hoặc đúng một lần tùy giao thức/triển khai. Kênh TCP có thể được bảo vệ với TLS/SSL, đồng thời có cơ chế xác thực SASL. Hai thành phần cốt lõi được slide nêu:
\begin{itemize}
\item \textbf{Exchange:} nhận bản tin và chọn (định tuyến) queue phù hợp dựa trên các quy tắc đã đặt.
\item \textbf{Message queue:} lưu bản tin chờ được chuyển tới phía tiêu thụ/nhận.
\end{itemize}
Kiến trúc này phù hợp hệ thống có xử lý hàng đợi, định tuyến thông điệp linh hoạt và tích hợp backend. \sourcepage{79}

\begin{figure}[htbp]\centering
\begin{tikzpicture}[>=Latex,font=\small,box/.style={draw=navy,rounded corners,minimum width=26mm,minimum height=10mm,fill=blue!5}]
\node[box](p) at (0,0){Producer};\node[box](ex) at (4,0){Exchange};\node[box](q) at (8,0){Queue};\node[box](c) at (12,0){Consumer};
\draw[->,thick,teal](p)--(ex);\draw[->,thick,teal](ex)--node[above]{Routing}(q);\draw[->,thick,teal](q)--(c);
\end{tikzpicture}
\caption{Luồng thông điệp cơ bản AMQP theo slide trang 79.}\end{figure}

\subsection{XMPP: Extensible Messaging and Presence Protocol}
XMPP là giao thức tiêu chuẩn hóa qua IETF để trao đổi thông tin trên Internet, hoạt động gần thời gian thực, độ trễ thấp theo bối cảnh ứng dụng và có tính độc lập nền tảng/phân tán. Theo slide, XMPP dựa trên TCP, hỗ trợ cả publish/subscribe và request/response; client kết nối tới server và trao đổi bản tin dựa trên \term{XML}{Extensible Markup Language}. Nhờ cấu trúc thông điệp mở rộng, nó có thể phục vụ ứng dụng liên lạc IoT, nhưng XML thường tạo thêm overhead so với các định dạng nhị phân nhỏ gọn. \sourcepage{80}

\subsection{Bảng so sánh các giao thức dữ liệu}
\begingroup\footnotesize\sloppy\begin{longtable}{p{1.8cm}p{3.45cm}p{1.8cm}p{4.15cm}}
\caption{Tổng hợp từ các slide trang 56--80; một số nhận xét là phân tích hệ thống hóa.}\\
\toprule Giao thức & Mô hình & Vận chuyển & Ghi chú trong IoT\\\midrule\endfirsthead
\toprule Giao thức & Mô hình & Vận chuyển & Ghi chú trong IoT\\\midrule\endhead
HTTP & Request/response, client/server & TCP & Phổ biến REST, stateless; overhead kết nối/yêu cầu có thể cao\\
CoAP & REST, request/response; có hỗ trợ quan sát & UDP & Nhẹ, CON/NON/ACK/RST, URI, proxy/cache, DTLS\\
MQTT & Publish/subscribe qua broker & TCP & Chủ đề, phiên, QoS 0/1/2\\
WebSocket & Kết nối song công dài hạn & TCP & Nâng cấp từ HTTP; trao đổi hai chiều liên tục\\
AMQP & Trao đổi thông điệp, hàng đợi & TCP & Exchange, queue, routing, TLS/SASL\\
XMPP & Cả publish/subscribe và request/response & TCP & Bản tin XML, mở rộng, gần thời gian thực\\\bottomrule
\end{longtable}\endgroup

\section{Tổng kết kiến thức và gợi ý ứng dụng}
\begin{ghinho}[title={Những khái niệm cần phân biệt}]
\textbf{Định danh/địa chỉ:} URI, EPC, uCode, IPv4, IPv6.\par
\textbf{Kết nối vật lý và mạng:} PLC/X10, NFC, UWB, Wi-Fi/HaLow, BLE, Zigbee, Z-Wave, Weightless, Sigfox, LoRaWAN, Li-Fi, 3G/4G/5G, LTE-M, NB-IoT.\par
\textbf{Bảo vệ, định tuyến:} IPsec, TLS/DTLS, IEEE 1888.3, RPL.\par
\textbf{Kiến trúc mạng:} cellular, LAN/PAN, LPWAN, mesh.\par
\textbf{Cách trao đổi dữ liệu:} request/response và publish/subscribe; HTTP, CoAP, MQTT, WebSocket, AMQP, XMPP.
\end{ghinho}

Một ví dụ tổng hợp: cảm biến chạy pin trong nhà có thể kết nối BLE hoặc Zigbee tới gateway; gateway kết nối Wi-Fi hoặc Ethernet lên máy chủ; dữ liệu được xuất bản qua MQTT để nhiều dịch vụ nhận cùng thời điểm. Một cảm biến cần REST tối giản có thể dùng CoAP/UDP; hệ thống Web có thể dùng HTTP để truy vấn và WebSocket để cập nhật liên tục. Nút ở vị trí xa không có Wi-Fi có thể cân nhắc LoRaWAN hoặc LTE-M/NB-IoT. Các lựa chọn này là \emph{minh họa thiết kế được suy ra từ bài học}, không phải kịch bản thí nghiệm bắt buộc trong slide.

\section{Đối chiếu toàn bộ 81 trang slide}
Bảng dưới đây xác định trang tài liệu gốc liên quan tới từng mục, để người đọc rà soát các hình và bảng của bài giảng. Trang 1 là bìa; trang 2 là mục lục; trang 81 là lời cảm ơn. Các trang mang hình không trích xuất được thành chữ được mô tả và minh họa lại trong phần tương ứng.
\begin{longtable}{p{3.0cm}p{10cm}}
\toprule Trang PDF & Nội dung có trong tài liệu/chương\\\midrule\endhead
1--2 & Bìa, học phần, giảng viên, mục lục\\
3--8 & Định nghĩa giao tiếp, H2M, M2M, M2H, H2H và các ví dụ\\
9--10 & Tổng quan giao thức IoT; định danh nút mạng\\
11--15 & URI/EPC/uCode; IPv4 và IPv6; hai hình so sánh địa chỉ và lợi ích IPv6\\
16--18 & Các lợi ích của IPv6; biểu tượng các chuẩn kết nối IoT\\
19--21 & PLC và X10; bảng thông số X10\\
22--23 & NFC; bảng tần số, tầm xa, chế độ và mã hóa\\
24--26 & UWB; Wi-Fi, Wi-Fi Direct và bảng tính năng Wi-Fi HaLow\\
27--28 & BLE và bảng so sánh Bluetooth Classic/BLE\\
29--31 & Zigbee, tính chất giao thức, hình sao--cây--mesh\\
32--35 & Z-Wave, Weightless, Sigfox, LoRaWAN\\
36--39 & Li-Fi, mạng tế bào, 5G, bảng 4G so với 5G\\
40--42 & LTE-M, NB-IoT, bảng so sánh đa công nghệ vô tuyến\\
43--46 & IPsec/TLS/DTLS/IEEE 1888.3, RPL và hai sơ đồ phân loại/định tuyến\\
47--55 & Bốn loại mạng IoT, ví dụ cellular/LAN-PAN/LPWAN/mesh, biểu đồ tầm xa--thông lượng\\
56--60 & Giao thức dữ liệu, đăng ký nút, publish/subscribe, request/response, liệt kê giao thức\\
61--65 & HTTP, phương thức HTTP, hạn chế IoT, sơ đồ HTTP và TCP so với UDP\\
66--71 & CoAP, proxy, REST, CON/NON/ACK/RST, luồng gói và ví dụ GET nhiệt độ\\
72--76 & MQTT, kiến trúc broker, bốn giai đoạn phiên, QoS 0--2, ví dụ áp dụng\\
77--78 & WebSocket, tương tác hai chiều và so sánh hình với HTTP\\
79--80 & AMQP (exchanges/queues) và XMPP/XML\\
81 & Kết thúc và lời cảm ơn\\\bottomrule
\end{longtable}

\begin{nguonslide}[title={Nguồn và phương pháp biên soạn}]
Tài liệu chính: TS. Nguyễn Đắc Cử, \textit{Cơ sở và thiết kế hệ thống IoT -- Chapter 3: Networking technologies in IoT}, Khoa Điện -- Điện tử, Trường Đại học Phenikaa, bản PDF 81 trang. Cách trình bày ở đây là \textbf{bản tổng hợp tiếng Việt theo chủ đề} chứ không phải bản dịch đối chiếu từng trang. Các sơ đồ đã được vẽ lại, dữ liệu theo bảng slide được ghi chú riêng; những điểm không nhất quán trong slide (như đơn vị gói X10, thiết lập NFC và phạm vi cellular trong bảng so sánh) được chỉ rõ chứ không tự ý thay thế. Khi tệp PDF gốc có mặt tại \texttt{source/chapter3-slides.pdf}, có thể đối chiếu trực tiếp với bảng trang ở trên.
\end{nguonslide}
\chapter{Công nghệ mạng trong Internet vạn vật}
\label{chap:iot-networking}
\begin{ghinho}[title={Mục tiêu và phạm vi}]
Chương này tổng hợp bốn phần của 81 trang slide: (1) giao tiếp trong IoT; (2) công nghệ và giao thức truyền thông, gồm định danh, địa chỉ hóa, công nghệ truyền dẫn, bảo mật và định tuyến; (3) phân loại mạng IoT; (4) giao thức dữ liệu HTTP, CoAP, MQTT, WebSocket, AMQP và XMPP. Các bảng thông số được ghi theo chính slide, không được hiểu là mức cam kết cho mọi phiên bản, thiết bị hoặc quốc gia. 
\end{ghinho}

\section{Giao tiếp trong hệ thống IoT}
\subsection{Khái niệm và vai trò của truyền thông}
Internet vạn vật (Internet of Things, IoT) kết nối các thiết bị thông minh qua mạng, cho phép chúng trao đổi dữ liệu và thực hiện công việc với ít hoặc không cần sự can thiệp trực tiếp của con người. Những thiết bị này thường tích hợp linh kiện điện tử, cảm biến, phần mềm, kết nối mạng và cơ cấu chấp hành. Truyền thông là điều kiện để dữ liệu đo ở một thiết bị trở thành thông tin có ích cho các thiết bị, dịch vụ và người sử dụng ở nơi khác. Không có liên kết đáng tin cậy giữa các thiết bị và nền tảng xử lý, một hệ thống cảm biến rời rạc chưa phát huy đầy đủ giá trị IoT. \sourcepage{3}

\subsection{Bốn hướng tương tác trong IoT}
\paragraph{Người đến máy (Human-to-Machine, H2M).} Con người gửi đầu vào cho hệ thống dưới dạng giọng nói, văn bản, hình ảnh hoặc thao tác trực tiếp; thiết bị thu nhận, phân tích và đáp ứng bằng văn bản, hình ảnh hoặc hành động. H2M phối hợp phần cứng và phần mềm để hỗ trợ các công việc hằng ngày. Ví dụ: nhận dạng giọng nói; nhận dạng khuôn mặt; chấm công sinh trắc học. \sourcepage{4}

\paragraph{Máy đến máy (Machine-to-Machine, M2M).} Hai hoặc nhiều máy trao đổi dữ liệu hoặc lệnh tự động, không đòi hỏi con người tham gia vào mỗi phiên liên lạc. Đường truyền có thể là hữu tuyến hoặc vô tuyến; một kết nối M2M điểm--điểm có thể sử dụng Ethernet hoặc mạng di động. IoT tiếp thu cách vận hành tự động của M2M và mở rộng sang các mạng thiết bị quy mô lớn có dịch vụ/nền tảng đám mây. Ví dụ: máy bán hàng tự động thông báo hết hàng cho nhà phân phối; máy giặt thông minh gửi thông báo hoàn tất; công tơ thông minh theo dõi điện năng và cảnh báo chủ sử dụng. \sourcepage{5--6}

\paragraph{Máy đến người (Machine-to-Human, M2H).} Thiết bị chủ động phát tín hiệu, văn bản, hình ảnh hoặc lời nói hướng tới con người, kể cả khi người đó không trực tiếp phát yêu cầu. Ví dụ: chuông báo cháy, đèn giao thông, thiết bị giám sát sức khỏe. \sourcepage{7}

\paragraph{Người đến người (Human-to-Human, H2H).} Con người liên lạc bằng lời nói, chữ viết, hình vẽ, nét mặt hoặc ngôn ngữ cơ thể. Trong một hệ thống IoT thực tế, H2H vẫn cần thiết để phối hợp xử lý sự cố, quản lý tình huống và đưa ra quyết định khi tự động hóa M2M chưa đáp ứng đủ. \sourcepage{8}

\begin{figure}[htbp]\centering
\begin{tikzpicture}[>=Latex,box/.style={draw=navy,rounded corners,fill=blue!5,minimum width=36mm,minimum height=13mm,align=center},node distance=20mm]
\node[box](human){Con người};\node[box,right=36mm of human](machine){Thiết bị IoT};
\draw[->,thick,teal,bend left=20](human) to node[above]{H2M: lệnh/đầu vào}(machine);
\draw[->,thick,orangeiot,bend left=20](machine) to node[below]{M2H: cảnh báo/kết quả}(human);
\node[box,below=19mm of machine](node2){Thiết bị IoT khác};
\draw[<->,thick,navy](machine)--node[right]{M2M}(node2);
\node[box,below=19mm of human](human2){Con người khác};
\draw[<->,thick,navy](human)--node[left]{H2H}(human2);
\end{tikzpicture}
\caption{Minh họa bốn loại hình tương tác H2M, M2M, M2H, H2H (vẽ lại theo các slide 4--8).}\end{figure}

\section{Công nghệ và giao thức truyền thông IoT}
\subsection{Tổng quan và yêu cầu lựa chọn}
Một hệ thống IoT có thể dùng Bluetooth, Wi-Fi, Wi-Fi Direct, giao thức vô tuyến công suất thấp, mạng di động LTE/LTE-A hoặc các công nghệ khác. Không có giao thức tốt nhất cho mọi ứng dụng: phải cân nhắc tầm phủ sóng, mức tiêu thụ năng lượng, băng thông, độ trễ, chi phí, tính di động, quy mô nút và mức bảo mật. Slide giới thiệu các họ công nghệ: PLC, X10, NFC, UWB, Wi-Fi, BLE, Zigbee, Z-Wave, Weightless, Sigfox, LoRaWAN, Li-Fi, mạng tế bào/5G, LTE-M và NB-IoT. \sourcepage{2,9,18}

\subsection{Định danh, đặt tên và địa chỉ hóa}
\paragraph{Định danh (identification).} Khi một nút gia nhập mạng, nó cần được nhận biết bằng một định danh. Cơ chế định danh giúp quản lý thiết bị, xác định đối tượng giao tiếp và hạn chế xử lý lưu lượng không cần thiết trong mạng rất lớn. Hệ thống còn cần tên và địa chỉ để tìm đúng thiết bị, đúng tài nguyên. Lưu ý về mặt kỹ thuật: việc cấp định danh tự nó không chặn được tấn công hay lưu lượng tràn ngập; phải kết hợp xác thực, phân quyền và cơ chế điều khiển lưu lượng. \sourcepage{10}

\paragraph{Đặt tên tài nguyên và đối tượng.} Ba cơ chế được nêu trong slide gồm:
\begin{itemize}
\item \textbf{URI (Uniform Resource Identifier):} chuỗi nhận dạng tài nguyên trên mạng. Ví dụ trong slide: \texttt{http://www.example.com/temperature/}.
\item \textbf{EPC (Electronic Product Code):} mã định danh đối tượng vật lý, thường mã hóa trên thẻ RFID thụ động. Ví dụ: \texttt{urn:epc:id:sgtin:}\allowbreak\texttt{0614141.012345.62852}.
\item \textbf{uCode (Ubiquitous Code):} mã định danh của hệ thống định danh đối tượng; ví dụ slide: \texttt{00001C0000000000}\allowbreak\texttt{0001000285E7A6E3}. Trang giới thiệu: \url{https://www.uidcenter.org}.
\end{itemize}
\sourcepage{11}

\paragraph{IPv4.} Giao thức Internet phiên bản 4 dùng địa chỉ 32 bit, biểu diễn thường gặp bằng bốn nhóm số thập phân phân tách bằng dấu chấm; tổng không gian lý thuyết là $2^{32}$ (xấp xỉ 4,3 tỷ) địa chỉ. Địa chỉ có thể chia mạng con, định tuyến và được cấu hình thủ công hoặc qua DHCP. Ví dụ slide minh họa dạng \texttt{192.168.5.18}. \sourcepage{12,14}

\paragraph{IPv6.} Giao thức Internet phiên bản 6 dùng địa chỉ 128 bit, tương ứng $2^{128}\approx 3{,}4\times 10^{38}$ địa chỉ; ví dụ ký hiệu rút gọn sử dụng các nhóm số hệ 16 và dấu hai chấm. Mục tiêu thiết kế bao gồm giải quyết hạn chế không gian địa chỉ của IPv4 và hỗ trợ số lượng rất lớn thiết bị. IPv6 có cơ chế tự cấu hình địa chỉ, hỗ trợ multicast và khả năng phối hợp với các cơ chế bảo vệ như IPsec. \sourcepage{13--17}

\begin{table}[htbp]\centering\small
\caption{So sánh IPv4 và IPv6 theo minh họa các slide 12--17.}
\begin{tabularx}{\textwidth}{p{3.15cm}XX}\toprule
Tiêu chí & IPv4 & IPv6\\\midrule
Độ dài địa chỉ & 32 bit & 128 bit\\
Số địa chỉ lý thuyết & $2^{32}\approx4{,}3\cdot10^9$ & $2^{128}\approx3{,}4\cdot10^{38}$\\
Dạng địa chỉ & 4 nhóm số thập phân (ví dụ 192.168.5.18) & Các nhóm số hệ 16, phân cách bằng dấu hai chấm\\
Cấu hình & Thủ công hoặc DHCP & Hỗ trợ tự cấu hình; có thể dùng DHCPv6\\
Ứng dụng trong IoT & Có thể bị hạn chế bởi số địa chỉ công cộng & Tạo không gian địa chỉ rất lớn cho nhiều nút\\\bottomrule
\end{tabularx}\end{table}

\paragraph{Các lợi ích IPv6 mà slide liệt kê.} (1) Dễ nhận diện và trao đổi dữ liệu giữa nhiều thiết bị hơn nhờ không gian địa chỉ lớn; (2) tự cấu hình đơn giản hơn ở quy mô lớn, có thể giảm công cấu hình và chi phí triển khai; (3) hỗ trợ truyền tải dữ liệu và hiệu quả sử dụng băng thông; (4) hỗ trợ các thuộc tính bảo mật về bí mật, xác thực và toàn vẹn khi cấu hình các cơ chế phù hợp; (5) multicast giúp tránh một số thao tác phát quảng bá thường xuyên, giảm số gói phải xử lý và có thể tiết kiệm pin; (6) cơ chế định tuyến và cấu trúc địa chỉ linh hoạt theo cách tổ chức mạng/ứng dụng. Đây là các lợi ích tiềm năng theo slide, không phải bảo đảm rằng bật IPv6 tự động tăng tốc hay bảo mật toàn bộ hệ thống. \sourcepage{16--17}

\begin{ghinho}[title={Phân biệt ba khái niệm}]
\textbf{Tên} cho biết tài nguyên hoặc đối tượng nào cần tìm; \textbf{định danh} phân biệt đối tượng trong hệ thống; \textbf{địa chỉ mạng} giúp chuyển gói đến vị trí mạng đang sử dụng. Tên, mã RFID và địa chỉ IP vì thế có thể cùng tồn tại trên một thiết bị IoT.
\end{ghinho}

\subsection{Công nghệ truyền dẫn qua dây điện: PLC và X10}
\paragraph{PLC (Power-Line Communication).} PLC sử dụng ngay đường dây điện hiện có để mang cả điện năng xoay chiều và tín hiệu dữ liệu. Một nút có thể nhận nguồn điện và trao đổi thông tin mà không cần tuyến dây dữ liệu riêng. Phù hợp cho thiết bị cố định, nhà thông minh, lưới điện thông minh, phương tiện hoặc môi trường có nhiễu vô tuyến cao. Hạn chế thực tế phụ thuộc chất lượng dây, tải và nhiễu từ thiết bị điện. \sourcepage{19}

\paragraph{X10.} Chuẩn điều khiển thiết bị trên đường điện ra đời khoảng thập niên 1970. Tín hiệu điều khiển được đưa lên dây điện bằng các xung sóng mang RF ngắn. Tốc độ và phạm vi thấp so với một số hệ PLC khác nhưng chi phí thấp; ứng dụng chính là bật/tắt, điều chỉnh đèn và điều khiển ổ cắm trong nhà thông minh qua bộ điều khiển hoặc ứng dụng. \sourcepage{20}

\begin{longtable}{p{.38\textwidth}p{.55\textwidth}}
\caption{Các thông số X10 theo bảng ở slide PDF trang 21.}\\\toprule
Thông số & Giá trị trong slide\\\midrule\endfirsthead
\toprule Thông số & Giá trị trong slide\\\midrule\endhead
Môi trường truyền & Đường dây điện (powerline only)\\
Loại thiết bị & Bộ điều khiển; thiết bị bị điều khiển; thiết bị có báo trạng thái\\
Điều chế lớp vật lý & OOK, sóng mang 120 kHz\\
Tốc độ dữ liệu & 120 bps tức thời; 60 bps duy trì; 24 bps tải hữu ích\\
Địa chỉ hóa & 16 mã nhà A--P và 16 mã thiết bị 1--16, tối đa 256 địa chỉ\\
Lệnh & 16 lệnh, như ON, OFF, DIM, BRIGHT, ALL UNITS OFF, ALL LIGHTS ON\\
Độ dài thông điệp & Slide ghi 100 bit (8,5 byte); hai trị số không tương đương nên cần kiểm tra nếu dùng thiết kế bản tin\\
Mã hóa & Manchester\\
Cấu trúc gói & Start Code 4 bit; House Code 8 bit; Key Code 10 bit (theo slide)\\
ACK và bảo mật & Không hỗ trợ theo bảng slide\\\bottomrule
\end{longtable}

\subsection{Giao tiếp khoảng cách rất ngắn: NFC và UWB}
\paragraph{NFC (Near-Field Communication).} Tập hợp giao thức vô tuyến cho các thiết bị tiếp xúc/đặt gần nhau, hướng tới xác thực, đọc/ghi thẻ, chia sẻ dữ liệu và kiểm soát ra vào; ví dụ thanh toán thông minh. Khác với cách dùng RFID chỉ để định danh trong nhiều hệ thống, NFC có các chế độ giao tiếp cho cả thiết bị và thẻ. \sourcepage{22}

\begin{longtable}{p{.47\textwidth}p{.45\textwidth}}
\caption{Thông số và chế độ NFC theo slide PDF trang 23.}\\\toprule
Đặc tính & Giá trị được nêu\\\midrule\endfirsthead
\toprule Đặc tính & Giá trị được nêu\\\midrule\endhead
Tần số mang & 13,56 MHz\\
Khoảng cách & Dưới 10 cm (bảng đối chiếu trang 42 sử dụng 4 cm)\\
Tốc độ & 106, 212 hoặc 424 kbit/s\\
Đối tượng & Thẻ (tag) và đầu đọc (reader)\\
Loại thẻ & Type 1 đến Type 5\\
Cấu hình mạng & Peer-to-peer\\
Chế độ thiết bị & Card emulation; reader/writer; peer-to-peer (active/passive)\\
Cách giao tiếp & Active--passive hoặc active--active\\
Thiết lập kết nối & Bảng slide ghi 'vài giây' và 'xấp xỉ dưới 0,1 giây'; cần xác định phép đo cụ thể\\
Mã hóa đường truyền & NRZ-L, Manchester, Modified Miller\\
Điều chế & ASK, BPSK (theo bảng)\\
Giải quyết va chạm & Có cơ chế chống va chạm\\
Tiêu chuẩn & ISO 14443A/B, ISO 18092, JIS X6319-4\\\bottomrule
\end{longtable}

\paragraph{UWB (Ultra-Wideband).} Vô tuyến sử dụng phổ tần rất rộng và xung phát năng lượng thấp để truyền dữ liệu và/hoặc định vị. Theo slide, khả năng hạn chế can nhiễu và theo dõi vị trí từng nút hỗ trợ mở rộng hệ thống; ứng dụng nổi bật là dịch vụ định vị và theo dõi tài sản. Không đồng nhất UWB với định vị GPS hay với mạng LPWAN: công năng và môi trường triển khai khác nhau. \sourcepage{24}

\subsection{Wi-Fi và Wi-Fi HaLow}
Wi-Fi phổ biến trong kết nối mạng cục bộ và có thể thiết lập kết nối ad hoc như Wi-Fi Direct mà không cần access point. Điểm bất lợi với cảm biến chạy pin là công suất tiêu thụ, dù nó không quan trọng bằng đối với nhiều thiết bị nhà thông minh có nguồn ổn định. \term{Wi-Fi HaLow}{IEEE 802.11ah} là hướng phát triển Wi-Fi tiết kiệm năng lượng dùng dải tần dưới 1 GHz, hướng tới thiết bị IoT. Bảng minh họa của slide nêu kênh OFDM băng hẹp, chế độ tiết kiệm năng lượng, IP gốc, bảo mật Wi-Fi, phạm vi khoảng 1 km, khả năng xuyên tường và hoạt động với pin nhỏ trong thời gian dài theo điều kiện ứng dụng. \sourcepage{25--26}

\subsection{Bluetooth Low Energy (BLE)}
BLE là kết nối vô tuyến tầm ngắn cho mạng cá nhân (PAN), chú trọng thiết bị cảm biến và wearable có nhu cầu trao đổi gói dữ liệu nhỏ, không liên tục. Từ Bluetooth 4.0, các cơ chế năng lượng thấp thích hợp để một nút hỗ trợ kết nối các nút cục bộ và chuyển dữ liệu về backend. Ứng dụng: tòa nhà thông minh, giao thông, thiết bị đeo. \sourcepage{27}

\begin{table}[htbp]\centering\small
\caption{Bảng Classic Bluetooth và BLE ở slide PDF trang 28 (số liệu theo phiên bản trong slide).}
\begin{tabularx}{\textwidth}{p{3.4cm}XX}\toprule
Thông số & Classic Bluetooth & BLE\\\midrule
Tầm hoạt động & 100 m & Hơn 100 m\\
Tốc độ dữ liệu & 1--3 Mb/s & 1 Mb/s\\
Thông lượng ứng dụng & 0,7--2,1 Mb/s & 0,27 Mb/s\\
Tần số & 2,4 GHz & 2,4 GHz\\
Bảo mật & 56/128 bit (bảng nguồn) & AES 128 bit với CCM\\
Độ trễ & 100 ms & 6 ms\\
Thời gian thiết lập & 100 ms & 3 ms\\
Âm thanh & Có & Không (bảng theo đời cũ)\\
Mô hình mạng & Sao & Sao\\
Công suất tiêu thụ (theo slide) & 1 W & 0,01--0,5 W\\
Dòng đỉnh (theo slide) & Dưới 30 mA & Dưới 15 mA\\\bottomrule
\end{tabularx}\end{table}

\begin{nguonslide}[title={Lưu ý với bảng Bluetooth}]
Bảng trang 28 phản ánh cách so sánh trong thời điểm soạn slide, không đại diện cho mọi phiên bản Bluetooth mới hơn. Đặc biệt mục \emph{Voice capable: No} áp dụng theo bảng lịch sử và không mô tả đầy đủ LE Audio. Nên dùng datasheet thiết bị cụ thể để thiết kế công suất.
\end{nguonslide}

\subsection{Zigbee}
Zigbee là công nghệ phù hợp mạng cảm biến/điều khiển công suất thấp với yêu cầu an toàn, độ bền kết nối và nhiều nút; được sử dụng trong ứng dụng công nghiệp và tự động hóa. Zigbee 3.0 được slide mô tả là sự thống nhất các chuẩn Zigbee trước đó. \sourcepage{29}

Các tính chất giao thức được liệt kê: hỗ trợ điểm--điểm, điểm--đa điểm và mesh; chu kỳ hoạt động (duty cycle) thấp giúp kéo dài tuổi thọ pin; độ trễ thấp; trải phổ chuỗi trực tiếp (DSSS); khả năng mở rộng tới khoảng 65.000 nút/mạng theo slide; mã hóa AES 128 bit; tránh xung đột, phát lại và xác nhận. Các giới hạn triển khai thực tế phụ thuộc cấu hình và thiết bị. \sourcepage{30}

\begin{figure}[htbp]\centering
\begin{tikzpicture}[node distance=12mm, every node/.style={font=\small}, coord/.style={circle,draw=orangeiot,fill=orange!25,minimum size=8mm},router/.style={circle,draw=teal,fill=cyan!20,minimum size=8mm},end/.style={circle,draw=navy,fill=blue!12,minimum size=7mm}]
\node[coord](c){C}; \node[router,below left=14mm and 22mm of c](r1){R};\node[router,below right=14mm and 22mm of c](r2){R};\node[router,below=30mm of c](r3){R};
\node[end,below left=13mm and 7mm of r1](a){E}; \node[end,left=17mm of r1](b){E}; \node[end,right=17mm of r2](d){E};\node[end,below right=13mm and 7mm of r2](e){E};
\draw[thick] (c)--(r1)--(r3)--(r2)--(c);\draw[thick] (r1)--(r2);\draw (r1)--(a) (r1)--(b) (r2)--(d) (r2)--(e);
\end{tikzpicture}
\caption{Ví dụ mạng Zigbee mesh. C: coordinator; R: router; E: end device (vẽ lại từ sơ đồ các topo slide 31).}\end{figure}

\subsection{Z-Wave}
Z-Wave được ứng dụng chủ yếu trong nhà thông minh để gắn thiết bị điện, điều khiển từ Internet và hỗ trợ an ninh/tự động hóa. Theo slide: tần số 868,42 MHz tại châu Âu, 908,42 MHz tại Hoa Kỳ; các tốc độ 9,6; 40; 200 kbit/s; điều chế GFSK ở dải gần 900 MHz; tầm khoảng 30 m trong nhà và 100 m ngoài trời. Thông số phụ thuộc khu vực pháp lý và phiên bản thiết bị. \sourcepage{32}

\subsection{Weightless, Sigfox và LoRaWAN}
\paragraph{Weightless.} Họ tiêu chuẩn mở LPWAN kết nối trạm gốc với hàng nghìn nút. Có Weightless-W, Weightless-P và Weightless-N cho các điều kiện khác nhau. Weightless-W hỗ trợ hai chiều ở phổ trắng truyền hình (TV white spaces), slide lấy ví dụ lĩnh vực dầu khí; Weightless-P cũng hai chiều, thích hợp mạng riêng và các trường hợp cần cả đường lên/đường xuống. Slide chỉ nêu tên Weightless-N mà không khai triển tính chất của biến thể này. \sourcepage{33}

\paragraph{Sigfox.} Dịch vụ kết nối LPWAN theo thuê bao tại các quốc gia được triển khai, tối ưu gói dữ liệu nhỏ gửi không thường xuyên. Ví dụ tiêu biểu: thiết bị đo thông minh hoặc cảm biến phát cảnh báo. \sourcepage{34}

\paragraph{LoRaWAN.} Công nghệ mạng diện rộng công suất thấp do hệ sinh thái LoRa Alliance thúc đẩy, cho truyền hai chiều, hướng tới chi phí thấp, pin lâu, mạng lớn và dịch vụ di động. Ứng dụng: thành phố thông minh và công nghiệp. Cần phân biệt điều chế vô tuyến LoRa ở lớp vật lý với giao thức mạng LoRaWAN. \sourcepage{35}

\subsection{Li-Fi và truyền thông ánh sáng}
Li-Fi (Light Fidelity) dùng ánh sáng nhìn thấy (Visible Light Communication, VLC) để truyền dữ liệu tốc độ cao, có thể tận dụng đèn LED hiệu suất cao. Điểm hạn chế trong mô tả slide là yêu cầu điều kiện truyền quang phù hợp, dễ bị cản bởi vật chắn hoặc mất đường nhìn; chưa được triển khai rộng rãi trong IoT tại thời điểm biên soạn. Tiềm năng: lớp học và phòng trình chiếu thông minh. \sourcepage{36}

\subsection{Mạng di động và 5G}
Các thế hệ mạng tế bào (3G, 4G, LTE-Advanced, 5G) có vùng phủ rộng và khả năng hỗ trợ thiết bị di động. LTE-Advanced được slide lấy làm ví dụ về tiến bộ tốc độ và chế độ tiết kiệm năng lượng: có biến thể/chế độ cho các nút có thể đạt thời gian pin rất dài tùy chu kỳ gửi dữ liệu; bản LTE-A khác có tốc độ tối đa ở mức rất cao. Hai trị số 10 Mb/s và 1 Gb/s trong slide mô tả các trường hợp khác nhau, không phải một cấu hình chung. \sourcepage{37}

Slide 5G (có trích dẫn dự báo trước năm 2020) nêu tăng tốc độ, độ phủ, số người dùng và mức cơ động so với 4G. Bảng so sánh 4G/5G ở slide PDF trang 39 có các ý: dải tần 4G khoảng 600 MHz--5,925 GHz; 5G gồm FR1 (khoảng 450 MHz--6 GHz) và FR2 (khoảng 24,25--52,6 GHz, tùy quy định); băng thông 4G tới 20 MHz, 5G FR1 tới 100 MHz/FR2 tới 400 MHz; khoảng cách sóng mang 4G cố định 15 kHz, 5G mở rộng; tốc độ đỉnh downlink theo slide là 1 Gb/s và 20 Gb/s; độ trễ slide lần lượt 10 ms và dưới 1 ms; quy mô kết nối 2.000 so với 1 triệu thiết bị/km$^2$; tốc độ di động 350 so với 500 km/h. Đây là mục tiêu/giá trị minh họa trong bảng, không phải hiệu năng bảo đảm mọi vùng phủ. \sourcepage{38--39}

\subsection{LTE-M và NB-IoT}
\paragraph{LTE-M.} LTE-M, tên đầy đủ LTE Machine-Type Communication, gồm eMTC (enhanced MTC), là công nghệ kết nối di động LPWAN do 3GPP thiết kế cho thiết bị M2M/IoT. eMTC LTE Cat M1 đã được đặc tả ở 3GPP Release 13 (tháng 6 năm 2016 theo slide). Hệ 3GPP còn có NB-IoT và EC-GSM-IoT. \sourcepage{40}

\paragraph{NB-IoT.} Narrowband IoT là tiêu chuẩn LPWAN cho các thiết bị truyền dữ liệu ít, chú trọng mức tiêu thụ điện thấp, dung lượng mạng, hiệu quả phổ và vùng phủ sâu (ví dụ trong công trình). Slide cho rằng thời gian dùng pin hơn 10 năm khả thi ở một số kịch bản; điều này phụ thuộc dữ liệu gửi, vùng phủ, pin và cấu hình hoạt động. \sourcepage{41}

\begin{table}[htbp]\centering\small
\caption{So sánh khái quát hai công nghệ LTE-M và NB-IoT từ slide PDF trang 50.}
\begin{tabularx}{\textwidth}{p{3.4cm}XX}\toprule
Đặc tính & LTE-M & NB-IoT\\\midrule
Tên chuẩn & LTE Cat M1 (theo slide) & LTE Cat NB1\\
Băng thông & Khoảng 1,08 MHz, trong kênh LTE 1,4 MHz & Tới 180 kHz\\
Thông lượng & 36 kb/s (giá trị trong bảng slide) & 30/60 kb/s (giá trị trong bảng slide)\\
Triển khai & Trong các mạng LTE hiện hữu & In-band, guard-band hoặc độc lập trong một số phổ GSM\\
Khả năng giọng nói & Có hỗ trợ thoại/VoLTE theo cấu hình & Không được thiết kế cho thoại\\
Di động/chuyển giao & Có & Hạn chế\\
Thời lượng pin theo bảng & Có thể tới 10 năm & Có thể tới 10 năm\\\bottomrule
\end{tabularx}\end{table}

\subsection{Bảng tổng hợp phổ, tầm phủ và tốc độ các công nghệ}
\begingroup\small\sloppy\begin{longtable}{p{2.6cm}p{3.5cm}p{2.25cm}p{3.6cm}}
\caption{Bảng đối chiếu slide PDF trang 42; số liệu tham khảo theo slide, không đại diện mọi phiên bản.}\\
\toprule Công nghệ & Tần số theo slide & Tầm phủ & Tốc độ theo slide\\\midrule\endfirsthead
\toprule Công nghệ & Tần số theo slide & Tầm phủ & Tốc độ\\\midrule\endhead
Bluetooth & 2,4 GHz & 100 m & 125 kbit/s--Mb/s\\
Wi-Fi & 2,4 hoặc 5 GHz & 50 m & 150--600 Mb/s\\
NFC & 13,56 MHz & 4 cm & 100--420 kbit/s\\
LoRaWAN & 867--869 MHz (Âu); 902--928 MHz (Bắc Mỹ) & 15 km & 0,3--50 kbit/s\\
Cellular & 900/1800/\allowbreak 1900/2100 MHz & Bảng nguồn ghi 30 m (nút--trạm), không phải tầm phủ mạng tế bào nói chung & 21 Mb/s (3G+); 600 Mb/s (4G)\\
Z-Wave & 865--926 MHz (ISM) & 100 m & 100 kbit/s\\
Zigbee & 2,4 GHz (ISM) & 100 m & 20--250 kbit/s\\
Sigfox & 900 MHz & 3--50 km & 10--1000 bit/s\\\bottomrule
\end{longtable}\endgroup

\subsection{Bảo mật ở lớp truyền thông}
Các giao thức được slide nêu bao gồm:
\begin{itemize}
\item \textbf{IPsec (Internet Protocol Security):} cơ chế bảo mật ở lớp mạng IP, cung cấp các phương thức xác thực, toàn vẹn và mã hóa khi được triển khai/cấu hình đúng. Được sử dụng với IPv6 nhưng không nên hiểu IPsec mặc nhiên được bật trên mọi mạng IPv6.
\item \textbf{TLS (Transport Layer Security):} bảo vệ kênh kết nối dựa trên TCP, ví dụ HTTPS hoặc MQTT qua TLS.
\item \textbf{DTLS (Datagram Transport Layer Security):} giải pháp bảo mật cho truyền thông kiểu datagram như UDP, ví dụ CoAP/DTLS.
\item \textbf{IEEE 1888.3:} thành viên họ tiêu chuẩn dành cho mạng điều khiển cộng đồng xanh phổ cập (Ubiquitous Green Community Control Network). Mục đích theo slide là hạn chế truy cập trái phép, ngăn rò rỉ dữ liệu và bảo đảm tính bí mật/toàn vẹn dữ liệu truyền.
\end{itemize}
Các thuộc tính cốt lõi là bảo mật nội dung, xác thực đối tác và phát hiện sửa đổi; hệ thống IoT vẫn phải quản lý khóa, phân quyền nút và cập nhật phần mềm. \sourcepage{43}

\subsection{Định tuyến và giao thức RPL}
Định tuyến (routing) là quá trình chọn đường chuyển lưu lượng trong một mạng hoặc qua nhiều mạng. Với rất nhiều thiết bị có tài nguyên hạn chế, chạy pin và liên kết chất lượng không ổn định, lựa chọn đường tốt có ảnh hưởng lớn tới khả năng mở rộng, thời gian gửi dữ liệu và thời lượng pin.

\term{RPL}{Routing Protocol for Low-Power and Lossy Networks} là chuẩn định tuyến IPv6 dành cho mạng công suất thấp và liên kết dễ mất gói. Slide nêu có thể sử dụng trên IEEE 802.15.4, hỗ trợ các kiểu truyền một--một, nhiều--một và một--nhiều. Hình trang 45 biểu diễn hai cấu hình \emph{DODAG} (Destination-Oriented Directed Acyclic Graph): non-storing và storing mode; các nút tổ chức tương quan với nút gốc, còn phương thức duy trì đường đi phụ thuộc chế độ hoạt động. \sourcepage{44--45}

\begin{figure}[htbp]\centering
\begin{tikzpicture}[>=Latex, every node/.style={font=\small}, leaf/.style={circle,draw=orangeiot,fill=orange!18,minimum size=8mm}, root/.style={circle,draw=navy,fill=blue!20,minimum size=10mm}]
\node[root](r) at (0,3){Gốc};
\node[leaf](a) at (-2,1.7){A};\node[leaf](b) at (0,1.7){B};\node[leaf](c) at (2,1.7){C};
\node[leaf](d) at (-2.6,0){D};\node[leaf](e) at (-.9,0){E};\node[leaf](f) at (1,0){F};\node[leaf](g) at (2.7,0){G};
\foreach \x/\y in {a/r,b/r,c/r,d/a,e/a,e/b,f/b,f/c,g/c}{\draw[->,teal,thick](\x)--(\y);}
\end{tikzpicture}
\caption{Một đồ thị định tuyến hướng về nút gốc kiểu DODAG của RPL (minh họa vẽ lại từ slide trang 45).}\end{figure}

\section{Các loại mạng IoT}
Việc kết nối tin cậy giữa cảm biến, thiết bị và nền tảng IoT quyết định hoạt động của ứng dụng. Slide chia các mạng đang dùng thành bốn nhóm: mạng tế bào, mạng cục bộ/cá nhân (LAN/PAN), mạng diện rộng công suất thấp (LPWAN), và mạng mắt lưới (mesh). Đây là các góc nhìn phân loại theo phạm vi/vận hành và topo; một hệ thống thực tế có thể kết hợp chúng. \sourcepage{47}

\subsection{Mạng tế bào (Cellular Networks)}
Mạng tế bào sử dụng cơ sở hạ tầng nhà khai thác di động giống điện thoại thông minh để kết nối thiết bị IoT. Ưu điểm là vùng phủ địa lý rộng và hỗ trợ tính di động; hạn chế là điện năng và phí dịch vụ, cũng như khả năng phủ sóng kém ở hầm sâu, tủ kỹ thuật kín, mỏ hoặc công trình nhiều lớp vật cản. LTE-M và NB-IoT là hai lựa chọn IoT di động được đề cập; slide cho rằng NB-IoT có thể tiêu thụ ít năng lượng và giá thành thấp hơn LTE-M, nhưng phạm vi hỗ trợ dịch vụ phụ thuộc mạng triển khai ở thời điểm khảo sát. \sourcepage{48--50}

\subsection{Mạng cục bộ và cá nhân (LAN/PAN)}
PAN (Personal Area Network) dùng cho kết nối các thiết bị gần cá nhân; LAN (Local Area Network) kết nối các nút trong nhà, khu vực hoặc cơ sở. Phạm vi ngắn, chi phí sử dụng thường thấp; độ tin cậy phụ thuộc nhiễu, vật cản, mật độ nút và cấu hình. Hai ví dụ chính trong slide là Wi-Fi và Bluetooth. Có thể dùng gateway chuyển dữ liệu của mạng LAN/PAN ra mạng diện rộng. \sourcepage{51--52}

\subsection{Mạng diện rộng công suất thấp (LPWAN)}
LPWAN cho tầm xa hơn các kết nối Wi-Fi/Bluetooth thông dụng nhưng thường dùng ít năng lượng hơn nhiều công nghệ di động băng rộng. Đổi lại, tốc độ thấp và phù hợp các bản tin cảm biến nhỏ, gửi cách quãng. LoRaWAN, Sigfox, Weightless và các biến thể LPWAN tế bào như LTE-M/NB-IoT là các ví dụ đã giới thiệu trong phần trên. \sourcepage{53}

\subsection{Mạng mắt lưới (Mesh Networks)}
Trong mạng mesh, các nút có thể phối hợp chuyển tiếp dữ liệu cho nhau để đến gateway. Tầm mỗi bước nhảy có thể ngắn nhưng tổng vùng phủ được kéo dài qua nhiều nút trung gian. Điều này hữu ích cho cảm biến phân tán trong một tòa nhà; đổi lại cần đủ nút, bộ chuyển tiếp và cấu hình đường đi để tránh các vùng không phủ. Sơ đồ Zigbee và RPL ở các phần trước là những minh họa thích hợp cho kiến trúc đa chặng. \sourcepage{54}

\begin{figure}[htbp]\centering
\begin{tikzpicture}[>=Latex,scale=1,transform shape]
\draw[->,thick] (0,0)--(12,0) node[right]{Tầm phủ};
\draw[->,thick] (0,0)--(0,5) node[above,align=center]{Tốc độ dữ liệu};
\fill[blue!12] (.6,2.8) rectangle (3.3,4.6);\node[align=center] at (1.95,3.8){Wi-Fi\\Bluetooth};
\fill[cyan!16] (7,.65) rectangle (10.15,2.45);\node[align=center] at (8.55,1.55){LPWAN\\LoRa, Sigfox, NB-IoT};
\fill[orange!15] (4.0,2.5) rectangle (10.7,4.65);\node[align=center] at (7.45,3.65){Mạng tế bào\\2G/3G/4G/5G};
\end{tikzpicture}
\caption{So sánh định tính tầm phủ--tốc độ theo bố cục slide trang 55; vị trí vùng chỉ mang tính mô tả, không theo thang số.}\end{figure}

\begin{table}[htbp]\centering\small
\caption{Tóm tắt bốn nhóm mạng IoT.}
\begin{tabularx}{\textwidth}{p{2.6cm}XXX}\toprule
Loại mạng & Điểm mạnh & Hạn chế & Ví dụ\\\midrule
Cellular & Phủ rộng, di động & Năng lượng/chi phí; vùng phủ trong nhà sâu & 4G, LTE-M, NB-IoT, 5G\\
LAN/PAN & Chi phí triển khai thấp, tốc độ tùy công nghệ & Tầm ngắn, dễ bị cản/nhiễu & Wi-Fi, BLE\\
LPWAN & Tầm xa, phù hợp cảm biến chạy pin & Thông lượng thấp, hạn chế payload/chu kỳ & LoRaWAN, Sigfox, Weightless\\
Mesh & Mở rộng đa chặng, có tuyến thay thế & Cần nút chuyển tiếp, phức tạp định tuyến & Zigbee mesh, RPL\\\bottomrule
\end{tabularx}\end{table}

\subsection{Sơ đồ phân nhóm các công nghệ mạng}
Slide trang 46 sắp xếp mạng theo \emph{cellular}, \emph{wired}, \emph{satellite}, \emph{short-range wireless}, và \emph{LP-WAN / LP wireless}. Các mục được nêu gồm: 2G/2.5G, 3G, 4G/LTE, 5G; POTS, Fiber, Ethernet, ATM/Frame Relay, MPLS, Powerline; C-Band và L-Band; Wi-Fi/802.11, Bluetooth/BLE, Zigbee/802.15.4, 900 MHz proprietary, 6LoWPAN, Wi-Fi HaLow; WiMAX, Sigfox, LoRa, NB-IoT, LTE-M, Fixed Microwave và WiSUN. Hình này bổ sung các loại truyền dẫn chứ không thay đổi cách phân nhóm bốn loại của slide 47.

\section{Giao thức dữ liệu của IoT}
\subsection{Vai trò của giao thức ứng dụng và mô hình trao đổi}
Slide bắt đầu bằng việc phân biệt truyền thông dữ liệu giữa các thiết bị với khả năng truy cập Internet, nhấn mạnh các nút công suất thấp có thể giao tiếp trực tiếp bằng liên kết có dây hoặc vô tuyến mà không phải lúc nào cũng cần truy cập Internet. Trong ngăn xếp giao thức, Wi-Fi/BLE/LoRa thuộc nhóm công nghệ kết nối bên dưới, còn HTTP, CoAP, MQTT, WebSocket, AMQP và XMPP cung cấp cách thức trao đổi thông điệp ở tầng ứng dụng. \sourcepage{56,60}

Một nút IoT có thể thu thập dữ liệu, phân tích dữ liệu và điều khiển nút khác. Tùy hệ thống, trước khi tương tác với dịch vụ, nút cần được nhận dạng và đăng ký. Slide gọi quá trình đó là \emph{node subscription}. Trong ngữ cảnh truyền dữ liệu, cần phân biệt \emph{đăng ký một nút với dịch vụ} và \emph{đăng ký theo dõi một chủ đề dữ liệu (subscribe topic)}. \sourcepage{57}

\paragraph{Publish/subscribe (xuất bản/đăng ký).} Nút xuất bản (publisher) gửi bản tin lên một chủ đề/kênh; các nút đăng ký (subscriber) nhận dữ liệu tương ứng khi bản tin được xuất bản. Thành phần trung gian có thể là broker; cách tiếp cận tách người gửi và người nhận và thuận lợi cho nhiều nút trao đổi thường xuyên. \sourcepage{58--59}

\paragraph{Request/response (yêu cầu/phản hồi).} Client chủ động yêu cầu dữ liệu hoặc tác vụ và server trả kết quả. Mô hình quen thuộc với ứng dụng client/server và các dịch vụ REST. Slide nhận xét publish/subscribe thường thuận lợi hơn với triển khai IoT lớn trao đổi dữ liệu nhiều, trong khi request/response dễ tích hợp ứng dụng client/server. \sourcepage{59}

\begin{figure}[htbp]\centering
\begin{tikzpicture}[>=Latex,font=\small,block/.style={draw=navy,rounded corners,minimum width=29mm,minimum height=10mm,fill=blue!6}]
\node[block](cli) at (0,2.7){Client};\node[block](serv) at (5,2.7){Server};
\draw[->,thick,teal] (cli)--node[above]{Request}(serv);
\draw[<-,thick,orangeiot] (cli) to[bend right=22] node[below]{Response}(serv);
\node[block](pub) at (0,0){Publisher};\node[block](broker) at (5,0){Broker};\node[block](sub) at (10,0){Subscriber};
\draw[->,thick,teal](pub)--node[above]{Publish}(broker);
\draw[->,thick,teal](broker)--node[above]{Push chủ đề}(sub);
\end{tikzpicture}
\caption{Hai mô hình trao đổi dữ liệu ở slide trang 57--59.}\end{figure}

\subsection{HTTP: Hypertext Transfer Protocol}
HTTP là giao thức ứng dụng phổ biến của Web, thường theo mô hình \emph{pull}: client gửi yêu cầu lấy tài nguyên, server trả phản hồi. HTTP không lưu trạng thái yêu cầu trước đó ở bản thân giao thức (stateless), hỗ trợ thương lượng kiểu nội dung và biểu diễn tài nguyên để các hệ thống được phát triển độc lập vẫn trao đổi dữ liệu. \sourcepage{61--62}

Những phương thức HTTP phổ biến: \texttt{GET} (đọc tài nguyên), \texttt{POST} (tạo/xử lý dữ liệu), \texttt{PUT} (thay thế/cập nhật tài nguyên), \texttt{PATCH} (cập nhật một phần), \texttt{DELETE} (xóa). CRUD: Create--Read--Update--Delete là các thao tác ứng dụng tương ứng ở mức khái niệm, không phải ánh xạ duy nhất cho mọi API. \sourcepage{62}

Theo slide, HTTP có hạn chế đối với một số hệ thống IoT: quá nhiều cặp yêu cầu--phản hồi khi nhiều cảm biến hoạt động khiến phát sinh độ trễ và năng lượng; thông thường client khởi xướng yêu cầu thay vì server tự đẩy bản tin như mô hình publish/subscribe; HTTP/TCP có overhead và chi phí kết nối, không tối ưu cho một số cảm biến dùng pin. Những hạn chế này được mô tả cho cách triển khai HTTP phổ biến; các biến thể HTTP mới, keep-alive, server push hoặc kỹ thuật khác có thể thay đổi chi phí thực tế. \sourcepage{63--64}

\begin{figure}[htbp]\centering
\begin{tikzpicture}[>=Latex,font=\small]
\node at (0,4.8){Client};\node at (6,4.8){Server};
\draw[dashed](0,4.4)--(0,0);\draw[dashed](6,4.4)--(6,0);
\draw[->,thick,navy](0,4)--node[above]{TCP: SYN}(6,3.5);
\draw[<-,thick,navy](0,2.9)--node[above]{SYN/ACK}(6,3.4);
\draw[->,thick,navy](0,2.8)--node[above]{ACK}(6,2.4);
\draw[->,thick,teal](0,2)--node[above]{HTTP GET /temperature}(6,1.6);
\draw[<-,thick,orangeiot](0,.8)--node[above]{HTTP 200 OK + dữ liệu}(6,1.3);
\end{tikzpicture}
\caption{Chuỗi bắt tay TCP và yêu cầu--phản hồi HTTP, dựng theo ý tưởng hình slide trang 63.}\end{figure}

\subsection{TCP và UDP}
\textbf{TCP} là giao thức hướng kết nối, sử dụng bắt tay thiết lập phiên, đánh số thứ tự, xác nhận/điều khiển truyền lại để cung cấp luồng byte tin cậy. \textbf{UDP} truyền datagram không thiết lập phiên TCP và không tự đảm bảo xác nhận, thứ tự hoặc truyền lại. Do đó TCP thường có chi phí thiết lập/trạng thái cao hơn, còn UDP thích hợp với giao thức tự quản lý độ tin cậy hoặc yêu cầu thông điệp ngắn như CoAP. Slide trang 65 minh họa TCP bằng SYN--SYN/ACK--ACK và UDP bằng các datagram yêu cầu--phản hồi không có bắt tay. Việc nói 'UDP không tin cậy' chỉ có nghĩa UDP không tự bảo đảm chuyển phát, không có nghĩa ứng dụng UDP không thể cài đặt xác nhận. \sourcepage{65}

\subsection{CoAP: Constrained Application Protocol}
CoAP là giao thức truyền tài nguyên kiểu Web thiết kế cho nút và mạng bị hạn chế về băng thông, độ sẵn sàng và tài nguyên xử lý. Nó sử dụng mô hình client/server tương tự REST/HTTP: server công bố tài nguyên qua URI, client dùng GET, POST, PUT, DELETE. Bản tin CoAP được thiết kế gọn và thường dựa trên UDP; việc xác nhận chuyển phát phải do cơ chế CoAP xử lý khi cần. \sourcepage{66}

\begin{figure}[htbp]\centering
\begin{tikzpicture}[>=Latex,font=\small,box/.style={draw=navy,fill=blue!5,rounded corners,minimum width=23mm,minimum height=10mm}]
\node[box](web) at (0,0){HTTP client};\node[box](proxy) at (4.2,0){HTTP/CoAP proxy};\node[box](sensor) at (8.8,0){Cảm biến CoAP};
\draw[<->,thick,navy](web)--node[above]{HTTP / TCP}(proxy);
\draw[<->,thick,teal](proxy)--node[above]{CoAP / UDP}(sensor);
\end{tikzpicture}
\caption{Liên kết client Web với mạng cảm biến giới hạn qua proxy CoAP (theo slide 67).}\end{figure}

\paragraph{Đặc trưng CoAP.} Slide nêu năm nhóm ưu điểm: phục vụ giao tiếp M2M bị ràng buộc; trao đổi bản tin không đồng bộ; overhead thấp và dễ phân tích cú pháp; hỗ trợ URI và kiểu nội dung; hỗ trợ proxy và cache. CoAP có thể làm cầu nối với Web bằng gateway/proxy. \sourcepage{67--68}

\paragraph{Bốn loại bản tin.} CoAP định nghĩa:
\begin{enumerate}[leftmargin=2em]
\item \textbf{Confirmable (CON):} bản tin cần ACK; nếu không nhận được sẽ có cơ chế gửi lại.
\item \textbf{Non-confirmable (NON):} bản tin không đòi ACK ở lớp thông điệp CoAP (\emph{fire-and-forget}).
\item \textbf{Acknowledgement (ACK):} xác nhận đã nhận bản tin CON.
\item \textbf{Reset (RST):} từ chối/đặt lại bản tin hoặc xử lý quan sát không còn hợp lệ.
\end{enumerate}

CoAP dùng Message ID 16 bit cho ghép bản tin và phát hiện bản sao, cùng cơ chế truyền lại có back-off đối với CON. Bảo mật có thể sử dụng DTLS. \sourcepage{69}

\begin{figure}[htbp]\centering
\begin{tikzpicture}[>=Latex,font=\small]
\node at (0,4.8){Client};\node at (7,4.8){Server};
\draw[dashed](0,4.4)--(0,.1);\draw[dashed](7,4.4)--(7,.1);
\draw[->,thick,teal](0,3.9)--node[above]{CON, MID=0x4A51, GET}(7,3.35);
\draw[<-,thick,orangeiot](0,2.3)--node[above]{ACK, MID=0x4A51, 2.05 Content}(7,2.9);
\draw[->,thick,teal](0,1.7)--node[above]{NON, MID=0x4A52}(7,1.2);
\node[align=left] at (3.5,.45){ACK cho CON; không ACK cho NON.};
\end{tikzpicture}
\caption{Ví dụ luồng CoAP confirmable và non-confirmable, vẽ lại theo slide trang 70--71.}\end{figure}

\paragraph{Ví dụ tra cứu tài nguyên.} Các hình slide 70--71 cho thấy bản tin \texttt{CON GET /temperature} với \texttt{Token} được phản hồi bằng \texttt{ACK 2.05 Content} chứa nhiệt độ \texttt{20.1 C}, hoặc phản hồi lỗi \texttt{4.04 Not Found} khi tài nguyên không tồn tại. Hình trang 70 cũng thể hiện CON--ACK, CON--RST và NON một chiều, để làm rõ khác nhau giữa xác nhận và từ chối. \sourcepage{70--71}

\subsection{MQTT: Message Queuing Telemetry Transport}
MQTT là giao thức gửi thông điệp nhẹ, hữu ích để phân phối dữ liệu telemetry tới các client có tài nguyên và kết nối mạng hạn chế. Nó dùng mô hình publish/subscribe và rất phổ biến trong giao tiếp M2M/IoT. Thành phần trung tâm là \textbf{broker}: publisher gửi dữ liệu theo chủ đề (topic); broker xác định người đăng ký để chuyển tiếp; subscriber nhận bản tin không cần duy trì kết nối trực tiếp với từng publisher. \sourcepage{72--73}

\begin{figure}[htbp]\centering
\begin{tikzpicture}[>=Latex,font=\small,box/.style={draw=navy,rounded corners,fill=blue!5,minimum height=12mm,minimum width=28mm}]
\node[box](p) at (0,0){Publisher};\node[box,fill=cyan!15](b) at (4.5,0){MQTT Broker};
\node[box](s1) at (9,1.4){Subscriber 1};\node[box](s2) at (9,-1.4){Subscriber 2};
\draw[->,thick,teal](p)--node[above]{PUBLISH topic}(b);
\draw[->,thick,orangeiot](b)--node[above,sloped]{PUBLISH}(s1);
\draw[->,thick,orangeiot](b)--node[below,sloped]{PUBLISH}(s2);
\end{tikzpicture}
\caption{Kiến trúc MQTT publisher--broker--subscriber (dựng lại từ hình slide trang 73).}\end{figure}

\paragraph{Bốn giai đoạn phiên MQTT.} Theo slide: (1) \textbf{Connection}: client thiết lập kết nối TCP/IP đến broker (cổng chuẩn hoặc cổng riêng do server định nghĩa); (2) \textbf{Authentication}: xác thực nếu broker yêu cầu; (3) \textbf{Communication}: publish và subscribe bản tin; (4) \textbf{Termination}: đóng phiên/kết nối. Với Client ID và thiết lập phiên thích hợp, broker có thể tiếp tục một phiên trước đó. \sourcepage{74}

\paragraph{Ba mức Quality of Service (QoS).}
\begin{table}[htbp]\centering\small
\caption{Ba mức QoS trong MQTT (slide trang 75).}
\begin{tabularx}{\textwidth}{p{1.4cm}p{3cm}X}\toprule
QoS & Ý nghĩa & Đặc điểm\\\midrule
0 & At most once (tối đa một lần) & Không bảo đảm thông điệp đến đích; không có cơ chế xác nhận chuyển phát QoS 0\\
1 & At least once (ít nhất một lần) & Bảo đảm theo cơ chế xác nhận/gửi lại; có thể trùng lặp bản tin\\
2 & Exactly once (đúng một lần) & Có quy trình bắt tay để bảo đảm xử lý chuyển phát đúng một lần theo ngữ nghĩa MQTT\\\bottomrule
\end{tabularx}\end{table}

\paragraph{Ví dụ ứng dụng ở thời điểm giảng dạy.} Slide nêu Facebook sử dụng MQTT trong hệ thống nhắn tin để tiết kiệm pin và truyền tin nhanh trên mạng không ổn định; đồng thời nêu các nền tảng AWS, Google Cloud, IBM Bluemix, Microsoft Azure và ThingOn IoT Platform có hỗ trợ MQTT. Những khẳng định mang yếu tố thời điểm này được giữ theo tư liệu giảng dạy, không coi là xác nhận về sản phẩm hiện tại. \sourcepage{76}

\subsection{WebSocket}
WebSocket liên quan HTTP trong bước thiết lập kết nối: trình duyệt/client thực hiện cơ chế bắt tay nâng cấp kết nối và duy trì một kênh truyền trên TCP để cả hai phía trao đổi bản tin cho đến khi đóng. So với cơ chế nhiều lượt request/response, WebSocket cho phép truyền hai chiều trên kết nối lâu dài, thuận lợi cho ứng dụng thời gian thực; slide nhận xét nó vẫn có overhead đáng kể đối với các nút IoT rất hạn chế tài nguyên. Hình trang 78 minh họa cột HTTP với những cặp request--response lặp lại so với cột WebSocket với bước handshake rồi truyền hai chiều trực tiếp. \sourcepage{77--78}

\begin{figure}[htbp]\centering
\begin{tikzpicture}[>=Latex,font=\small]
\node at (0,3.5){Client};\node at (6,3.5){Server};\draw[dashed](0,3.2)--(0,0);\draw[dashed](6,3.2)--(6,0);
\draw[->,thick,navy](0,2.9)--node[above]{HTTP Upgrade: websocket}(6,2.65);
\draw[<-,thick,navy](0,2.22)--node[above]{101 Switching Protocols}(6,2.5);
\draw[->,thick,teal](0,1.6)--node[above]{Thông điệp client $\to$ server}(6,1.4);
\draw[<-,thick,orangeiot](0,.8)--node[above]{Thông điệp server $\to$ client}(6,1.0);
\end{tikzpicture}
\caption{Thiết lập và truyền hai chiều với WebSocket, vẽ lại theo slide 77--78.}\end{figure}

\subsection{AMQP: Advanced Message Queuing Protocol}
AMQP là tiêu chuẩn mở thiên về giao tiếp hướng thông điệp, hỗ trợ publish/subscribe. Tương tự cách mô tả MQTT trong slide, các ứng dụng có thể sử dụng các ngữ nghĩa chuyển phát tối đa một lần, ít nhất một lần hoặc đúng một lần tùy giao thức/triển khai. Kênh TCP có thể được bảo vệ với TLS/SSL, đồng thời có cơ chế xác thực SASL. Hai thành phần cốt lõi được slide nêu:
\begin{itemize}
\item \textbf{Exchange:} nhận bản tin và chọn (định tuyến) queue phù hợp dựa trên các quy tắc đã đặt.
\item \textbf{Message queue:} lưu bản tin chờ được chuyển tới phía tiêu thụ/nhận.
\end{itemize}
Kiến trúc này phù hợp hệ thống có xử lý hàng đợi, định tuyến thông điệp linh hoạt và tích hợp backend. \sourcepage{79}

\begin{figure}[htbp]\centering
\begin{tikzpicture}[>=Latex,font=\small,box/.style={draw=navy,rounded corners,minimum width=26mm,minimum height=10mm,fill=blue!5}]
\node[box](p) at (0,0){Producer};\node[box](ex) at (4,0){Exchange};\node[box](q) at (8,0){Queue};\node[box](c) at (12,0){Consumer};
\draw[->,thick,teal](p)--(ex);\draw[->,thick,teal](ex)--node[above]{Routing}(q);\draw[->,thick,teal](q)--(c);
\end{tikzpicture}
\caption{Luồng thông điệp cơ bản AMQP theo slide trang 79.}\end{figure}

\subsection{XMPP: Extensible Messaging and Presence Protocol}
XMPP là giao thức tiêu chuẩn hóa qua IETF để trao đổi thông tin trên Internet, hoạt động gần thời gian thực, độ trễ thấp theo bối cảnh ứng dụng và có tính độc lập nền tảng/phân tán. Theo slide, XMPP dựa trên TCP, hỗ trợ cả publish/subscribe và request/response; client kết nối tới server và trao đổi bản tin dựa trên \term{XML}{Extensible Markup Language}. Nhờ cấu trúc thông điệp mở rộng, nó có thể phục vụ ứng dụng liên lạc IoT, nhưng XML thường tạo thêm overhead so với các định dạng nhị phân nhỏ gọn. \sourcepage{80}

\subsection{Bảng so sánh các giao thức dữ liệu}
\begingroup\footnotesize\sloppy\begin{longtable}{p{1.8cm}p{3.45cm}p{1.8cm}p{4.15cm}}
\caption{Tổng hợp từ các slide trang 56--80; một số nhận xét là phân tích hệ thống hóa.}\\
\toprule Giao thức & Mô hình & Vận chuyển & Ghi chú trong IoT\\\midrule\endfirsthead
\toprule Giao thức & Mô hình & Vận chuyển & Ghi chú trong IoT\\\midrule\endhead
HTTP & Request/response, client/server & TCP & Phổ biến REST, stateless; overhead kết nối/yêu cầu có thể cao\\
CoAP & REST, request/response; có hỗ trợ quan sát & UDP & Nhẹ, CON/NON/ACK/RST, URI, proxy/cache, DTLS\\
MQTT & Publish/subscribe qua broker & TCP & Chủ đề, phiên, QoS 0/1/2\\
WebSocket & Kết nối song công dài hạn & TCP & Nâng cấp từ HTTP; trao đổi hai chiều liên tục\\
AMQP & Trao đổi thông điệp, hàng đợi & TCP & Exchange, queue, routing, TLS/SASL\\
XMPP & Cả publish/subscribe và request/response & TCP & Bản tin XML, mở rộng, gần thời gian thực\\\bottomrule
\end{longtable}\endgroup

\section{Tổng kết kiến thức và gợi ý ứng dụng}
\begin{ghinho}[title={Những khái niệm cần phân biệt}]
\textbf{Định danh/địa chỉ:} URI, EPC, uCode, IPv4, IPv6.\par
\textbf{Kết nối vật lý và mạng:} PLC/X10, NFC, UWB, Wi-Fi/HaLow, BLE, Zigbee, Z-Wave, Weightless, Sigfox, LoRaWAN, Li-Fi, 3G/4G/5G, LTE-M, NB-IoT.\par
\textbf{Bảo vệ, định tuyến:} IPsec, TLS/DTLS, IEEE 1888.3, RPL.\par
\textbf{Kiến trúc mạng:} cellular, LAN/PAN, LPWAN, mesh.\par
\textbf{Cách trao đổi dữ liệu:} request/response và publish/subscribe; HTTP, CoAP, MQTT, WebSocket, AMQP, XMPP.
\end{ghinho}

Một ví dụ tổng hợp: cảm biến chạy pin trong nhà có thể kết nối BLE hoặc Zigbee tới gateway; gateway kết nối Wi-Fi hoặc Ethernet lên máy chủ; dữ liệu được xuất bản qua MQTT để nhiều dịch vụ nhận cùng thời điểm. Một cảm biến cần REST tối giản có thể dùng CoAP/UDP; hệ thống Web có thể dùng HTTP để truy vấn và WebSocket để cập nhật liên tục. Nút ở vị trí xa không có Wi-Fi có thể cân nhắc LoRaWAN hoặc LTE-M/NB-IoT. Các lựa chọn này là \emph{minh họa thiết kế được suy ra từ bài học}, không phải kịch bản thí nghiệm bắt buộc trong slide.

\section{Đối chiếu toàn bộ 81 trang slide}
Bảng dưới đây xác định trang tài liệu gốc liên quan tới từng mục, để người đọc rà soát các hình và bảng của bài giảng. Trang 1 là bìa; trang 2 là mục lục; trang 81 là lời cảm ơn. Các trang mang hình không trích xuất được thành chữ được mô tả và minh họa lại trong phần tương ứng.
\begin{longtable}{p{3.0cm}p{10cm}}
\toprule Trang PDF & Nội dung có trong tài liệu/chương\\\midrule\endhead
1--2 & Bìa, học phần, giảng viên, mục lục\\
3--8 & Định nghĩa giao tiếp, H2M, M2M, M2H, H2H và các ví dụ\\
9--10 & Tổng quan giao thức IoT; định danh nút mạng\\
11--15 & URI/EPC/uCode; IPv4 và IPv6; hai hình so sánh địa chỉ và lợi ích IPv6\\
16--18 & Các lợi ích của IPv6; biểu tượng các chuẩn kết nối IoT\\
19--21 & PLC và X10; bảng thông số X10\\
22--23 & NFC; bảng tần số, tầm xa, chế độ và mã hóa\\
24--26 & UWB; Wi-Fi, Wi-Fi Direct và bảng tính năng Wi-Fi HaLow\\
27--28 & BLE và bảng so sánh Bluetooth Classic/BLE\\
29--31 & Zigbee, tính chất giao thức, hình sao--cây--mesh\\
32--35 & Z-Wave, Weightless, Sigfox, LoRaWAN\\
36--39 & Li-Fi, mạng tế bào, 5G, bảng 4G so với 5G\\
40--42 & LTE-M, NB-IoT, bảng so sánh đa công nghệ vô tuyến\\
43--46 & IPsec/TLS/DTLS/IEEE 1888.3, RPL và hai sơ đồ phân loại/định tuyến\\
47--55 & Bốn loại mạng IoT, ví dụ cellular/LAN-PAN/LPWAN/mesh, biểu đồ tầm xa--thông lượng\\
56--60 & Giao thức dữ liệu, đăng ký nút, publish/subscribe, request/response, liệt kê giao thức\\
61--65 & HTTP, phương thức HTTP, hạn chế IoT, sơ đồ HTTP và TCP so với UDP\\
66--71 & CoAP, proxy, REST, CON/NON/ACK/RST, luồng gói và ví dụ GET nhiệt độ\\
72--76 & MQTT, kiến trúc broker, bốn giai đoạn phiên, QoS 0--2, ví dụ áp dụng\\
77--78 & WebSocket, tương tác hai chiều và so sánh hình với HTTP\\
79--80 & AMQP (exchanges/queues) và XMPP/XML\\
81 & Kết thúc và lời cảm ơn\\\bottomrule
\end{longtable}

\begin{nguonslide}[title={Nguồn và phương pháp biên soạn}]
Tài liệu chính: TS. Nguyễn Đắc Cử, \textit{Cơ sở và thiết kế hệ thống IoT -- Chapter 3: Networking technologies in IoT}, Khoa Điện -- Điện tử, Trường Đại học Phenikaa, bản PDF 81 trang. Cách trình bày ở đây là \textbf{bản tổng hợp tiếng Việt theo chủ đề} chứ không phải bản dịch đối chiếu từng trang. Các sơ đồ đã được vẽ lại, dữ liệu theo bảng slide được ghi chú riêng; những điểm không nhất quán trong slide (như đơn vị gói X10, thiết lập NFC và phạm vi cellular trong bảng so sánh) được chỉ rõ chứ không tự ý thay thế. Khi tệp PDF gốc có mặt tại \texttt{source/chapter3-slides.pdf}, có thể đối chiếu trực tiếp với bảng trang ở trên.
\end{nguonslide}
% Phụ lục đầy đủ 50 slide gốc, tùy chọn khi tác giả có quyền sử dụng nguồn.
\IfFileExists{chapter1_source.pdf}{%
  \clearpage\chapter*{Phụ lục: 50 slide bài giảng gốc}\addcontentsline{toc}{chapter}{Phụ lục: 50 slide bài giảng gốc}
  \noindent Phụ lục này giữ nguyên hình minh họa trong toàn bộ 50 trang nguồn để đối chiếu với bản tổng hợp tiếng Việt.
  \clearpage\includepdf[pages=-]{chapter1_source.pdf}%
}{}
\end{document}